%% Generated by Sphinx.
\def\sphinxdocclass{report}
\documentclass[letterpaper,10pt,french]{sphinxmanual}
\ifdefined\pdfpxdimen
   \let\sphinxpxdimen\pdfpxdimen\else\newdimen\sphinxpxdimen
\fi \sphinxpxdimen=.75bp\relax
\ifdefined\pdfimageresolution
    \pdfimageresolution= \numexpr \dimexpr1in\relax/\sphinxpxdimen\relax
\fi
\newdimen\sphinxremdimen\sphinxremdimen = 10pt
%% let collapsible pdf bookmarks panel have high depth per default
\PassOptionsToPackage{bookmarksdepth=5}{hyperref}

\PassOptionsToPackage{booktabs}{sphinx}
\PassOptionsToPackage{colorrows}{sphinx}

\PassOptionsToPackage{warn}{textcomp}
\usepackage[utf8]{inputenc}
\ifdefined\DeclareUnicodeCharacter
% support both utf8 and utf8x syntaxes
  \ifdefined\DeclareUnicodeCharacterAsOptional
    \def\sphinxDUC#1{\DeclareUnicodeCharacter{"#1}}
  \else
    \let\sphinxDUC\DeclareUnicodeCharacter
  \fi
  \sphinxDUC{00A0}{\nobreakspace}
  \sphinxDUC{2500}{\sphinxunichar{2500}}
  \sphinxDUC{2502}{\sphinxunichar{2502}}
  \sphinxDUC{2514}{\sphinxunichar{2514}}
  \sphinxDUC{251C}{\sphinxunichar{251C}}
  \sphinxDUC{2572}{\textbackslash}
\fi
\usepackage{cmap}
\usepackage[T1]{fontenc}
\usepackage{amsmath,amssymb,amstext}
\usepackage{babel}



\usepackage{tgtermes}
\usepackage{tgheros}
\renewcommand{\ttdefault}{txtt}



\usepackage[Sonny]{fncychap}
\ChNameVar{\Large\normalfont\sffamily}
\ChTitleVar{\Large\normalfont\sffamily}
\usepackage{sphinx}

\fvset{fontsize=auto}
\usepackage{geometry}


% Include hyperref last.
\usepackage{hyperref}
% Fix anchor placement for figures with captions.
\usepackage{hypcap}% it must be loaded after hyperref.
% Set up styles of URL: it should be placed after hyperref.
\urlstyle{same}

\addto\captionsfrench{\renewcommand{\contentsname}{Référence API}}

\usepackage{sphinxmessages}
\setcounter{tocdepth}{1}



\title{Read\_For\_Me}
\date{01 avril 2026}
\release{1.0}
\author{Read\_For\_Me Team}
\newcommand{\sphinxlogo}{\vbox{}}
\renewcommand{\releasename}{Version}
\makeindex
\begin{document}

\ifdefined\shorthandoff
  \ifnum\catcode`\=\string=\active\shorthandoff{=}\fi
  \ifnum\catcode`\"=\active\shorthandoff{"}\fi
\fi

\pagestyle{empty}
\sphinxmaketitle
\pagestyle{plain}
\sphinxtableofcontents
\pagestyle{normal}
\phantomsection\label{\detokenize{index::doc}}


\sphinxAtStartPar
Cette documentation est générée automatiquement à partir des docstrings
du code Python du projet \sphinxcode{\sphinxupquote{Read\_For\_Me}}.

\sphinxAtStartPar
Elle décrit les modules principaux de l’application :
\begin{itemize}
\item {} 
\sphinxAtStartPar
\sphinxcode{\sphinxupquote{main}}

\item {} 
\sphinxAtStartPar
\sphinxcode{\sphinxupquote{config}}

\item {} 
\sphinxAtStartPar
\sphinxcode{\sphinxupquote{core}}

\item {} 
\sphinxAtStartPar
\sphinxcode{\sphinxupquote{hardware}}

\item {} 
\sphinxAtStartPar
\sphinxcode{\sphinxupquote{modes}}

\end{itemize}

\sphinxstepscope


\chapter{Read\_For\_Me}
\label{\detokenize{api/modules:read-for-me}}\label{\detokenize{api/modules::doc}}
\sphinxstepscope


\section{config package}
\label{\detokenize{api/config:config-package}}\label{\detokenize{api/config::doc}}

\subsection{Submodules}
\label{\detokenize{api/config:submodules}}

\subsection{config.config\_loader module}
\label{\detokenize{api/config:module-config.config_loader}}\label{\detokenize{api/config:config-config-loader-module}}\index{module@\spxentry{module}!config.config\_loader@\spxentry{config.config\_loader}}\index{config.config\_loader@\spxentry{config.config\_loader}!module@\spxentry{module}}
\sphinxAtStartPar
Charge et valide la configuration centrale depuis \sphinxcode{\sphinxupquote{config/config.toml}}.

\sphinxAtStartPar
Le loader renvoie un dictionnaire unique pour toute l’application avec
valeurs par défaut fusionnées puis validées.
\index{load\_config() (dans le module config.config\_loader)@\spxentry{load\_config()}\spxextra{dans le module config.config\_loader}}

\begin{fulllineitems}
\phantomsection\label{\detokenize{api/config:config.config_loader.load_config}}
\pysigstartsignatures
\pysiglinewithargsret
{\sphinxcode{\sphinxupquote{config.config\_loader.}}\sphinxbfcode{\sphinxupquote{load\_config}}}
{}
{{ $\rightarrow$ Dict\DUrole{p}{{[}}str\DUrole{p}{,}\DUrole{w}{ }Any\DUrole{p}{{]}}}}
\pysigstopsignatures
\sphinxAtStartPar
Charge \sphinxtitleref{config.toml}, fusionne les valeurs par défaut puis valide le résultat.

\end{fulllineitems}

\index{get\_app\_config() (dans le module config.config\_loader)@\spxentry{get\_app\_config()}\spxextra{dans le module config.config\_loader}}

\begin{fulllineitems}
\phantomsection\label{\detokenize{api/config:config.config_loader.get_app_config}}
\pysigstartsignatures
\pysiglinewithargsret
{\sphinxcode{\sphinxupquote{config.config\_loader.}}\sphinxbfcode{\sphinxupquote{get\_app\_config}}}
{}
{{ $\rightarrow$ Dict\DUrole{p}{{[}}str\DUrole{p}{,}\DUrole{w}{ }Any\DUrole{p}{{]}}}}
\pysigstopsignatures
\sphinxAtStartPar
Retourne la configuration complète validée de l’application.

\end{fulllineitems}

\index{get\_speaker\_config() (dans le module config.config\_loader)@\spxentry{get\_speaker\_config()}\spxextra{dans le module config.config\_loader}}

\begin{fulllineitems}
\phantomsection\label{\detokenize{api/config:config.config_loader.get_speaker_config}}
\pysigstartsignatures
\pysiglinewithargsret
{\sphinxcode{\sphinxupquote{config.config\_loader.}}\sphinxbfcode{\sphinxupquote{get\_speaker\_config}}}
{}
{{ $\rightarrow$ Dict\DUrole{p}{{[}}str\DUrole{p}{,}\DUrole{w}{ }Any\DUrole{p}{{]}}}}
\pysigstopsignatures
\sphinxAtStartPar
Retourne uniquement la section \sphinxtitleref{speaker} de la configuration validée.

\end{fulllineitems}



\subsection{Module contents}
\label{\detokenize{api/config:module-config}}\label{\detokenize{api/config:module-contents}}\index{module@\spxentry{module}!config@\spxentry{config}}\index{config@\spxentry{config}!module@\spxentry{module}}
\sphinxstepscope


\section{core package}
\label{\detokenize{api/core:core-package}}\label{\detokenize{api/core::doc}}

\subsection{Submodules}
\label{\detokenize{api/core:submodules}}

\subsection{core.config\_loader module}
\label{\detokenize{api/core:module-core.config_loader}}\label{\detokenize{api/core:core-config-loader-module}}\index{module@\spxentry{module}!core.config\_loader@\spxentry{core.config\_loader}}\index{core.config\_loader@\spxentry{core.config\_loader}!module@\spxentry{module}}
\sphinxAtStartPar
Pont de compatibilité vers le loader central de configuration.

\sphinxAtStartPar
Le point d’entrée officiel est \sphinxcode{\sphinxupquote{config.config\_loader}}.
\index{get\_config() (dans le module core.config\_loader)@\spxentry{get\_config()}\spxextra{dans le module core.config\_loader}}

\begin{fulllineitems}
\phantomsection\label{\detokenize{api/core:core.config_loader.get_config}}
\pysigstartsignatures
\pysiglinewithargsret
{\sphinxcode{\sphinxupquote{core.config\_loader.}}\sphinxbfcode{\sphinxupquote{get\_config}}}
{}
{{ $\rightarrow$ Dict\DUrole{p}{{[}}str\DUrole{p}{,}\DUrole{w}{ }Any\DUrole{p}{{]}}}}
\pysigstopsignatures
\sphinxAtStartPar
Alias lisible pour récupérer la configuration complète.

\end{fulllineitems}



\subsection{core.mode\_base module}
\label{\detokenize{api/core:module-core.mode_base}}\label{\detokenize{api/core:core-mode-base-module}}\index{module@\spxentry{module}!core.mode\_base@\spxentry{core.mode\_base}}\index{core.mode\_base@\spxentry{core.mode\_base}!module@\spxentry{module}}

\subsubsection{core.mode\_base}
\label{\detokenize{api/core:core-mode-base}}
\sphinxAtStartPar
Définit la classe abstraite \sphinxtitleref{Mode}, base commune à \sphinxstylestrong{tous les modes
fonctionnels} de l’assistant (date/heure, multimètre, mode test, etc.).


\paragraph{Rôle de cette classe}
\label{\detokenize{api/core:role-de-cette-classe}}
\sphinxAtStartPar
Un \sphinxstyleemphasis{mode} représente un « contexte » d’utilisation de l’assistant :
par exemple « Mode Date et heure », « Mode Multimètre », etc.
À un instant donné, \sphinxstylestrong{un seul mode est actif}. Le changement de mode
est piloté par le sélecteur rotatif et orchestré par \sphinxtitleref{ModeManager}.

\sphinxAtStartPar
Cette classe définit le \sphinxstylestrong{contrat minimal} qu’un mode doit respecter :
\begin{itemize}
\item {} 
\sphinxAtStartPar
un \sphinxstylestrong{nom} lisible (\sphinxtitleref{name}) pour l’annonce vocale du mode,

\item {} \begin{description}
\sphinxlineitem{des hooks de cycle de vie :}\begin{itemize}
\item {} 
\sphinxAtStartPar
\sphinxtitleref{on\_enter()} : appelé quand on arrive sur le mode,

\item {} 
\sphinxAtStartPar
\sphinxtitleref{on\_exit()} : appelé quand on quitte le mode,

\end{itemize}

\end{description}

\item {} \begin{description}
\sphinxlineitem{des callbacks liés au \sphinxstylestrong{bouton poussoir} :}\begin{itemize}
\item {} 
\sphinxAtStartPar
\sphinxtitleref{on\_short\_press()} : appui court,

\item {} 
\sphinxAtStartPar
\sphinxtitleref{on\_long\_press()}  : appui long,

\item {} 
\sphinxAtStartPar
\sphinxtitleref{on\_double\_press()} : double appui rapide (optionnel, valeur par défaut : ne fait rien).

\end{itemize}

\end{description}

\end{itemize}

\sphinxAtStartPar
Cette classe ne connaît rien :
\sphinxhyphen{} ni de la synthèse vocale (\sphinxtitleref{Speaker}),
\sphinxhyphen{} ni du matériel (rotary, BLE, etc.),
\sphinxhyphen{} ni du \sphinxtitleref{ModeManager} qui l’utilise.

\sphinxAtStartPar
Elle sert uniquement de \sphinxstylestrong{base commune}, pour garder une interface
cohérente entre tous les modes.
\index{Mode (classe dans core.mode\_base)@\spxentry{Mode}\spxextra{classe dans core.mode\_base}}

\begin{fulllineitems}
\phantomsection\label{\detokenize{api/core:core.mode_base.Mode}}
\pysigstartsignatures
\pysiglinewithargsret
{\sphinxbfcode{\sphinxupquote{\DUrole{k}{class}\DUrole{w}{ }}}\sphinxcode{\sphinxupquote{core.mode\_base.}}\sphinxbfcode{\sphinxupquote{Mode}}}
{\sphinxparam{\DUrole{n}{name}\DUrole{p}{:}\DUrole{w}{ }\DUrole{n}{str}}}
{}
\pysigstopsignatures
\sphinxAtStartPar
Bases : \sphinxcode{\sphinxupquote{ABC}}

\sphinxAtStartPar
Classe de base abstraite pour tous les modes de l’assistant.

\sphinxAtStartPar
Chaque mode concret (ex: \sphinxtitleref{ModeDateHeure}, \sphinxtitleref{ModeMultimetre}, etc.)
doit \sphinxstylestrong{hériter} de cette classe et implémenter au minimum :
\begin{itemize}
\item {} 
\sphinxAtStartPar
\sphinxtitleref{on\_short\_press()} : réaction à un appui court sur le bouton,

\item {} 
\sphinxAtStartPar
\sphinxtitleref{on\_long\_press()}  : réaction à un appui long sur le bouton.

\end{itemize}

\sphinxAtStartPar
Les hooks suivants sont \sphinxstylestrong{optionnels} (no\sphinxhyphen{}op par défaut) :
\begin{itemize}
\item {} 
\sphinxAtStartPar
\sphinxtitleref{on\_enter()} / \sphinxtitleref{on\_exit()} : lifecycle du mode (drivers, timers…).

\item {} 
\sphinxAtStartPar
\sphinxtitleref{on\_double\_press()} : double appui sur le bouton rotatif.

\item {} 
\sphinxAtStartPar
\sphinxtitleref{on\_key\_pressed(key)} : touche clavier reçue du ModeManager.

\item {} 
\sphinxAtStartPar
\sphinxtitleref{on\_keypad\_key(key)} : alias de compatibilité de \sphinxtitleref{on\_key\_pressed}.
Seules les touches NON\sphinxhyphen{}globales sont transmises ici — les touches
réservées au volume et à la vitesse sont interceptées par le
ModeManager avant d’arriver dans les modes.

\end{itemize}
\index{\_\_init\_\_() (méthode core.mode\_base.Mode)@\spxentry{\_\_init\_\_()}\spxextra{méthode core.mode\_base.Mode}}

\begin{fulllineitems}
\phantomsection\label{\detokenize{api/core:core.mode_base.Mode.__init__}}
\pysigstartsignatures
\pysiglinewithargsret
{\sphinxbfcode{\sphinxupquote{\_\_init\_\_}}}
{\sphinxparam{\DUrole{n}{name}\DUrole{p}{:}\DUrole{w}{ }\DUrole{n}{str}}}
{}
\pysigstopsignatures
\sphinxAtStartPar
Initialise un mode abstrait.
\begin{quote}\begin{description}
\sphinxlineitem{Paramètres}
\sphinxAtStartPar
\sphinxstyleliteralstrong{\sphinxupquote{name}} \textendash{} Nom lisible du mode, utilisé pour les annonces vocales
(exemples : « Machine à lire », « Multimètre », « Date et heure »).

\end{description}\end{quote}

\end{fulllineitems}

\index{on\_enter() (méthode core.mode\_base.Mode)@\spxentry{on\_enter()}\spxextra{méthode core.mode\_base.Mode}}

\begin{fulllineitems}
\phantomsection\label{\detokenize{api/core:core.mode_base.Mode.on_enter}}
\pysigstartsignatures
\pysiglinewithargsret
{\sphinxbfcode{\sphinxupquote{on\_enter}}}
{}
{{ $\rightarrow$ None}}
\pysigstopsignatures
\sphinxAtStartPar
Hook appelé quand on \sphinxstylestrong{arrive} sur ce mode.

\sphinxAtStartPar
Ce hook est invoqué par le \sphinxtitleref{ModeManager} lorsque le sélecteur
rotatif pointe sur ce mode et que ce mode devient le mode actif.

\sphinxAtStartPar
Par défaut, cette méthode ne fait rien. Les sous\sphinxhyphen{}classes peuvent
la surcharger pour, par exemple :
\begin{itemize}
\item {} 
\sphinxAtStartPar
démarrer une tâche ou un driver matériel,

\item {} 
\sphinxAtStartPar
initialiser un état interne,

\item {} 
\sphinxAtStartPar
annoncer un message spécifique au mode, etc.

\end{itemize}

\end{fulllineitems}

\index{on\_exit() (méthode core.mode\_base.Mode)@\spxentry{on\_exit()}\spxextra{méthode core.mode\_base.Mode}}

\begin{fulllineitems}
\phantomsection\label{\detokenize{api/core:core.mode_base.Mode.on_exit}}
\pysigstartsignatures
\pysiglinewithargsret
{\sphinxbfcode{\sphinxupquote{on\_exit}}}
{}
{{ $\rightarrow$ None}}
\pysigstopsignatures
\sphinxAtStartPar
Hook appelé quand on \sphinxstylestrong{quitte} ce mode.

\sphinxAtStartPar
Ce hook est invoqué par le \sphinxtitleref{ModeManager} juste avant de basculer
vers un autre mode (lors d’une rotation stable du sélecteur).

\sphinxAtStartPar
Par défaut, cette méthode ne fait rien. Les sous\sphinxhyphen{}classes peuvent
la surcharger pour, par exemple :
\begin{itemize}
\item {} 
\sphinxAtStartPar
arrêter proprement un driver,

\item {} 
\sphinxAtStartPar
libérer des ressources,

\item {} 
\sphinxAtStartPar
annuler des timers, etc.

\end{itemize}

\end{fulllineitems}

\index{on\_short\_press() (méthode core.mode\_base.Mode)@\spxentry{on\_short\_press()}\spxextra{méthode core.mode\_base.Mode}}

\begin{fulllineitems}
\phantomsection\label{\detokenize{api/core:core.mode_base.Mode.on_short_press}}
\pysigstartsignatures
\pysiglinewithargsret
{\sphinxbfcode{\sphinxupquote{\DUrole{k}{abstractmethod}\DUrole{w}{ }}}\sphinxbfcode{\sphinxupquote{on\_short\_press}}}
{}
{{ $\rightarrow$ None}}
\pysigstopsignatures
\sphinxAtStartPar
Appui \sphinxstylestrong{court} sur le bouton poussoir lorsque ce mode est actif.

\sphinxAtStartPar
Cette méthode \sphinxstylestrong{doit} être implémentée par chaque mode concret.

\sphinxAtStartPar
Exemple d’usage typique :
\begin{itemize}
\item {} 
\sphinxAtStartPar
dans un mode « Multimètre » : annoncer la dernière mesure,

\item {} 
\sphinxAtStartPar
dans un mode « Date et heure » : lire l’heure actuelle,

\item {} 
\sphinxAtStartPar
dans un mode « Lecture » : lire la ligne ou l’élément suivant.

\end{itemize}

\end{fulllineitems}

\index{on\_long\_press() (méthode core.mode\_base.Mode)@\spxentry{on\_long\_press()}\spxextra{méthode core.mode\_base.Mode}}

\begin{fulllineitems}
\phantomsection\label{\detokenize{api/core:core.mode_base.Mode.on_long_press}}
\pysigstartsignatures
\pysiglinewithargsret
{\sphinxbfcode{\sphinxupquote{\DUrole{k}{abstractmethod}\DUrole{w}{ }}}\sphinxbfcode{\sphinxupquote{on\_long\_press}}}
{}
{{ $\rightarrow$ None}}
\pysigstopsignatures
\sphinxAtStartPar
Appui \sphinxstylestrong{long} sur le bouton poussoir lorsque ce mode est actif.

\sphinxAtStartPar
Cette méthode \sphinxstylestrong{doit} être implémentée par chaque mode concret.

\sphinxAtStartPar
Exemple d’usage typique :
\begin{itemize}
\item {} 
\sphinxAtStartPar
afficher ou annoncer un état détaillé,

\item {} 
\sphinxAtStartPar
activer une fonction secondaire,

\item {} 
\sphinxAtStartPar
basculer dans un sous\sphinxhyphen{}mode, etc.

\end{itemize}

\end{fulllineitems}

\index{on\_double\_press() (méthode core.mode\_base.Mode)@\spxentry{on\_double\_press()}\spxextra{méthode core.mode\_base.Mode}}

\begin{fulllineitems}
\phantomsection\label{\detokenize{api/core:core.mode_base.Mode.on_double_press}}
\pysigstartsignatures
\pysiglinewithargsret
{\sphinxbfcode{\sphinxupquote{on\_double\_press}}}
{}
{{ $\rightarrow$ None}}
\pysigstopsignatures
\sphinxAtStartPar
\sphinxstylestrong{Double appui rapide} sur le bouton poussoir lorsque ce mode est actif.

\sphinxAtStartPar
Ce hook est \sphinxstylestrong{optionnel} : par défaut, il ne fait rien.
Il peut être surchargé par un mode qui souhaite proposer un
comportement spécifique sur double clic.

\sphinxAtStartPar
Exemple d’usage typique :
\begin{itemize}
\item {} 
\sphinxAtStartPar
dans le mode « Multimètre » : activer / désactiver la
lecture automatique toutes les X secondes,

\item {} 
\sphinxAtStartPar
dans un autre mode : activer un mode « favori » ou une fonction
avancée sans surcharger le simple appui court.

\end{itemize}

\end{fulllineitems}

\index{on\_key\_pressed() (méthode core.mode\_base.Mode)@\spxentry{on\_key\_pressed()}\spxextra{méthode core.mode\_base.Mode}}

\begin{fulllineitems}
\phantomsection\label{\detokenize{api/core:core.mode_base.Mode.on_key_pressed}}
\pysigstartsignatures
\pysiglinewithargsret
{\sphinxbfcode{\sphinxupquote{on\_key\_pressed}}}
{\sphinxparam{\DUrole{n}{key}\DUrole{p}{:}\DUrole{w}{ }\DUrole{n}{str}}}
{{ $\rightarrow$ None}}
\pysigstopsignatures
\sphinxAtStartPar
Point d’entrée recommandé pour les touches keypad non\sphinxhyphen{}globales.

\sphinxAtStartPar
Par défaut, délègue vers \sphinxtitleref{on\_keypad\_key} pour préserver la
compatibilité des modes existants.

\end{fulllineitems}

\index{on\_keypad\_key() (méthode core.mode\_base.Mode)@\spxentry{on\_keypad\_key()}\spxextra{méthode core.mode\_base.Mode}}

\begin{fulllineitems}
\phantomsection\label{\detokenize{api/core:core.mode_base.Mode.on_keypad_key}}
\pysigstartsignatures
\pysiglinewithargsret
{\sphinxbfcode{\sphinxupquote{on\_keypad\_key}}}
{\sphinxparam{\DUrole{n}{key}\DUrole{p}{:}\DUrole{w}{ }\DUrole{n}{str}}}
{{ $\rightarrow$ None}}
\pysigstopsignatures
\sphinxAtStartPar
Alias de compatibilité pour la gestion des touches keypad.

\sphinxAtStartPar
Les nouveaux modes devraient préférer surcharger \sphinxtitleref{on\_key\_pressed}.
Les modes existants qui surchargent \sphinxtitleref{on\_keypad\_key} restent
entièrement compatibles.


\subsubsection{Règle importante}
\label{\detokenize{api/core:regle-importante}}
\sphinxAtStartPar
Seules les touches NON réservées par le ModeManager sont
transmises ici. Les touches globales (volume+, volume\sphinxhyphen{},
vitesse+, vitesse\sphinxhyphen{}) sont interceptées en amont et ne
parviennent jamais dans cette méthode.


\subsubsection{Paramètres}
\label{\detokenize{api/core:parametres}}\begin{description}
\sphinxlineitem{key :}
\sphinxAtStartPar
Chaîne représentant la touche pressée (« 1 »\textendash{}« 9 », « 0 »,
« A »\textendash{}« D », « * », « \# »).

\end{description}


\subsubsection{Exemple d’usage}
\label{\detokenize{api/core:exemple-dusage}}
\sphinxAtStartPar
Dans ModeReadingMachine :
\begin{quote}
\begin{description}
\sphinxlineitem{def on\_keypad\_key(self, key: str) \sphinxhyphen{}\textgreater{} None:}\begin{description}
\sphinxlineitem{if key == « 1 »:}
\sphinxAtStartPar
self.\_start\_capture\_pipeline()

\sphinxlineitem{elif key == « 5 »:}
\sphinxAtStartPar
self.\_toggle\_pause()

\end{description}

\sphinxAtStartPar
…

\end{description}
\end{quote}

\end{fulllineitems}


\end{fulllineitems}



\subsection{core.mode\_manager module}
\label{\detokenize{api/core:module-core.mode_manager}}\label{\detokenize{api/core:core-mode-manager-module}}\index{module@\spxentry{module}!core.mode\_manager@\spxentry{core.mode\_manager}}\index{core.mode\_manager@\spxentry{core.mode\_manager}!module@\spxentry{module}}

\subsubsection{core.mode\_manager}
\label{\detokenize{api/core:core-mode-manager}}
\sphinxAtStartPar
Orchestrateur central des modes de l’assistant.


\paragraph{Responsabilités}
\label{\detokenize{api/core:responsabilites}}\begin{enumerate}
\sphinxsetlistlabels{\arabic}{enumi}{enumii}{}{.}%
\item {} 
\sphinxAtStartPar
Gérer la liste des modes et le mode courant.

\item {} 
\sphinxAtStartPar
Changer de mode quand le sélecteur rotatif tourne (set\_index).

\item {} 
\sphinxAtStartPar
Router les événements bouton rotatif vers le mode actif
(handle\_short\_press, handle\_long\_press, handle\_double\_press).

\item {} 
\sphinxAtStartPar
Optionnellement gérer un clavier matriciel 4x4 (Keypad4x4) :
\sphinxhyphen{} Les touches « globales » (volume +/\sphinxhyphen{}, vitesse +/\sphinxhyphen{}) sont
\begin{quote}

\sphinxAtStartPar
interceptées ici et traitées directement via le Speaker.
Elles sont actives dans TOUS les modes, tout le temps.
\end{quote}
\begin{itemize}
\item {} 
\sphinxAtStartPar
Les autres touches sont transmises au mode actif via
Mode.on\_keypad\_key(key). Seuls les modes qui le souhaitent
surchargent cette méthode (sinon : no\sphinxhyphen{}op).

\end{itemize}

\end{enumerate}


\paragraph{Architecture du clavier}
\label{\detokenize{api/core:architecture-du-clavier}}
\sphinxAtStartPar
Le Keypad4x4 est démarré UNE SEULE FOIS dans \_\_init\_\_ et tourne
en continu, contrairement à l’ancienne approche où chaque mode
le démarrait/arrêtait dans on\_enter/on\_exit.

\sphinxAtStartPar
Cela permet aux touches globales d’être disponibles dans n’importe
quel mode, sans aucune gestion spécifique dans les modes eux\sphinxhyphen{}mêmes.


\paragraph{Touches globales (non\sphinxhyphen{}surchargeables par les modes)}
\label{\detokenize{api/core:touches-globales-non-surchargeables-par-les-modes}}\begin{quote}

\sphinxAtStartPar
KEY\_GLOBAL\_VOL\_UP   = « 8 »   \(\rightarrow\) volume + VOLUME\_STEP
KEY\_GLOBAL\_VOL\_DOWN = « 2 »   \(\rightarrow\) volume \sphinxhyphen{} VOLUME\_STEP
KEY\_GLOBAL\_SPEED\_UP = « 9 »   \(\rightarrow\) vitesse × SPEED\_FACTOR\_UP
KEY\_GLOBAL\_SPEED\_DOWN = « 7 » \(\rightarrow\) vitesse × SPEED\_FACTOR\_DOWN
\end{quote}

\sphinxAtStartPar
Un bip de feedback est joué à chaque changement (aigu=+, grave=\sphinxhyphen{}).
\index{ModeManager (classe dans core.mode\_manager)@\spxentry{ModeManager}\spxextra{classe dans core.mode\_manager}}

\begin{fulllineitems}
\phantomsection\label{\detokenize{api/core:core.mode_manager.ModeManager}}
\pysigstartsignatures
\pysiglinewithargsret
{\sphinxbfcode{\sphinxupquote{\DUrole{k}{class}\DUrole{w}{ }}}\sphinxcode{\sphinxupquote{core.mode\_manager.}}\sphinxbfcode{\sphinxupquote{ModeManager}}}
{\sphinxparam{\DUrole{n}{modes}\DUrole{p}{:}\DUrole{w}{ }\DUrole{n}{List\DUrole{p}{{[}}{\hyperref[\detokenize{api/core:core.mode_base.Mode}]{\sphinxcrossref{Mode}}}\DUrole{p}{{]}}}}\sphinxparamcomma \sphinxparam{\DUrole{n}{speaker}\DUrole{p}{:}\DUrole{w}{ }\DUrole{n}{{\hyperref[\detokenize{api/core:core.speaker.Speaker}]{\sphinxcrossref{Speaker}}}}}\sphinxparamcomma \sphinxparam{\DUrole{n}{keypad}\DUrole{o}{=}\DUrole{default_value}{None}}\sphinxparamcomma \sphinxparam{\DUrole{n}{key\_global\_vol\_up}\DUrole{p}{:}\DUrole{w}{ }\DUrole{n}{str}\DUrole{w}{ }\DUrole{o}{=}\DUrole{w}{ }\DUrole{default_value}{\textquotesingle{}8\textquotesingle{}}}\sphinxparamcomma \sphinxparam{\DUrole{n}{key\_global\_vol\_down}\DUrole{p}{:}\DUrole{w}{ }\DUrole{n}{str}\DUrole{w}{ }\DUrole{o}{=}\DUrole{w}{ }\DUrole{default_value}{\textquotesingle{}2\textquotesingle{}}}\sphinxparamcomma \sphinxparam{\DUrole{n}{key\_global\_speed\_up}\DUrole{p}{:}\DUrole{w}{ }\DUrole{n}{str}\DUrole{w}{ }\DUrole{o}{=}\DUrole{w}{ }\DUrole{default_value}{\textquotesingle{}9\textquotesingle{}}}\sphinxparamcomma \sphinxparam{\DUrole{n}{key\_global\_speed\_down}\DUrole{p}{:}\DUrole{w}{ }\DUrole{n}{str}\DUrole{w}{ }\DUrole{o}{=}\DUrole{w}{ }\DUrole{default_value}{\textquotesingle{}7\textquotesingle{}}}\sphinxparamcomma \sphinxparam{\DUrole{n}{volume\_step}\DUrole{p}{:}\DUrole{w}{ }\DUrole{n}{int}\DUrole{w}{ }\DUrole{o}{=}\DUrole{w}{ }\DUrole{default_value}{4}}\sphinxparamcomma \sphinxparam{\DUrole{n}{speed\_factor\_up}\DUrole{p}{:}\DUrole{w}{ }\DUrole{n}{float}\DUrole{w}{ }\DUrole{o}{=}\DUrole{w}{ }\DUrole{default_value}{1.25}}\sphinxparamcomma \sphinxparam{\DUrole{n}{speed\_factor\_down}\DUrole{p}{:}\DUrole{w}{ }\DUrole{n}{float}\DUrole{w}{ }\DUrole{o}{=}\DUrole{w}{ }\DUrole{default_value}{0.8}}}
{}
\pysigstopsignatures
\sphinxAtStartPar
Bases : \sphinxcode{\sphinxupquote{object}}

\sphinxAtStartPar
Chef d’orchestre des modes de l’assistant.


\subsubsection{Paramètres}
\label{\detokenize{api/core:id1}}\begin{description}
\sphinxlineitem{modes :}
\sphinxAtStartPar
Liste ordonnée de modes (index 0 = position 0 du rotatif).

\sphinxlineitem{speaker :}
\sphinxAtStartPar
Speaker central — utilisé pour les annonces et les réglages
globaux (volume, vitesse, bip).

\sphinxlineitem{keypad :}
\sphinxAtStartPar
Instance de Keypad4x4 (optionnel). Si fourni, le ModeManager
démarre le clavier et gère les touches globales + la délégation
vers le mode courant. Si None, le clavier est ignoré.

\sphinxlineitem{volume\_step :}
\sphinxAtStartPar
Nombre de points ALSA ajoutés/retirés par pression (0\textendash{}100).

\sphinxlineitem{speed\_factor\_up :}
\sphinxAtStartPar
Multiplicateur appliqué à la vitesse pour une augmentation.

\sphinxlineitem{speed\_factor\_down :}
\sphinxAtStartPar
Multiplicateur appliqué à la vitesse pour une diminution.

\end{description}
\index{KEY\_GLOBAL\_VOL\_UP (attribut core.mode\_manager.ModeManager)@\spxentry{KEY\_GLOBAL\_VOL\_UP}\spxextra{attribut core.mode\_manager.ModeManager}}

\begin{fulllineitems}
\phantomsection\label{\detokenize{api/core:core.mode_manager.ModeManager.KEY_GLOBAL_VOL_UP}}
\pysigstartsignatures
\pysigline
{\sphinxbfcode{\sphinxupquote{KEY\_GLOBAL\_VOL\_UP}}\sphinxbfcode{\sphinxupquote{\DUrole{w}{ }\DUrole{p}{=}\DUrole{w}{ }\textquotesingle{}8\textquotesingle{}}}}
\pysigstopsignatures
\end{fulllineitems}

\index{KEY\_GLOBAL\_VOL\_DOWN (attribut core.mode\_manager.ModeManager)@\spxentry{KEY\_GLOBAL\_VOL\_DOWN}\spxextra{attribut core.mode\_manager.ModeManager}}

\begin{fulllineitems}
\phantomsection\label{\detokenize{api/core:core.mode_manager.ModeManager.KEY_GLOBAL_VOL_DOWN}}
\pysigstartsignatures
\pysigline
{\sphinxbfcode{\sphinxupquote{KEY\_GLOBAL\_VOL\_DOWN}}\sphinxbfcode{\sphinxupquote{\DUrole{w}{ }\DUrole{p}{=}\DUrole{w}{ }\textquotesingle{}2\textquotesingle{}}}}
\pysigstopsignatures
\end{fulllineitems}

\index{KEY\_GLOBAL\_SPEED\_UP (attribut core.mode\_manager.ModeManager)@\spxentry{KEY\_GLOBAL\_SPEED\_UP}\spxextra{attribut core.mode\_manager.ModeManager}}

\begin{fulllineitems}
\phantomsection\label{\detokenize{api/core:core.mode_manager.ModeManager.KEY_GLOBAL_SPEED_UP}}
\pysigstartsignatures
\pysigline
{\sphinxbfcode{\sphinxupquote{KEY\_GLOBAL\_SPEED\_UP}}\sphinxbfcode{\sphinxupquote{\DUrole{w}{ }\DUrole{p}{=}\DUrole{w}{ }\textquotesingle{}9\textquotesingle{}}}}
\pysigstopsignatures
\end{fulllineitems}

\index{KEY\_GLOBAL\_SPEED\_DOWN (attribut core.mode\_manager.ModeManager)@\spxentry{KEY\_GLOBAL\_SPEED\_DOWN}\spxextra{attribut core.mode\_manager.ModeManager}}

\begin{fulllineitems}
\phantomsection\label{\detokenize{api/core:core.mode_manager.ModeManager.KEY_GLOBAL_SPEED_DOWN}}
\pysigstartsignatures
\pysigline
{\sphinxbfcode{\sphinxupquote{KEY\_GLOBAL\_SPEED\_DOWN}}\sphinxbfcode{\sphinxupquote{\DUrole{w}{ }\DUrole{p}{=}\DUrole{w}{ }\textquotesingle{}7\textquotesingle{}}}}
\pysigstopsignatures
\end{fulllineitems}

\index{GLOBAL\_KEYS (attribut core.mode\_manager.ModeManager)@\spxentry{GLOBAL\_KEYS}\spxextra{attribut core.mode\_manager.ModeManager}}

\begin{fulllineitems}
\phantomsection\label{\detokenize{api/core:core.mode_manager.ModeManager.GLOBAL_KEYS}}
\pysigstartsignatures
\pysigline
{\sphinxbfcode{\sphinxupquote{GLOBAL\_KEYS}}\sphinxbfcode{\sphinxupquote{\DUrole{w}{ }\DUrole{p}{=}\DUrole{w}{ }frozenset(\{\textquotesingle{}2\textquotesingle{}, \textquotesingle{}7\textquotesingle{}, \textquotesingle{}8\textquotesingle{}, \textquotesingle{}9\textquotesingle{}\})}}}
\pysigstopsignatures
\end{fulllineitems}

\index{VOLUME\_STEP (attribut core.mode\_manager.ModeManager)@\spxentry{VOLUME\_STEP}\spxextra{attribut core.mode\_manager.ModeManager}}

\begin{fulllineitems}
\phantomsection\label{\detokenize{api/core:core.mode_manager.ModeManager.VOLUME_STEP}}
\pysigstartsignatures
\pysigline
{\sphinxbfcode{\sphinxupquote{VOLUME\_STEP}}\sphinxbfcode{\sphinxupquote{\DUrole{p}{:}\DUrole{w}{ }int}}\sphinxbfcode{\sphinxupquote{\DUrole{w}{ }\DUrole{p}{=}\DUrole{w}{ }4}}}
\pysigstopsignatures
\end{fulllineitems}

\index{SPEED\_FACTOR\_UP (attribut core.mode\_manager.ModeManager)@\spxentry{SPEED\_FACTOR\_UP}\spxextra{attribut core.mode\_manager.ModeManager}}

\begin{fulllineitems}
\phantomsection\label{\detokenize{api/core:core.mode_manager.ModeManager.SPEED_FACTOR_UP}}
\pysigstartsignatures
\pysigline
{\sphinxbfcode{\sphinxupquote{SPEED\_FACTOR\_UP}}\sphinxbfcode{\sphinxupquote{\DUrole{p}{:}\DUrole{w}{ }float}}\sphinxbfcode{\sphinxupquote{\DUrole{w}{ }\DUrole{p}{=}\DUrole{w}{ }1.25}}}
\pysigstopsignatures
\end{fulllineitems}

\index{SPEED\_FACTOR\_DOWN (attribut core.mode\_manager.ModeManager)@\spxentry{SPEED\_FACTOR\_DOWN}\spxextra{attribut core.mode\_manager.ModeManager}}

\begin{fulllineitems}
\phantomsection\label{\detokenize{api/core:core.mode_manager.ModeManager.SPEED_FACTOR_DOWN}}
\pysigstartsignatures
\pysigline
{\sphinxbfcode{\sphinxupquote{SPEED\_FACTOR\_DOWN}}\sphinxbfcode{\sphinxupquote{\DUrole{p}{:}\DUrole{w}{ }float}}\sphinxbfcode{\sphinxupquote{\DUrole{w}{ }\DUrole{p}{=}\DUrole{w}{ }0.8}}}
\pysigstopsignatures
\end{fulllineitems}

\index{\_\_init\_\_() (méthode core.mode\_manager.ModeManager)@\spxentry{\_\_init\_\_()}\spxextra{méthode core.mode\_manager.ModeManager}}

\begin{fulllineitems}
\phantomsection\label{\detokenize{api/core:core.mode_manager.ModeManager.__init__}}
\pysigstartsignatures
\pysiglinewithargsret
{\sphinxbfcode{\sphinxupquote{\_\_init\_\_}}}
{\sphinxparam{\DUrole{n}{modes}\DUrole{p}{:}\DUrole{w}{ }\DUrole{n}{List\DUrole{p}{{[}}{\hyperref[\detokenize{api/core:core.mode_base.Mode}]{\sphinxcrossref{Mode}}}\DUrole{p}{{]}}}}\sphinxparamcomma \sphinxparam{\DUrole{n}{speaker}\DUrole{p}{:}\DUrole{w}{ }\DUrole{n}{{\hyperref[\detokenize{api/core:core.speaker.Speaker}]{\sphinxcrossref{Speaker}}}}}\sphinxparamcomma \sphinxparam{\DUrole{n}{keypad}\DUrole{o}{=}\DUrole{default_value}{None}}\sphinxparamcomma \sphinxparam{\DUrole{n}{key\_global\_vol\_up}\DUrole{p}{:}\DUrole{w}{ }\DUrole{n}{str}\DUrole{w}{ }\DUrole{o}{=}\DUrole{w}{ }\DUrole{default_value}{\textquotesingle{}8\textquotesingle{}}}\sphinxparamcomma \sphinxparam{\DUrole{n}{key\_global\_vol\_down}\DUrole{p}{:}\DUrole{w}{ }\DUrole{n}{str}\DUrole{w}{ }\DUrole{o}{=}\DUrole{w}{ }\DUrole{default_value}{\textquotesingle{}2\textquotesingle{}}}\sphinxparamcomma \sphinxparam{\DUrole{n}{key\_global\_speed\_up}\DUrole{p}{:}\DUrole{w}{ }\DUrole{n}{str}\DUrole{w}{ }\DUrole{o}{=}\DUrole{w}{ }\DUrole{default_value}{\textquotesingle{}9\textquotesingle{}}}\sphinxparamcomma \sphinxparam{\DUrole{n}{key\_global\_speed\_down}\DUrole{p}{:}\DUrole{w}{ }\DUrole{n}{str}\DUrole{w}{ }\DUrole{o}{=}\DUrole{w}{ }\DUrole{default_value}{\textquotesingle{}7\textquotesingle{}}}\sphinxparamcomma \sphinxparam{\DUrole{n}{volume\_step}\DUrole{p}{:}\DUrole{w}{ }\DUrole{n}{int}\DUrole{w}{ }\DUrole{o}{=}\DUrole{w}{ }\DUrole{default_value}{4}}\sphinxparamcomma \sphinxparam{\DUrole{n}{speed\_factor\_up}\DUrole{p}{:}\DUrole{w}{ }\DUrole{n}{float}\DUrole{w}{ }\DUrole{o}{=}\DUrole{w}{ }\DUrole{default_value}{1.25}}\sphinxparamcomma \sphinxparam{\DUrole{n}{speed\_factor\_down}\DUrole{p}{:}\DUrole{w}{ }\DUrole{n}{float}\DUrole{w}{ }\DUrole{o}{=}\DUrole{w}{ }\DUrole{default_value}{0.8}}}
{{ $\rightarrow$ None}}
\pysigstopsignatures
\sphinxAtStartPar
Initialise le ModeManager.

\sphinxAtStartPar
Démarre avec le premier mode (index 0) :
\sphinxhyphen{} annonce vocale du mode,
\sphinxhyphen{} appel de on\_enter() sur ce mode.

\sphinxAtStartPar
Si un keypad est fourni, il est démarré immédiatement et tourne
en continu pour toute la durée de vie de l’assistant.

\end{fulllineitems}

\index{current\_index (propriété core.mode\_manager.ModeManager)@\spxentry{current\_index}\spxextra{propriété core.mode\_manager.ModeManager}}

\begin{fulllineitems}
\phantomsection\label{\detokenize{api/core:core.mode_manager.ModeManager.current_index}}
\pysigstartsignatures
\pysigline
{\sphinxbfcode{\sphinxupquote{\DUrole{k}{property}\DUrole{w}{ }}}\sphinxbfcode{\sphinxupquote{current\_index}}\sphinxbfcode{\sphinxupquote{\DUrole{p}{:}\DUrole{w}{ }int}}}
\pysigstopsignatures
\sphinxAtStartPar
Index du mode courant (0\sphinxhyphen{}based).

\end{fulllineitems}

\index{current\_mode (propriété core.mode\_manager.ModeManager)@\spxentry{current\_mode}\spxextra{propriété core.mode\_manager.ModeManager}}

\begin{fulllineitems}
\phantomsection\label{\detokenize{api/core:core.mode_manager.ModeManager.current_mode}}
\pysigstartsignatures
\pysigline
{\sphinxbfcode{\sphinxupquote{\DUrole{k}{property}\DUrole{w}{ }}}\sphinxbfcode{\sphinxupquote{current\_mode}}\sphinxbfcode{\sphinxupquote{\DUrole{p}{:}\DUrole{w}{ }{\hyperref[\detokenize{api/core:core.mode_base.Mode}]{\sphinxcrossref{Mode}}}}}}
\pysigstopsignatures
\sphinxAtStartPar
Instance du mode courant.

\end{fulllineitems}

\index{handle\_key\_pressed() (méthode core.mode\_manager.ModeManager)@\spxentry{handle\_key\_pressed()}\spxextra{méthode core.mode\_manager.ModeManager}}

\begin{fulllineitems}
\phantomsection\label{\detokenize{api/core:core.mode_manager.ModeManager.handle_key_pressed}}
\pysigstartsignatures
\pysiglinewithargsret
{\sphinxbfcode{\sphinxupquote{handle\_key\_pressed}}}
{\sphinxparam{\DUrole{n}{key}\DUrole{p}{:}\DUrole{w}{ }\DUrole{n}{str}}}
{{ $\rightarrow$ None}}
\pysigstopsignatures
\sphinxAtStartPar
Point d’entrée unique pour toutes les touches du clavier.


\paragraph{Traitement}
\label{\detokenize{api/core:traitement}}\begin{enumerate}
\sphinxsetlistlabels{\arabic}{enumi}{enumii}{}{.}%
\item {} 
\sphinxAtStartPar
Si la touche est « globale » \(\rightarrow\) action immédiate (volume/vitesse)
+ bip de feedback. Ne parvient JAMAIS au mode courant.

\item {} 
\sphinxAtStartPar
Sinon \(\rightarrow\) délégation à self.\_current\_mode.on\_keypad\_key(key).
Si le mode ne surcharge pas cette méthode, c’est un no\sphinxhyphen{}op.

\end{enumerate}

\end{fulllineitems}

\index{set\_index() (méthode core.mode\_manager.ModeManager)@\spxentry{set\_index()}\spxextra{méthode core.mode\_manager.ModeManager}}

\begin{fulllineitems}
\phantomsection\label{\detokenize{api/core:core.mode_manager.ModeManager.set_index}}
\pysigstartsignatures
\pysiglinewithargsret
{\sphinxbfcode{\sphinxupquote{set\_index}}}
{\sphinxparam{\DUrole{n}{new\_index}\DUrole{p}{:}\DUrole{w}{ }\DUrole{n}{int}}\sphinxparamcomma \sphinxparam{\DUrole{n}{direction}\DUrole{p}{:}\DUrole{w}{ }\DUrole{n}{str\DUrole{w}{ }\DUrole{p}{|}\DUrole{w}{ }None}\DUrole{w}{ }\DUrole{o}{=}\DUrole{w}{ }\DUrole{default_value}{None}}}
{{ $\rightarrow$ None}}
\pysigstopsignatures
\sphinxAtStartPar
Change le mode actif.

\sphinxAtStartPar
Appelé par RotarySelector lorsqu’une rotation stable est détectée.


\paragraph{Paramètres}
\label{\detokenize{api/core:id2}}\begin{description}
\sphinxlineitem{new\_index :}
\sphinxAtStartPar
Nouvel index (0 ≤ index \textless{} len(modes)).

\sphinxlineitem{direction :}
\sphinxAtStartPar
Sens de rotation (« sens horaire » / « sens inverse ») — utilisé
pour les logs. Pas d’effet fonctionnel pour l’instant.

\end{description}


\paragraph{Comportement}
\label{\detokenize{api/core:comportement}}\begin{enumerate}
\sphinxsetlistlabels{\arabic}{enumi}{enumii}{}{.}%
\item {} 
\sphinxAtStartPar
Coupe tout audio en cours (stop\_all\_audio).

\item {} 
\sphinxAtStartPar
on\_exit() sur l’ancien mode.

\item {} 
\sphinxAtStartPar
Annonce le nouveau mode (nom mis en queue en premier).

\item {} 
\sphinxAtStartPar
on\_enter() sur le nouveau mode.

\end{enumerate}

\end{fulllineitems}

\index{handle\_short\_press() (méthode core.mode\_manager.ModeManager)@\spxentry{handle\_short\_press()}\spxextra{méthode core.mode\_manager.ModeManager}}

\begin{fulllineitems}
\phantomsection\label{\detokenize{api/core:core.mode_manager.ModeManager.handle_short_press}}
\pysigstartsignatures
\pysiglinewithargsret
{\sphinxbfcode{\sphinxupquote{handle\_short\_press}}}
{}
{{ $\rightarrow$ None}}
\pysigstopsignatures
\sphinxAtStartPar
Appui court sur le bouton rotatif.

\sphinxAtStartPar
Délégué au mode courant via on\_short\_press().

\end{fulllineitems}

\index{handle\_long\_press() (méthode core.mode\_manager.ModeManager)@\spxentry{handle\_long\_press()}\spxextra{méthode core.mode\_manager.ModeManager}}

\begin{fulllineitems}
\phantomsection\label{\detokenize{api/core:core.mode_manager.ModeManager.handle_long_press}}
\pysigstartsignatures
\pysiglinewithargsret
{\sphinxbfcode{\sphinxupquote{handle\_long\_press}}}
{}
{{ $\rightarrow$ None}}
\pysigstopsignatures
\sphinxAtStartPar
Appui long sur le bouton rotatif.

\sphinxAtStartPar
Délégué au mode courant via on\_long\_press().

\end{fulllineitems}

\index{handle\_double\_press() (méthode core.mode\_manager.ModeManager)@\spxentry{handle\_double\_press()}\spxextra{méthode core.mode\_manager.ModeManager}}

\begin{fulllineitems}
\phantomsection\label{\detokenize{api/core:core.mode_manager.ModeManager.handle_double_press}}
\pysigstartsignatures
\pysiglinewithargsret
{\sphinxbfcode{\sphinxupquote{handle\_double\_press}}}
{}
{{ $\rightarrow$ None}}
\pysigstopsignatures
\sphinxAtStartPar
Double appui rapide sur le bouton rotatif.

\sphinxAtStartPar
Délégué au mode courant via on\_double\_press().

\end{fulllineitems}


\end{fulllineitems}



\subsection{core.mode\_registry module}
\label{\detokenize{api/core:module-core.mode_registry}}\label{\detokenize{api/core:core-mode-registry-module}}\index{module@\spxentry{module}!core.mode\_registry@\spxentry{core.mode\_registry}}\index{core.mode\_registry@\spxentry{core.mode\_registry}!module@\spxentry{module}}

\subsubsection{core.mode\_registry}
\label{\detokenize{api/core:core-mode-registry}}
\sphinxAtStartPar
Ce module définit le \sphinxstylestrong{registre des modes} disponibles dans l’application.


\paragraph{Pourquoi un registre ?}
\label{\detokenize{api/core:pourquoi-un-registre}}
\sphinxAtStartPar
On veut éviter de modifier \sphinxtitleref{main.py} à chaque ajout de mode.

\sphinxAtStartPar
Avec un registre :
\sphinxhyphen{} \sphinxtitleref{config/settings.toml} décide quels modes sont actifs et dans quel ordre.
\sphinxhyphen{} \sphinxtitleref{main.py} se contente de :
\begin{itemize}
\item {} 
\sphinxAtStartPar
lire la config,

\item {} 
\sphinxAtStartPar
instancier les modes via ce registre.

\end{itemize}

\sphinxAtStartPar
Cela rend l’architecture :
\sphinxhyphen{} \sphinxstylestrong{extensible} (ajouter un mode = 1 import + 1 ligne dans le registre),
\sphinxhyphen{} \sphinxstylestrong{configurable} (activer/désactiver/re\sphinxhyphen{}ordonner via TOML),
\sphinxhyphen{} \sphinxstylestrong{propre} (pas de if/else ou de liste codée en dur dans main).


\paragraph{Contrat}
\label{\detokenize{api/core:contrat}}\begin{itemize}
\item {} 
\sphinxAtStartPar
Les clés (ex: « datetime ») doivent correspondre aux valeurs déclarées
dans \sphinxtitleref{{[}modes{]}.enabled} du fichier \sphinxtitleref{config/config.toml}.

\item {} 
\sphinxAtStartPar
Chaque valeur du registre est une \sphinxstylestrong{classe} dérivée de \sphinxtitleref{Mode}
(et non une instance).

\end{itemize}


\paragraph{Exemple}
\label{\detokenize{api/core:exemple}}
\sphinxAtStartPar
Dans \sphinxtitleref{config/config.toml} :

\sphinxAtStartPar
{[}modes{]}
enabled = {[}« datetime », « dummy », « multimeter »{]}

\sphinxAtStartPar
Alors \sphinxtitleref{main.py} instanciera dans cet ordre :
ModeDateHeure, ModeDummy, ModeMultimetre.


\subsection{core.speaker module}
\label{\detokenize{api/core:module-core.speaker}}\label{\detokenize{api/core:core-speaker-module}}\index{module@\spxentry{module}!core.speaker@\spxentry{core.speaker}}\index{core.speaker@\spxentry{core.speaker}!module@\spxentry{module}}

\subsubsection{core.speaker}
\label{\detokenize{api/core:core-speaker}}
\sphinxAtStartPar
Sortie audio centrale de l’assistant — deux canaux indépendants :
\begin{enumerate}
\sphinxsetlistlabels{\arabic}{enumi}{enumii}{}{.}%
\item {} 
\sphinxAtStartPar
CANAL ANNONCES  (say / speak)
─────────────────────────────
\sphinxhyphen{} File d’attente (Queue) pour ordonner les annonces.
\sphinxhyphen{} Chaque annonce = WAV Piper joué par un sous\sphinxhyphen{}processus mplayer
\begin{quote}

\sphinxAtStartPar
BLOQUANT (non\sphinxhyphen{}slave). Le thread queue attend la fin réelle du
fichier avant de passer au suivant.
\end{quote}
\begin{itemize}
\item {} 
\sphinxAtStartPar
Interruptible instantanément via stop\_all\_audio() \(\rightarrow\)
terminate() sur le sous\sphinxhyphen{}processus en cours.

\end{itemize}

\item {} 
\sphinxAtStartPar
CANAL LECTURE LONGUE  (play / PlaybackHandle)
──────────────────────────────────────────────
\sphinxhyphen{} mplayer en mode slave (\sphinxhyphen{}slave \sphinxhyphen{}idle) pour contrôle fin :
\begin{quote}

\sphinxAtStartPar
pause, reprise, seek, vitesse.
\end{quote}
\begin{itemize}
\item {} 
\sphinxAtStartPar
Utilisé exclusivement par ModeReadingMachine.

\item {} 
\sphinxAtStartPar
Indépendant de la queue.

\end{itemize}

\end{enumerate}
\begin{description}
\sphinxlineitem{VOLUME GLOBAL}\begin{itemize}
\item {} 
\sphinxAtStartPar
Contrôlé via amixer (ALSA) \(\rightarrow\) affecte tout (annonces + lecture).

\item {} 
\sphinxAtStartPar
set\_volume\_global() est appelable depuis n’importe quel mode.

\end{itemize}

\sphinxlineitem{ISOLATION DES MODES}\begin{itemize}
\item {} 
\sphinxAtStartPar
stop\_all\_audio() : vide la queue + tue le sous\sphinxhyphen{}processus d’annonce
en cours + envoie stop au player long\sphinxhyphen{}form.

\item {} 
\sphinxAtStartPar
Appelé par ModeManager avant chaque changement de mode.

\end{itemize}

\end{description}
\index{PlaybackHandle (classe dans core.speaker)@\spxentry{PlaybackHandle}\spxextra{classe dans core.speaker}}

\begin{fulllineitems}
\phantomsection\label{\detokenize{api/core:core.speaker.PlaybackHandle}}
\pysigstartsignatures
\pysiglinewithargsret
{\sphinxbfcode{\sphinxupquote{\DUrole{k}{class}\DUrole{w}{ }}}\sphinxcode{\sphinxupquote{core.speaker.}}\sphinxbfcode{\sphinxupquote{PlaybackHandle}}}
{\sphinxparam{\DUrole{n}{\_player}\DUrole{p}{:}\DUrole{w}{ }\DUrole{n}{\_AudioPlayer}}}
{}
\pysigstopsignatures
\sphinxAtStartPar
Bases : \sphinxcode{\sphinxupquote{object}}

\sphinxAtStartPar
Contrôle d’une lecture longue (machine à lire).

\sphinxAtStartPar
Méthodes disponibles : pause, resume, stop, seek, set\_speed.
\index{pause() (méthode core.speaker.PlaybackHandle)@\spxentry{pause()}\spxextra{méthode core.speaker.PlaybackHandle}}

\begin{fulllineitems}
\phantomsection\label{\detokenize{api/core:core.speaker.PlaybackHandle.pause}}
\pysigstartsignatures
\pysiglinewithargsret
{\sphinxbfcode{\sphinxupquote{pause}}}
{}
{{ $\rightarrow$ None}}
\pysigstopsignatures
\end{fulllineitems}

\index{resume() (méthode core.speaker.PlaybackHandle)@\spxentry{resume()}\spxextra{méthode core.speaker.PlaybackHandle}}

\begin{fulllineitems}
\phantomsection\label{\detokenize{api/core:core.speaker.PlaybackHandle.resume}}
\pysigstartsignatures
\pysiglinewithargsret
{\sphinxbfcode{\sphinxupquote{resume}}}
{}
{{ $\rightarrow$ None}}
\pysigstopsignatures
\end{fulllineitems}

\index{stop() (méthode core.speaker.PlaybackHandle)@\spxentry{stop()}\spxextra{méthode core.speaker.PlaybackHandle}}

\begin{fulllineitems}
\phantomsection\label{\detokenize{api/core:core.speaker.PlaybackHandle.stop}}
\pysigstartsignatures
\pysiglinewithargsret
{\sphinxbfcode{\sphinxupquote{stop}}}
{}
{{ $\rightarrow$ None}}
\pysigstopsignatures
\end{fulllineitems}

\index{seek() (méthode core.speaker.PlaybackHandle)@\spxentry{seek()}\spxextra{méthode core.speaker.PlaybackHandle}}

\begin{fulllineitems}
\phantomsection\label{\detokenize{api/core:core.speaker.PlaybackHandle.seek}}
\pysigstartsignatures
\pysiglinewithargsret
{\sphinxbfcode{\sphinxupquote{seek}}}
{\sphinxparam{\DUrole{n}{seconds}\DUrole{p}{:}\DUrole{w}{ }\DUrole{n}{int}}}
{{ $\rightarrow$ None}}
\pysigstopsignatures
\end{fulllineitems}

\index{set\_speed() (méthode core.speaker.PlaybackHandle)@\spxentry{set\_speed()}\spxextra{méthode core.speaker.PlaybackHandle}}

\begin{fulllineitems}
\phantomsection\label{\detokenize{api/core:core.speaker.PlaybackHandle.set_speed}}
\pysigstartsignatures
\pysiglinewithargsret
{\sphinxbfcode{\sphinxupquote{set\_speed}}}
{\sphinxparam{\DUrole{n}{speed}\DUrole{p}{:}\DUrole{w}{ }\DUrole{n}{float}}}
{{ $\rightarrow$ None}}
\pysigstopsignatures
\end{fulllineitems}


\end{fulllineitems}

\index{Speaker (classe dans core.speaker)@\spxentry{Speaker}\spxextra{classe dans core.speaker}}

\begin{fulllineitems}
\phantomsection\label{\detokenize{api/core:core.speaker.Speaker}}
\pysigstartsignatures
\pysigline
{\sphinxbfcode{\sphinxupquote{\DUrole{k}{class}\DUrole{w}{ }}}\sphinxcode{\sphinxupquote{core.speaker.}}\sphinxbfcode{\sphinxupquote{Speaker}}}
\pysigstopsignatures
\sphinxAtStartPar
Bases : \sphinxcode{\sphinxupquote{object}}

\sphinxAtStartPar
Synthèse vocale centrale de l’assistant.


\subsubsection{Interface publique}
\label{\detokenize{api/core:interface-publique}}
\sphinxAtStartPar
say(text, blocking=False)      \(\rightarrow\) annonce courte (queue)
speak(text)                    \(\rightarrow\) alias say() non\sphinxhyphen{}bloquant
play(path) \(\rightarrow\) PlaybackHandle    \(\rightarrow\) lecture longue (machine à lire)
play\_sound(path, blocking)     \(\rightarrow\) son court hors\sphinxhyphen{}queue (aplay)
stop\_all\_audio()               \(\rightarrow\) coupe tout (changement de mode)
set\_volume\_global(percent)     \(\rightarrow\) volume ALSA global
set\_global\_speed(speed)        \(\rightarrow\) vitesse globale (0.5 \textendash{} 2.0)
close()                        \(\rightarrow\) arrêt propre
\index{synthesize\_to\_file() (méthode core.speaker.Speaker)@\spxentry{synthesize\_to\_file()}\spxextra{méthode core.speaker.Speaker}}

\begin{fulllineitems}
\phantomsection\label{\detokenize{api/core:core.speaker.Speaker.synthesize_to_file}}
\pysigstartsignatures
\pysiglinewithargsret
{\sphinxbfcode{\sphinxupquote{synthesize\_to\_file}}}
{\sphinxparam{\DUrole{n}{text}\DUrole{p}{:}\DUrole{w}{ }\DUrole{n}{str}}\sphinxparamcomma \sphinxparam{\DUrole{n}{output\_path}\DUrole{p}{:}\DUrole{w}{ }\DUrole{n}{str}}}
{{ $\rightarrow$ None}}
\pysigstopsignatures
\sphinxAtStartPar
Synthétise “text” via Piper et écrit le WAV à “output\_path”.

\sphinxAtStartPar
Méthode publique utilisée par les modules extérieurs (ex: pipeline
machine à lire) qui veulent utiliser la voix Piper du Speaker sans
charger un second modèle en mémoire.


\paragraph{Paramètres}
\label{\detokenize{api/core:id3}}\begin{description}
\sphinxlineitem{text :}
\sphinxAtStartPar
Texte à synthétiser.

\sphinxlineitem{output\_path :}
\sphinxAtStartPar
Chemin absolu du fichier WAV de sortie.
Le dossier parent doit exister.

\end{description}
\begin{quote}\begin{description}
\sphinxlineitem{raises ValueError}
\sphinxAtStartPar
Si le texte est vide.

\end{description}\end{quote}

\end{fulllineitems}

\index{play\_notification() (méthode core.speaker.Speaker)@\spxentry{play\_notification()}\spxextra{méthode core.speaker.Speaker}}

\begin{fulllineitems}
\phantomsection\label{\detokenize{api/core:core.speaker.Speaker.play_notification}}
\pysigstartsignatures
\pysiglinewithargsret
{\sphinxbfcode{\sphinxupquote{play\_notification}}}
{\sphinxparam{\DUrole{n}{blocking}\DUrole{p}{:}\DUrole{w}{ }\DUrole{n}{bool}\DUrole{w}{ }\DUrole{o}{=}\DUrole{w}{ }\DUrole{default_value}{False}}}
{{ $\rightarrow$ bool}}
\pysigstopsignatures
\sphinxAtStartPar
Joue le son de notification de volume/vitesse.
\begin{quote}\begin{description}
\sphinxlineitem{Renvoie}
\sphinxAtStartPar
True si un fichier de notification a été trouvé et tenté,
False sinon.

\sphinxlineitem{Type renvoyé}
\sphinxAtStartPar
bool

\end{description}\end{quote}

\end{fulllineitems}

\index{beep() (méthode core.speaker.Speaker)@\spxentry{beep()}\spxextra{méthode core.speaker.Speaker}}

\begin{fulllineitems}
\phantomsection\label{\detokenize{api/core:core.speaker.Speaker.beep}}
\pysigstartsignatures
\pysiglinewithargsret
{\sphinxbfcode{\sphinxupquote{beep}}}
{\sphinxparam{\DUrole{n}{direction}\DUrole{p}{:}\DUrole{w}{ }\DUrole{n}{str}\DUrole{w}{ }\DUrole{o}{=}\DUrole{w}{ }\DUrole{default_value}{\textquotesingle{}up\textquotesingle{}}}}
{{ $\rightarrow$ None}}
\pysigstopsignatures
\sphinxAtStartPar
Alias de compatibilité : joue le son de notification.

\end{fulllineitems}

\index{say() (méthode core.speaker.Speaker)@\spxentry{say()}\spxextra{méthode core.speaker.Speaker}}

\begin{fulllineitems}
\phantomsection\label{\detokenize{api/core:core.speaker.Speaker.say}}
\pysigstartsignatures
\pysiglinewithargsret
{\sphinxbfcode{\sphinxupquote{say}}}
{\sphinxparam{\DUrole{n}{text}\DUrole{p}{:}\DUrole{w}{ }\DUrole{n}{str}}\sphinxparamcomma \sphinxparam{\DUrole{n}{blocking}\DUrole{p}{:}\DUrole{w}{ }\DUrole{n}{bool}\DUrole{w}{ }\DUrole{o}{=}\DUrole{w}{ }\DUrole{default_value}{False}}}
{{ $\rightarrow$ None}}
\pysigstopsignatures
\sphinxAtStartPar
Annonce courte — synthétisée par Piper, lue via la queue.


\paragraph{Paramètres}
\label{\detokenize{api/core:id4}}\begin{description}
\sphinxlineitem{text :}
\sphinxAtStartPar
Texte à synthétiser.

\sphinxlineitem{blocking :}
\sphinxAtStartPar
Si True, attend que CETTE annonce (et toutes celles déjà
en queue avant elle) soit terminée avant de retourner.
Utiliser pour s’assurer qu’une annonce est entendue avant
de poursuivre l’initialisation.

\end{description}

\end{fulllineitems}

\index{speak() (méthode core.speaker.Speaker)@\spxentry{speak()}\spxextra{méthode core.speaker.Speaker}}

\begin{fulllineitems}
\phantomsection\label{\detokenize{api/core:core.speaker.Speaker.speak}}
\pysigstartsignatures
\pysiglinewithargsret
{\sphinxbfcode{\sphinxupquote{speak}}}
{\sphinxparam{\DUrole{n}{text}\DUrole{p}{:}\DUrole{w}{ }\DUrole{n}{str}}}
{{ $\rightarrow$ None}}
\pysigstopsignatures
\sphinxAtStartPar
Alias non\sphinxhyphen{}bloquant de say() — compatibilité avec l’ancien code.

\end{fulllineitems}

\index{play() (méthode core.speaker.Speaker)@\spxentry{play()}\spxextra{méthode core.speaker.Speaker}}

\begin{fulllineitems}
\phantomsection\label{\detokenize{api/core:core.speaker.Speaker.play}}
\pysigstartsignatures
\pysiglinewithargsret
{\sphinxbfcode{\sphinxupquote{play}}}
{\sphinxparam{\DUrole{n}{path}\DUrole{p}{:}\DUrole{w}{ }\DUrole{n}{str}}}
{{ $\rightarrow$ {\hyperref[\detokenize{api/core:core.speaker.PlaybackHandle}]{\sphinxcrossref{PlaybackHandle}}}}}
\pysigstopsignatures
\sphinxAtStartPar
Lecture longue (machine à lire).

\sphinxAtStartPar
Coupe les annonces en cours et lance la lecture via mplayer slave.
Retourne un PlaybackHandle pour contrôle fin.


\paragraph{Paramètres}
\label{\detokenize{api/core:id5}}\begin{description}
\sphinxlineitem{path :}
\sphinxAtStartPar
Chemin vers le fichier WAV à lire.

\end{description}

\end{fulllineitems}

\index{play\_sound() (méthode core.speaker.Speaker)@\spxentry{play\_sound()}\spxextra{méthode core.speaker.Speaker}}

\begin{fulllineitems}
\phantomsection\label{\detokenize{api/core:core.speaker.Speaker.play_sound}}
\pysigstartsignatures
\pysiglinewithargsret
{\sphinxbfcode{\sphinxupquote{play\_sound}}}
{\sphinxparam{\DUrole{n}{path}\DUrole{p}{:}\DUrole{w}{ }\DUrole{n}{str}}\sphinxparamcomma \sphinxparam{\DUrole{n}{blocking}\DUrole{p}{:}\DUrole{w}{ }\DUrole{n}{bool}\DUrole{w}{ }\DUrole{o}{=}\DUrole{w}{ }\DUrole{default_value}{False}}}
{{ $\rightarrow$ None}}
\pysigstopsignatures
\sphinxAtStartPar
Son court hors\sphinxhyphen{}queue (effets sonores : shutter, erreur, etc.).

\sphinxAtStartPar
Utilise aplay (ALSA, très léger) — aucune interférence avec
la queue ni le player long\sphinxhyphen{}form.

\end{fulllineitems}

\index{stop\_all\_audio() (méthode core.speaker.Speaker)@\spxentry{stop\_all\_audio()}\spxextra{méthode core.speaker.Speaker}}

\begin{fulllineitems}
\phantomsection\label{\detokenize{api/core:core.speaker.Speaker.stop_all_audio}}
\pysigstartsignatures
\pysiglinewithargsret
{\sphinxbfcode{\sphinxupquote{stop\_all\_audio}}}
{}
{{ $\rightarrow$ None}}
\pysigstopsignatures
\sphinxAtStartPar
Coupe tout l’audio immédiatement.


\paragraph{Actions}
\label{\detokenize{api/core:actions}}\begin{enumerate}
\sphinxsetlistlabels{\arabic}{enumi}{enumii}{}{.}%
\item {} 
\sphinxAtStartPar
Vide la queue (et supprime les WAVs en attente).

\item {} 
\sphinxAtStartPar
Tue le sous\sphinxhyphen{}processus d’annonce en cours (terminate).

\item {} 
\sphinxAtStartPar
Arrête le player long\sphinxhyphen{}form (mplayer slave \(\rightarrow\) stop).

\end{enumerate}

\sphinxAtStartPar
Appelé par ModeManager à chaque changement de mode.

\end{fulllineitems}

\index{set\_volume\_global() (méthode core.speaker.Speaker)@\spxentry{set\_volume\_global()}\spxextra{méthode core.speaker.Speaker}}

\begin{fulllineitems}
\phantomsection\label{\detokenize{api/core:core.speaker.Speaker.set_volume_global}}
\pysigstartsignatures
\pysiglinewithargsret
{\sphinxbfcode{\sphinxupquote{set\_volume\_global}}}
{\sphinxparam{\DUrole{n}{percent}\DUrole{p}{:}\DUrole{w}{ }\DUrole{n}{int}}\sphinxparamcomma \sphinxparam{\DUrole{n}{mixer}\DUrole{p}{:}\DUrole{w}{ }\DUrole{n}{str\DUrole{w}{ }\DUrole{p}{|}\DUrole{w}{ }None}\DUrole{w}{ }\DUrole{o}{=}\DUrole{w}{ }\DUrole{default_value}{None}}}
{{ $\rightarrow$ None}}
\pysigstopsignatures
\sphinxAtStartPar
Règle le volume ALSA global (0 \textendash{} 100).

\sphinxAtStartPar
Affecte toutes les sorties audio (annonces + lecture longue).
Le mixer est détecté automatiquement si non précisé.
La valeur est mémorisée dans \_current\_volume pour permettre
aux modes de lire le niveau courant via get\_volume().

\end{fulllineitems}

\index{get\_volume() (méthode core.speaker.Speaker)@\spxentry{get\_volume()}\spxextra{méthode core.speaker.Speaker}}

\begin{fulllineitems}
\phantomsection\label{\detokenize{api/core:core.speaker.Speaker.get_volume}}
\pysigstartsignatures
\pysiglinewithargsret
{\sphinxbfcode{\sphinxupquote{get\_volume}}}
{}
{{ $\rightarrow$ int}}
\pysigstopsignatures
\sphinxAtStartPar
Retourne le volume global courant (0 \textendash{} 100).

\sphinxAtStartPar
Valeur mémorisée à chaque appel de set\_volume\_global().
Utilisée par les modes pour effectuer des incréments relatifs
(ex : volume +4, volume \sphinxhyphen{}4) sans interroger ALSA.

\end{fulllineitems}

\index{get\_speed() (méthode core.speaker.Speaker)@\spxentry{get\_speed()}\spxextra{méthode core.speaker.Speaker}}

\begin{fulllineitems}
\phantomsection\label{\detokenize{api/core:core.speaker.Speaker.get_speed}}
\pysigstartsignatures
\pysiglinewithargsret
{\sphinxbfcode{\sphinxupquote{get\_speed}}}
{}
{{ $\rightarrow$ float}}
\pysigstopsignatures
\sphinxAtStartPar
Retourne la vitesse de lecture globale courante (0.5 \textendash{} 2.0).

\sphinxAtStartPar
Valeur mémorisée à chaque appel de set\_global\_speed().
Utilisée par les modes pour effectuer des incréments relatifs.

\end{fulllineitems}

\index{set\_global\_speed() (méthode core.speaker.Speaker)@\spxentry{set\_global\_speed()}\spxextra{méthode core.speaker.Speaker}}

\begin{fulllineitems}
\phantomsection\label{\detokenize{api/core:core.speaker.Speaker.set_global_speed}}
\pysigstartsignatures
\pysiglinewithargsret
{\sphinxbfcode{\sphinxupquote{set\_global\_speed}}}
{\sphinxparam{\DUrole{n}{speed}\DUrole{p}{:}\DUrole{w}{ }\DUrole{n}{float}}}
{{ $\rightarrow$ None}}
\pysigstopsignatures
\sphinxAtStartPar
Règle la vitesse globale (0.5 \textendash{} 2.0).

\sphinxAtStartPar
S’applique :
\sphinxhyphen{} Aux prochaines annonces (utilisé au moment du subprocess).
\sphinxhyphen{} Au player long\sphinxhyphen{}form immédiatement (commande mplayer).

\end{fulllineitems}

\index{close() (méthode core.speaker.Speaker)@\spxentry{close()}\spxextra{méthode core.speaker.Speaker}}

\begin{fulllineitems}
\phantomsection\label{\detokenize{api/core:core.speaker.Speaker.close}}
\pysigstartsignatures
\pysiglinewithargsret
{\sphinxbfcode{\sphinxupquote{close}}}
{}
{{ $\rightarrow$ None}}
\pysigstopsignatures
\sphinxAtStartPar
Arrêt propre du Speaker.

\sphinxAtStartPar
Interrompt les lectures en cours, vide la queue, supprime
les fichiers temporaires.

\end{fulllineitems}


\end{fulllineitems}



\subsection{Module contents}
\label{\detokenize{api/core:module-core}}\label{\detokenize{api/core:module-contents}}\index{module@\spxentry{module}!core@\spxentry{core}}\index{core@\spxentry{core}!module@\spxentry{module}}
\sphinxstepscope


\section{hardware package}
\label{\detokenize{api/hardware:hardware-package}}\label{\detokenize{api/hardware::doc}}

\subsection{Subpackages}
\label{\detokenize{api/hardware:subpackages}}
\sphinxstepscope


\subsubsection{hardware.devices package}
\label{\detokenize{api/hardware.devices:hardware-devices-package}}\label{\detokenize{api/hardware.devices::doc}}

\paragraph{Submodules}
\label{\detokenize{api/hardware.devices:submodules}}

\paragraph{hardware.devices.caliper module}
\label{\detokenize{api/hardware.devices:module-hardware.devices.caliper}}\label{\detokenize{api/hardware.devices:hardware-devices-caliper-module}}\index{module@\spxentry{module}!hardware.devices.caliper@\spxentry{hardware.devices.caliper}}\index{hardware.devices.caliper@\spxentry{hardware.devices.caliper}!module@\spxentry{module}}
\sphinxAtStartPar
Driver pour le pied à coulisse numérique (version NRF24).


\subparagraph{Vue d’ensemble}
\label{\detokenize{api/hardware.devices:vue-d-ensemble}}
\sphinxAtStartPar
Ce module assure le lien entre le transport radio brut et la logique métier :
\begin{enumerate}
\sphinxsetlistlabels{\arabic}{enumi}{enumii}{}{.}%
\item {} \begin{itemize}
\item {} 
\sphinxAtStartPar
Structure les données décodées : valeur, unité (mm/pouce) et signe.

\end{itemize}

\item {} \begin{itemize}
\item {} 
\sphinxAtStartPar
Initialise le transport \sphinxtitleref{Nrf24Transport}.

\item {} 
\sphinxAtStartPar
Décode la structure binaire \sphinxtitleref{MessageStruct} envoyée par l’ATtiny85.

\item {} 
\sphinxAtStartPar
Gère les notifications d’état (prêt/erreur) vers le mode actif.

\end{itemize}

\end{enumerate}


\subparagraph{Protocole Binaire (Caliper\_TX\_24b\_ATtiny85)}
\label{\detokenize{api/hardware.devices:protocole-binaire-caliper-tx-24b-attiny85}}
\sphinxAtStartPar
La trame reçue via NRF24 comporte au minimum 6 octets (Little Endian) :
\sphinxhyphen{} Octets 0 à 3 : float (4 octets) \sphinxhyphen{}\textgreater{} Valeur numérique brute.
\sphinxhyphen{} Octet 4      : bool (1 octet)  \sphinxhyphen{}\textgreater{} Unité (0 = mm, 1 = pouce).
\sphinxhyphen{} Octet 5      : bool (1 octet)  \sphinxhyphen{}\textgreater{} Signe (0 = plus, 1 = moins).
\index{CaliperData (classe dans hardware.devices.caliper)@\spxentry{CaliperData}\spxextra{classe dans hardware.devices.caliper}}

\begin{fulllineitems}
\phantomsection\label{\detokenize{api/hardware.devices:hardware.devices.caliper.CaliperData}}
\pysigstartsignatures
\pysiglinewithargsret
{\sphinxbfcode{\sphinxupquote{\DUrole{k}{class}\DUrole{w}{ }}}\sphinxcode{\sphinxupquote{hardware.devices.caliper.}}\sphinxbfcode{\sphinxupquote{CaliperData}}}
{\sphinxparam{\DUrole{n}{value}\DUrole{p}{:}\DUrole{w}{ }\DUrole{n}{float}}\sphinxparamcomma \sphinxparam{\DUrole{n}{unit\_name}\DUrole{p}{:}\DUrole{w}{ }\DUrole{n}{str}}\sphinxparamcomma \sphinxparam{\DUrole{n}{is\_negative}\DUrole{p}{:}\DUrole{w}{ }\DUrole{n}{bool}}}
{}
\pysigstopsignatures
\sphinxAtStartPar
Bases : \sphinxcode{\sphinxupquote{object}}

\sphinxAtStartPar
Représente une mesure décodée issue du pied à coulisse.


\subparagraph{Attributs}
\label{\detokenize{api/hardware.devices:attributs}}\begin{description}
\sphinxlineitem{value}{[}float{]}
\sphinxAtStartPar
Valeur numérique de la mesure (ex: 12.45).

\sphinxlineitem{unit\_name}{[}str{]}
\sphinxAtStartPar
Nom de l’unité (« millimètre » ou « pouce »).

\sphinxlineitem{is\_negative}{[}bool{]}
\sphinxAtStartPar
True si la mesure est négative, False sinon.

\end{description}
\index{value (attribut hardware.devices.caliper.CaliperData)@\spxentry{value}\spxextra{attribut hardware.devices.caliper.CaliperData}}

\begin{fulllineitems}
\phantomsection\label{\detokenize{api/hardware.devices:hardware.devices.caliper.CaliperData.value}}
\pysigstartsignatures
\pysigline
{\sphinxbfcode{\sphinxupquote{value}}\sphinxbfcode{\sphinxupquote{\DUrole{p}{:}\DUrole{w}{ }float}}}
\pysigstopsignatures
\end{fulllineitems}

\index{unit\_name (attribut hardware.devices.caliper.CaliperData)@\spxentry{unit\_name}\spxextra{attribut hardware.devices.caliper.CaliperData}}

\begin{fulllineitems}
\phantomsection\label{\detokenize{api/hardware.devices:hardware.devices.caliper.CaliperData.unit_name}}
\pysigstartsignatures
\pysigline
{\sphinxbfcode{\sphinxupquote{unit\_name}}\sphinxbfcode{\sphinxupquote{\DUrole{p}{:}\DUrole{w}{ }str}}}
\pysigstopsignatures
\end{fulllineitems}

\index{is\_negative (attribut hardware.devices.caliper.CaliperData)@\spxentry{is\_negative}\spxextra{attribut hardware.devices.caliper.CaliperData}}

\begin{fulllineitems}
\phantomsection\label{\detokenize{api/hardware.devices:hardware.devices.caliper.CaliperData.is_negative}}
\pysigstartsignatures
\pysigline
{\sphinxbfcode{\sphinxupquote{is\_negative}}\sphinxbfcode{\sphinxupquote{\DUrole{p}{:}\DUrole{w}{ }bool}}}
\pysigstopsignatures
\end{fulllineitems}


\end{fulllineitems}

\index{CaliperDriver (classe dans hardware.devices.caliper)@\spxentry{CaliperDriver}\spxextra{classe dans hardware.devices.caliper}}

\begin{fulllineitems}
\phantomsection\label{\detokenize{api/hardware.devices:hardware.devices.caliper.CaliperDriver}}
\pysigstartsignatures
\pysiglinewithargsret
{\sphinxbfcode{\sphinxupquote{\DUrole{k}{class}\DUrole{w}{ }}}\sphinxcode{\sphinxupquote{hardware.devices.caliper.}}\sphinxbfcode{\sphinxupquote{CaliperDriver}}}
{\sphinxparam{\DUrole{n}{on\_status}\DUrole{p}{:}\DUrole{w}{ }\DUrole{n}{Callable\DUrole{p}{{[}}\DUrole{p}{{[}}str\DUrole{p}{{]}}\DUrole{p}{,}\DUrole{w}{ }None\DUrole{p}{{]}}}}\sphinxparamcomma \sphinxparam{\DUrole{n}{on\_new\_measure}\DUrole{p}{:}\DUrole{w}{ }\DUrole{n}{Callable\DUrole{p}{{[}}\DUrole{p}{{[}}{\hyperref[\detokenize{api/hardware.devices:hardware.devices.caliper.CaliperData}]{\sphinxcrossref{CaliperData}}}\DUrole{p}{{]}}\DUrole{p}{,}\DUrole{w}{ }None\DUrole{p}{{]}}}}}
{}
\pysigstopsignatures
\sphinxAtStartPar
Bases : \sphinxcode{\sphinxupquote{object}}

\sphinxAtStartPar
Driver haut niveau pour le pied à coulisse sans fil.

\sphinxAtStartPar
Ce driver utilise le transport NRF24 pour recevoir des paquets et les
transformer en objets \sphinxtitleref{CaliperData} utilisables par l’assistant vocal.
\index{\_\_init\_\_() (méthode hardware.devices.caliper.CaliperDriver)@\spxentry{\_\_init\_\_()}\spxextra{méthode hardware.devices.caliper.CaliperDriver}}

\begin{fulllineitems}
\phantomsection\label{\detokenize{api/hardware.devices:hardware.devices.caliper.CaliperDriver.__init__}}
\pysigstartsignatures
\pysiglinewithargsret
{\sphinxbfcode{\sphinxupquote{\_\_init\_\_}}}
{\sphinxparam{\DUrole{n}{on\_status}\DUrole{p}{:}\DUrole{w}{ }\DUrole{n}{Callable\DUrole{p}{{[}}\DUrole{p}{{[}}str\DUrole{p}{{]}}\DUrole{p}{,}\DUrole{w}{ }None\DUrole{p}{{]}}}}\sphinxparamcomma \sphinxparam{\DUrole{n}{on\_new\_measure}\DUrole{p}{:}\DUrole{w}{ }\DUrole{n}{Callable\DUrole{p}{{[}}\DUrole{p}{{[}}{\hyperref[\detokenize{api/hardware.devices:hardware.devices.caliper.CaliperData}]{\sphinxcrossref{CaliperData}}}\DUrole{p}{{]}}\DUrole{p}{,}\DUrole{w}{ }None\DUrole{p}{{]}}}}}
{}
\pysigstopsignatures
\sphinxAtStartPar
Initialise le driver du pied à coulisse.


\subparagraph{Paramètres}
\label{\detokenize{api/hardware.devices:parametres}}\begin{description}
\sphinxlineitem{on\_status}{[}Callable{[}{[}str{]}, None{]}{]}
\sphinxAtStartPar
Callback pour signaler l’état du matériel (ex: « Pied à coulisse prêt »).

\sphinxlineitem{on\_new\_measure}{[}Callable{[}{[}CaliperData{]}, None{]}{]}
\sphinxAtStartPar
Callback appelé automatiquement à chaque nouvelle mesure reçue.

\end{description}

\end{fulllineitems}

\index{start() (méthode hardware.devices.caliper.CaliperDriver)@\spxentry{start()}\spxextra{méthode hardware.devices.caliper.CaliperDriver}}

\begin{fulllineitems}
\phantomsection\label{\detokenize{api/hardware.devices:hardware.devices.caliper.CaliperDriver.start}}
\pysigstartsignatures
\pysiglinewithargsret
{\sphinxbfcode{\sphinxupquote{start}}}
{}
{{ $\rightarrow$ None}}
\pysigstopsignatures
\sphinxAtStartPar
Démarre le transport radio et prépare la réception.

\sphinxAtStartPar
Cette méthode tente d’initialiser le module NRF24. Le résultat
est communiqué via le callback \sphinxtitleref{on\_status}.

\end{fulllineitems}

\index{stop() (méthode hardware.devices.caliper.CaliperDriver)@\spxentry{stop()}\spxextra{méthode hardware.devices.caliper.CaliperDriver}}

\begin{fulllineitems}
\phantomsection\label{\detokenize{api/hardware.devices:hardware.devices.caliper.CaliperDriver.stop}}
\pysigstartsignatures
\pysiglinewithargsret
{\sphinxbfcode{\sphinxupquote{stop}}}
{}
{{ $\rightarrow$ None}}
\pysigstopsignatures
\sphinxAtStartPar
Arrête proprement le transport et libère les ressources radio.

\end{fulllineitems}

\index{get\_last\_measure() (méthode hardware.devices.caliper.CaliperDriver)@\spxentry{get\_last\_measure()}\spxextra{méthode hardware.devices.caliper.CaliperDriver}}

\begin{fulllineitems}
\phantomsection\label{\detokenize{api/hardware.devices:hardware.devices.caliper.CaliperDriver.get_last_measure}}
\pysigstartsignatures
\pysiglinewithargsret
{\sphinxbfcode{\sphinxupquote{get\_last\_measure}}}
{}
{{ $\rightarrow$ {\hyperref[\detokenize{api/hardware.devices:hardware.devices.caliper.CaliperData}]{\sphinxcrossref{CaliperData}}}\DUrole{w}{ }\DUrole{p}{|}\DUrole{w}{ }None}}
\pysigstopsignatures
\sphinxAtStartPar
Retourne la toute dernière mesure mémorisée.
\begin{quote}\begin{description}
\sphinxlineitem{Renvoie}
\sphinxAtStartPar
La mesure si elle existe, sinon None.

\sphinxlineitem{Type renvoyé}
\sphinxAtStartPar
{\hyperref[\detokenize{api/hardware.devices:hardware.devices.caliper.CaliperData}]{\sphinxcrossref{CaliperData}}} or None

\end{description}\end{quote}

\end{fulllineitems}


\end{fulllineitems}



\paragraph{hardware.devices.esp32\_thermometer module}
\label{\detokenize{api/hardware.devices:module-hardware.devices.esp32_thermometer}}\label{\detokenize{api/hardware.devices:hardware-devices-esp32-thermometer-module}}\index{module@\spxentry{module}!hardware.devices.esp32\_thermometer@\spxentry{hardware.devices.esp32\_thermometer}}\index{hardware.devices.esp32\_thermometer@\spxentry{hardware.devices.esp32\_thermometer}!module@\spxentry{module}}

\subparagraph{hardware.devices.esp32\_thermometer}
\label{\detokenize{api/hardware.devices:hardware-devices-esp32-thermometer}}
\sphinxAtStartPar
Driver pour le thermomètre connecté ESP32.


\subparagraph{Vue d’ensemble}
\label{\detokenize{api/hardware.devices:id1}}
\sphinxAtStartPar
Ce module fournit :
1. \sphinxtitleref{ThermometerData} : structure de donnée propre (valeur float).
2. \sphinxtitleref{ThermometerDriver} : pilote haut niveau utilisant \sphinxtitleref{WifiClient}.

\sphinxAtStartPar
Il est conçu pour être utilisé par \sphinxtitleref{ModeThermometre}.
\index{ThermometerData (classe dans hardware.devices.esp32\_thermometer)@\spxentry{ThermometerData}\spxextra{classe dans hardware.devices.esp32\_thermometer}}

\begin{fulllineitems}
\phantomsection\label{\detokenize{api/hardware.devices:hardware.devices.esp32_thermometer.ThermometerData}}
\pysigstartsignatures
\pysiglinewithargsret
{\sphinxbfcode{\sphinxupquote{\DUrole{k}{class}\DUrole{w}{ }}}\sphinxcode{\sphinxupquote{hardware.devices.esp32\_thermometer.}}\sphinxbfcode{\sphinxupquote{ThermometerData}}}
{\sphinxparam{\DUrole{n}{value}\DUrole{p}{:}\DUrole{w}{ }\DUrole{n}{float}}\sphinxparamcomma \sphinxparam{\DUrole{n}{unit}\DUrole{p}{:}\DUrole{w}{ }\DUrole{n}{str}\DUrole{w}{ }\DUrole{o}{=}\DUrole{w}{ }\DUrole{default_value}{\textquotesingle{}Celsius\textquotesingle{}}}\sphinxparamcomma \sphinxparam{\DUrole{n}{timestamp}\DUrole{p}{:}\DUrole{w}{ }\DUrole{n}{float}\DUrole{w}{ }\DUrole{o}{=}\DUrole{w}{ }\DUrole{default_value}{0.0}}}
{}
\pysigstopsignatures
\sphinxAtStartPar
Bases : \sphinxcode{\sphinxupquote{object}}

\sphinxAtStartPar
Représente une mesure de température.
\index{value (attribut hardware.devices.esp32\_thermometer.ThermometerData)@\spxentry{value}\spxextra{attribut hardware.devices.esp32\_thermometer.ThermometerData}}

\begin{fulllineitems}
\phantomsection\label{\detokenize{api/hardware.devices:hardware.devices.esp32_thermometer.ThermometerData.value}}
\pysigstartsignatures
\pysigline
{\sphinxbfcode{\sphinxupquote{value}}\sphinxbfcode{\sphinxupquote{\DUrole{p}{:}\DUrole{w}{ }float}}}
\pysigstopsignatures
\end{fulllineitems}

\index{unit (attribut hardware.devices.esp32\_thermometer.ThermometerData)@\spxentry{unit}\spxextra{attribut hardware.devices.esp32\_thermometer.ThermometerData}}

\begin{fulllineitems}
\phantomsection\label{\detokenize{api/hardware.devices:hardware.devices.esp32_thermometer.ThermometerData.unit}}
\pysigstartsignatures
\pysigline
{\sphinxbfcode{\sphinxupquote{unit}}\sphinxbfcode{\sphinxupquote{\DUrole{p}{:}\DUrole{w}{ }str}}\sphinxbfcode{\sphinxupquote{\DUrole{w}{ }\DUrole{p}{=}\DUrole{w}{ }\textquotesingle{}Celsius\textquotesingle{}}}}
\pysigstopsignatures
\end{fulllineitems}

\index{timestamp (attribut hardware.devices.esp32\_thermometer.ThermometerData)@\spxentry{timestamp}\spxextra{attribut hardware.devices.esp32\_thermometer.ThermometerData}}

\begin{fulllineitems}
\phantomsection\label{\detokenize{api/hardware.devices:hardware.devices.esp32_thermometer.ThermometerData.timestamp}}
\pysigstartsignatures
\pysigline
{\sphinxbfcode{\sphinxupquote{timestamp}}\sphinxbfcode{\sphinxupquote{\DUrole{p}{:}\DUrole{w}{ }float}}\sphinxbfcode{\sphinxupquote{\DUrole{w}{ }\DUrole{p}{=}\DUrole{w}{ }0.0}}}
\pysigstopsignatures
\end{fulllineitems}

\index{from\_raw() (méthode de la classe hardware.devices.esp32\_thermometer.ThermometerData)@\spxentry{from\_raw()}\spxextra{méthode de la classe hardware.devices.esp32\_thermometer.ThermometerData}}

\begin{fulllineitems}
\phantomsection\label{\detokenize{api/hardware.devices:hardware.devices.esp32_thermometer.ThermometerData.from_raw}}
\pysigstartsignatures
\pysiglinewithargsret
{\sphinxbfcode{\sphinxupquote{\DUrole{k}{classmethod}\DUrole{w}{ }}}\sphinxbfcode{\sphinxupquote{from\_raw}}}
{\sphinxparam{\DUrole{n}{raw\_str}\DUrole{p}{:}\DUrole{w}{ }\DUrole{n}{str}}}
{{ $\rightarrow$ {\hyperref[\detokenize{api/hardware.devices:hardware.devices.esp32_thermometer.ThermometerData}]{\sphinxcrossref{ThermometerData}}}}}
\pysigstopsignatures
\sphinxAtStartPar
Décode la chaîne brute “22.5”.

\end{fulllineitems}


\end{fulllineitems}

\index{ThermometerDriver (classe dans hardware.devices.esp32\_thermometer)@\spxentry{ThermometerDriver}\spxextra{classe dans hardware.devices.esp32\_thermometer}}

\begin{fulllineitems}
\phantomsection\label{\detokenize{api/hardware.devices:hardware.devices.esp32_thermometer.ThermometerDriver}}
\pysigstartsignatures
\pysiglinewithargsret
{\sphinxbfcode{\sphinxupquote{\DUrole{k}{class}\DUrole{w}{ }}}\sphinxcode{\sphinxupquote{hardware.devices.esp32\_thermometer.}}\sphinxbfcode{\sphinxupquote{ThermometerDriver}}}
{\sphinxparam{\DUrole{keyword-only-separator}{\DUrole{o}{\sphinxstyleabbreviation{*} (Séparateur de paramètres par mot\sphinxhyphen{}clé uniquement (PEP 3102))}}}\sphinxparamcomma \sphinxparam{\DUrole{n}{target\_url}\DUrole{p}{:}\DUrole{w}{ }\DUrole{n}{str}\DUrole{w}{ }\DUrole{o}{=}\DUrole{w}{ }\DUrole{default_value}{\textquotesingle{}http://thermo.local/api\textquotesingle{}}}\sphinxparamcomma \sphinxparam{\DUrole{n}{on\_status}\DUrole{p}{:}\DUrole{w}{ }\DUrole{n}{Callable\DUrole{p}{{[}}\DUrole{p}{{[}}str\DUrole{p}{{]}}\DUrole{p}{,}\DUrole{w}{ }None\DUrole{p}{{]}}\DUrole{w}{ }\DUrole{p}{|}\DUrole{w}{ }None}\DUrole{w}{ }\DUrole{o}{=}\DUrole{w}{ }\DUrole{default_value}{None}}\sphinxparamcomma \sphinxparam{\DUrole{n}{on\_new\_measure}\DUrole{p}{:}\DUrole{w}{ }\DUrole{n}{Callable\DUrole{p}{{[}}\DUrole{p}{{[}}{\hyperref[\detokenize{api/hardware.devices:hardware.devices.esp32_thermometer.ThermometerData}]{\sphinxcrossref{ThermometerData}}}\DUrole{p}{{]}}\DUrole{p}{,}\DUrole{w}{ }None\DUrole{p}{{]}}\DUrole{w}{ }\DUrole{p}{|}\DUrole{w}{ }None}\DUrole{w}{ }\DUrole{o}{=}\DUrole{w}{ }\DUrole{default_value}{None}}}
{}
\pysigstopsignatures
\sphinxAtStartPar
Bases : \sphinxcode{\sphinxupquote{object}}

\sphinxAtStartPar
Driver haut niveau pour le thermomètre WiFi.

\sphinxAtStartPar
Responsabilités :
\sphinxhyphen{} Piloter le transport WiFi.
\sphinxhyphen{} Décoder les données brutes.
\sphinxhyphen{} Fournir des messages de statut en français pour le Speaker.
\index{start() (méthode hardware.devices.esp32\_thermometer.ThermometerDriver)@\spxentry{start()}\spxextra{méthode hardware.devices.esp32\_thermometer.ThermometerDriver}}

\begin{fulllineitems}
\phantomsection\label{\detokenize{api/hardware.devices:hardware.devices.esp32_thermometer.ThermometerDriver.start}}
\pysigstartsignatures
\pysiglinewithargsret
{\sphinxbfcode{\sphinxupquote{start}}}
{}
{{ $\rightarrow$ None}}
\pysigstopsignatures
\sphinxAtStartPar
Démarre la surveillance.

\end{fulllineitems}

\index{stop() (méthode hardware.devices.esp32\_thermometer.ThermometerDriver)@\spxentry{stop()}\spxextra{méthode hardware.devices.esp32\_thermometer.ThermometerDriver}}

\begin{fulllineitems}
\phantomsection\label{\detokenize{api/hardware.devices:hardware.devices.esp32_thermometer.ThermometerDriver.stop}}
\pysigstartsignatures
\pysiglinewithargsret
{\sphinxbfcode{\sphinxupquote{stop}}}
{}
{{ $\rightarrow$ None}}
\pysigstopsignatures
\sphinxAtStartPar
Arrête la surveillance.

\end{fulllineitems}

\index{is\_connected (propriété hardware.devices.esp32\_thermometer.ThermometerDriver)@\spxentry{is\_connected}\spxextra{propriété hardware.devices.esp32\_thermometer.ThermometerDriver}}

\begin{fulllineitems}
\phantomsection\label{\detokenize{api/hardware.devices:hardware.devices.esp32_thermometer.ThermometerDriver.is_connected}}
\pysigstartsignatures
\pysigline
{\sphinxbfcode{\sphinxupquote{\DUrole{k}{property}\DUrole{w}{ }}}\sphinxbfcode{\sphinxupquote{is\_connected}}\sphinxbfcode{\sphinxupquote{\DUrole{p}{:}\DUrole{w}{ }bool}}}
\pysigstopsignatures
\sphinxAtStartPar
Indique si le capteur est joignable.

\end{fulllineitems}

\index{get\_last\_measure() (méthode hardware.devices.esp32\_thermometer.ThermometerDriver)@\spxentry{get\_last\_measure()}\spxextra{méthode hardware.devices.esp32\_thermometer.ThermometerDriver}}

\begin{fulllineitems}
\phantomsection\label{\detokenize{api/hardware.devices:hardware.devices.esp32_thermometer.ThermometerDriver.get_last_measure}}
\pysigstartsignatures
\pysiglinewithargsret
{\sphinxbfcode{\sphinxupquote{get\_last\_measure}}}
{}
{{ $\rightarrow$ Dict\DUrole{p}{{[}}str\DUrole{p}{,}\DUrole{w}{ }float\DUrole{w}{ }\DUrole{p}{|}\DUrole{w}{ }None\DUrole{p}{{]}}}}
\pysigstopsignatures
\sphinxAtStartPar
Retourne la dernière mesure sous forme de dict simple.
Format compatible pour une consommation facile par le Mode.

\end{fulllineitems}


\end{fulllineitems}



\paragraph{hardware.devices.owon\_multimeter module}
\label{\detokenize{api/hardware.devices:module-hardware.devices.owon_multimeter}}\label{\detokenize{api/hardware.devices:hardware-devices-owon-multimeter-module}}\index{module@\spxentry{module}!hardware.devices.owon\_multimeter@\spxentry{hardware.devices.owon\_multimeter}}\index{hardware.devices.owon\_multimeter@\spxentry{hardware.devices.owon\_multimeter}!module@\spxentry{module}}
\sphinxAtStartPar
Driver pour le multimètre OWON 16.


\subparagraph{Vue d’ensemble}
\label{\detokenize{api/hardware.devices:id2}}
\sphinxAtStartPar
Ce module fournit deux briques :
\begin{enumerate}
\sphinxsetlistlabels{\arabic}{enumi}{enumii}{}{.}%
\item {} \begin{itemize}
\item {} 
\sphinxAtStartPar
Représente \sphinxstylestrong{une trame de 6 octets} envoyée par le multimètre OWON 16.

\item {} 
\sphinxAtStartPar
Décode la structure binaire (bits décimale / overflow / fonction / flags / valeur / signe).

\item {} \begin{description}
\sphinxlineitem{Présente les informations dans un objet Python lisible :}\begin{itemize}
\item {} 
\sphinxAtStartPar
\sphinxtitleref{value}           : valeur numérique en float (avec signe et décimale appliqués),

\item {} 
\sphinxAtStartPar
\sphinxtitleref{unit\_name}       : nom de la fonction, ex. \sphinxtitleref{« MILLI VOLT DC »},

\item {} 
\sphinxAtStartPar
flags booléens    : \sphinxtitleref{overflow}, \sphinxtitleref{auto\_ranging}, \sphinxtitleref{data\_hold\_mode}, \sphinxtitleref{relative\_mode}, \sphinxtitleref{low\_battery},

\item {} 
\sphinxAtStartPar
\sphinxtitleref{raw\_data}        : la trame brute (6 octets) pour debug / log.

\end{itemize}

\end{description}

\end{itemize}

\item {} \begin{itemize}
\item {} 
\sphinxAtStartPar
Driver \sphinxstylestrong{haut niveau} pour le multimètre OWON basé sur le transport générique \sphinxtitleref{BleClient}.

\item {} \begin{description}
\sphinxlineitem{Responsabilités :}\begin{itemize}
\item {} 
\sphinxAtStartPar
gérer la recherche / connexion BLE (\sphinxtitleref{start()}, \sphinxtitleref{stop()}, \sphinxtitleref{is\_connected}),

\item {} 
\sphinxAtStartPar
décoder les trames brutes en \sphinxtitleref{OwonMultimeterData},

\item {} 
\sphinxAtStartPar
mémoriser la \sphinxstylestrong{dernière mesure} (\sphinxtitleref{get\_last\_measure()}),

\item {} 
\sphinxAtStartPar
annoncer intelligemment les \sphinxstylestrong{changements de mode de mesure} (fonction du multimètre)
via un callback \sphinxtitleref{on\_status(text)} (par ex. le \sphinxtitleref{Speaker}).

\end{itemize}

\end{description}

\item {} 
\sphinxAtStartPar
Il ne fait \sphinxstylestrong{aucune synthèse vocale directement} : il envoie du texte vers le Mode,
qui lui\sphinxhyphen{}même appelle le \sphinxtitleref{Speaker}.

\end{itemize}

\end{enumerate}


\subparagraph{Couche \& responsabilités}
\label{\detokenize{api/hardware.devices:couche-responsabilites}}\begin{itemize}
\item {} \begin{description}
\sphinxlineitem{Ce module connaît :}\begin{itemize}
\item {} 
\sphinxAtStartPar
le protocole binaire OWON (layout des 6 octets),

\item {} 
\sphinxAtStartPar
les codes de fonction (\sphinxtitleref{OWON\_FUNCTION}),

\item {} 
\sphinxAtStartPar
la logique d’annonce des changements de mode (hystérésis).

\end{itemize}

\end{description}

\item {} \begin{description}
\sphinxlineitem{Il ne connaît PAS :}\begin{itemize}
\item {} 
\sphinxAtStartPar
le sélecteur rotatif (géré ailleurs),

\item {} 
\sphinxAtStartPar
l’architecture globale des modes,

\item {} 
\sphinxAtStartPar
la façon dont le texte sera lu (voix, langue, etc.).

\end{itemize}

\end{description}

\end{itemize}
\index{OWON\_FUNCTION (dans le module hardware.devices.owon\_multimeter)@\spxentry{OWON\_FUNCTION}\spxextra{dans le module hardware.devices.owon\_multimeter}}

\begin{fulllineitems}
\phantomsection\label{\detokenize{api/hardware.devices:hardware.devices.owon_multimeter.OWON_FUNCTION}}
\pysigstartsignatures
\pysigline
{\sphinxcode{\sphinxupquote{hardware.devices.owon\_multimeter.}}\sphinxbfcode{\sphinxupquote{OWON\_FUNCTION}}\sphinxbfcode{\sphinxupquote{\DUrole{p}{:}\DUrole{w}{ }Dict\DUrole{p}{{[}}str\DUrole{p}{,}\DUrole{w}{ }int\DUrole{p}{{]}}}}\sphinxbfcode{\sphinxupquote{\DUrole{w}{ }\DUrole{p}{=}\DUrole{w}{ }\{\textquotesingle{}AMPERE\_AC\textquotesingle{}: 7708, \textquotesingle{}AMPERE\_DC\textquotesingle{}: 7700, \textquotesingle{}CELSIUS\textquotesingle{}: 7748, \textquotesingle{}CONTINUITY\_TEST\textquotesingle{}: 7772, \textquotesingle{}DIODE\_TEST\textquotesingle{}: 7764, \textquotesingle{}FAHRENHEIT\textquotesingle{}: 7756, \textquotesingle{}FARAD\textquotesingle{}: 7724, \textquotesingle{}HERTZ\textquotesingle{}: 7732, \textquotesingle{}KILO\_OHM\textquotesingle{}: 7717, \textquotesingle{}MEGA\_OHM\textquotesingle{}: 7718, \textquotesingle{}MICRO\_AMPERE\_AC\textquotesingle{}: 7706, \textquotesingle{}MICRO\_AMPERE\_DC\textquotesingle{}: 7698, \textquotesingle{}MICRO\_FARAD\textquotesingle{}: 7722, \textquotesingle{}MILLI\_AMPERE\_AC\textquotesingle{}: 7707, \textquotesingle{}MILLI\_AMPERE\_DC\textquotesingle{}: 7699, \textquotesingle{}MILLI\_FARAD\textquotesingle{}: 7723, \textquotesingle{}MILLI\_VOLT\_AC\textquotesingle{}: 7691, \textquotesingle{}MILLI\_VOLT\_DC\textquotesingle{}: 7683, \textquotesingle{}NANO\_FARAD\textquotesingle{}: 7721, \textquotesingle{}NEAR\_FIELD\textquotesingle{}: 7788, \textquotesingle{}OHM\_NORMAL\textquotesingle{}: 7716, \textquotesingle{}PERCENTAGE\textquotesingle{}: 7740, \textquotesingle{}VOLT\_AC\textquotesingle{}: 7692, \textquotesingle{}VOLT\_DC\textquotesingle{}: 7684\}}}}
\pysigstopsignatures
\sphinxAtStartPar
Dictionnaire de mapping « nom symbolique » \sphinxhyphen{}\textgreater{} code 13 bits « function\_selector »

\sphinxAtStartPar
Ces valeurs proviennent de la documentation / reverse\sphinxhyphen{}engineering du
multimètre OWON 16. Le champ « fonction » est encodé sur 13 bits dans les
deux premiers octets de chaque trame.


\subparagraph{Remarque importante sur les milliampères}
\label{\detokenize{api/hardware.devices:remarque-importante-sur-les-milliamperes}}
\sphinxAtStartPar
Dans la doc d’origine, les codes \sphinxstylestrong{AC/DC pour les milliampères} sont
inversés par rapport à ce que le multimètre envoie réellement.
Après tests sur le matériel, nous avons corrigé le mapping :
\begin{itemize}
\item {} 
\sphinxAtStartPar
7699  \sphinxhyphen{}\textgreater{} « MILLI\_AMPERE\_DC »

\item {} 
\sphinxAtStartPar
7707  \sphinxhyphen{}\textgreater{} « MILLI\_AMPERE\_AC »

\end{itemize}

\sphinxAtStartPar
Cela garantit que les annonces vocales et \sphinxtitleref{unit\_name} sont cohérents avec
ce qui est affiché sur l’écran du multimètre (évite de dire « AC » quand
l’écran est en « DC » et inversement).

\end{fulllineitems}

\index{OWON\_MODE\_LABEL\_FR (dans le module hardware.devices.owon\_multimeter)@\spxentry{OWON\_MODE\_LABEL\_FR}\spxextra{dans le module hardware.devices.owon\_multimeter}}

\begin{fulllineitems}
\phantomsection\label{\detokenize{api/hardware.devices:hardware.devices.owon_multimeter.OWON_MODE_LABEL_FR}}
\pysigstartsignatures
\pysigline
{\sphinxcode{\sphinxupquote{hardware.devices.owon\_multimeter.}}\sphinxbfcode{\sphinxupquote{OWON\_MODE\_LABEL\_FR}}\sphinxbfcode{\sphinxupquote{\DUrole{p}{:}\DUrole{w}{ }Dict\DUrole{p}{{[}}int\DUrole{p}{,}\DUrole{w}{ }str\DUrole{p}{{]}}}}\sphinxbfcode{\sphinxupquote{\DUrole{w}{ }\DUrole{p}{=}\DUrole{w}{ }\{7683: \textquotesingle{}millivoltmètre courant continu\textquotesingle{}, 7684: \textquotesingle{}voltmètre courant continu\textquotesingle{}, 7691: \textquotesingle{}millivoltmètre courant alternatif\textquotesingle{}, 7692: \textquotesingle{}voltmètre courant alternatif\textquotesingle{}, 7698: \textquotesingle{}microampèremètre courant continu\textquotesingle{}, 7699: \textquotesingle{}milliampèremètre courant continu\textquotesingle{}, 7700: \textquotesingle{}ampèremètre courant continu\textquotesingle{}, 7706: \textquotesingle{}microampèremètre courant alternatif\textquotesingle{}, 7707: \textquotesingle{}milliampèremètre courant alternatif\textquotesingle{}, 7708: \textquotesingle{}ampèremètre courant alternatif\textquotesingle{}, 7716: \textquotesingle{}ohmmètre\textquotesingle{}, 7717: \textquotesingle{}kilo ohmmètre\textquotesingle{}, 7718: \textquotesingle{}méga ohmmètre\textquotesingle{}, 7721: \textquotesingle{}mesure de capacité en nanofarad\textquotesingle{}, 7722: \textquotesingle{}mesure de capacité en microfarad\textquotesingle{}, 7723: \textquotesingle{}mesure de capacité en millifarad\textquotesingle{}, 7724: \textquotesingle{}mesure de capacité en farad\textquotesingle{}, 7732: \textquotesingle{}fréquencemètre\textquotesingle{}, 7740: \textquotesingle{}mesure de pourcentage\textquotesingle{}, 7748: \textquotesingle{}thermomètre en degrés Celsius\textquotesingle{}, 7756: \textquotesingle{}thermomètre en degrés Fahrenheit\textquotesingle{}, 7764: \textquotesingle{}test de diode\textquotesingle{}, 7772: \textquotesingle{}test de continuité\textquotesingle{}, 7788: \textquotesingle{}détection de champ électrique\textquotesingle{}\}}}}
\pysigstopsignatures
\sphinxAtStartPar
Mapping « code fonction OWON » \sphinxhyphen{}\textgreater{} phrase FR annoncée

\sphinxAtStartPar
Utilisé par \sphinxtitleref{OwonMultimeterDriver} pour annoncer les changements de mode
(ex : « Mode voltmètre courant continu. »). Ces libellés sont pensés pour
être prononcés directement par la synthèse vocale.

\end{fulllineitems}

\index{OwonMultimeterData (classe dans hardware.devices.owon\_multimeter)@\spxentry{OwonMultimeterData}\spxextra{classe dans hardware.devices.owon\_multimeter}}

\begin{fulllineitems}
\phantomsection\label{\detokenize{api/hardware.devices:hardware.devices.owon_multimeter.OwonMultimeterData}}
\pysigstartsignatures
\pysiglinewithargsret
{\sphinxbfcode{\sphinxupquote{\DUrole{k}{class}\DUrole{w}{ }}}\sphinxcode{\sphinxupquote{hardware.devices.owon\_multimeter.}}\sphinxbfcode{\sphinxupquote{OwonMultimeterData}}}
{\sphinxparam{\DUrole{n}{raw\_data}\DUrole{p}{:}\DUrole{w}{ }\DUrole{n}{List\DUrole{p}{{[}}int\DUrole{p}{{]}}}}\sphinxparamcomma \sphinxparam{\DUrole{n}{decimal\_places}\DUrole{p}{:}\DUrole{w}{ }\DUrole{n}{int}}\sphinxparamcomma \sphinxparam{\DUrole{n}{overflow}\DUrole{p}{:}\DUrole{w}{ }\DUrole{n}{int}}\sphinxparamcomma \sphinxparam{\DUrole{n}{unit}\DUrole{p}{:}\DUrole{w}{ }\DUrole{n}{int}}\sphinxparamcomma \sphinxparam{\DUrole{n}{unit\_name}\DUrole{p}{:}\DUrole{w}{ }\DUrole{n}{str}}\sphinxparamcomma \sphinxparam{\DUrole{n}{data\_hold\_mode}\DUrole{p}{:}\DUrole{w}{ }\DUrole{n}{int}}\sphinxparamcomma \sphinxparam{\DUrole{n}{relative\_mode}\DUrole{p}{:}\DUrole{w}{ }\DUrole{n}{int}}\sphinxparamcomma \sphinxparam{\DUrole{n}{auto\_ranging}\DUrole{p}{:}\DUrole{w}{ }\DUrole{n}{int}}\sphinxparamcomma \sphinxparam{\DUrole{n}{low\_battery}\DUrole{p}{:}\DUrole{w}{ }\DUrole{n}{int}}\sphinxparamcomma \sphinxparam{\DUrole{n}{value}\DUrole{p}{:}\DUrole{w}{ }\DUrole{n}{float}}\sphinxparamcomma \sphinxparam{\DUrole{n}{sign}\DUrole{p}{:}\DUrole{w}{ }\DUrole{n}{int}}}
{}
\pysigstopsignatures
\sphinxAtStartPar
Bases : \sphinxcode{\sphinxupquote{object}}

\sphinxAtStartPar
Représente une trame décodée du multimètre OWON 16 (6 octets).


\subparagraph{Structure binaire (vue « fabricant »)}
\label{\detokenize{api/hardware.devices:structure-binaire-vue-fabricant}}
\sphinxAtStartPar
Les 6 octets sont organisés comme suit (Little Endian) :
\begin{quote}

\sphinxAtStartPar
decimal\_places          0           2
overflow                2           1
function\_selector       3          13
data\_hold\_mode         16           1
relative\_mode          17           1
auto\_ranging           18           1
low\_battery            19           1
not\_used\_1             20           4
not\_used\_2             24           8
value                  32          15
sign                   47           1
\end{quote}


\subparagraph{Interprétation}
\label{\detokenize{api/hardware.devices:interpretation}}\begin{itemize}
\item {} 
\sphinxAtStartPar
\sphinxtitleref{decimal\_places} : position de la décimale (0, 1, 2 ou 3).

\item {} 
\sphinxAtStartPar
\sphinxtitleref{overflow}       : dépassement de gamme (0 = OK, 1 = overflow).

\item {} \begin{description}
\sphinxlineitem{\sphinxtitleref{function\_selector} (ici \sphinxtitleref{unit}) :}
\sphinxAtStartPar
code 13 bits mappé via \sphinxtitleref{OWON\_FUNCTION} vers un nom symbolique.

\end{description}

\item {} 
\sphinxAtStartPar
\sphinxtitleref{data\_hold\_mode} : flag HOLD (gel d’affichage).

\item {} 
\sphinxAtStartPar
\sphinxtitleref{relative\_mode}  : mode relatif (REL).

\item {} 
\sphinxAtStartPar
\sphinxtitleref{auto\_ranging}   : autorange actif ou non.

\item {} 
\sphinxAtStartPar
\sphinxtitleref{low\_battery}    : batterie faible.

\item {} 
\sphinxAtStartPar
\sphinxtitleref{value}          : valeur brute signée (15 bits + signe), mise à l’échelle.

\item {} 
\sphinxAtStartPar
\sphinxtitleref{sign}           : 0 = positif, 1 = négatif.

\end{itemize}


\subparagraph{Attributs principaux exposés}
\label{\detokenize{api/hardware.devices:attributs-principaux-exposes}}\begin{description}
\sphinxlineitem{raw\_data}{[}List{[}int{]}{]}
\sphinxAtStartPar
Trame brute reçue (6 octets).

\sphinxlineitem{decimal\_places}{[}int{]}
\sphinxAtStartPar
Nombre de décimales à appliquer à la valeur.

\sphinxlineitem{overflow}{[}int{]}
\sphinxAtStartPar
Flag overflow (0 ou 1).

\sphinxlineitem{unit}{[}int{]}
\sphinxAtStartPar
Code numérique de la fonction (13 bits).

\sphinxlineitem{unit\_name}{[}str{]}
\sphinxAtStartPar
Nom lisible de la fonction (ex: « MILLI VOLT DC », « OHM NORMAL », …).

\sphinxlineitem{data\_hold\_mode}{[}int{]}
\sphinxAtStartPar
Flag HOLD (0 ou 1).

\sphinxlineitem{relative\_mode}{[}int{]}
\sphinxAtStartPar
Flag REL (0 ou 1).

\sphinxlineitem{auto\_ranging}{[}int{]}
\sphinxAtStartPar
Flag AutoRange (0 ou 1).

\sphinxlineitem{low\_battery}{[}int{]}
\sphinxAtStartPar
Flag batterie faible (0 ou 1).

\sphinxlineitem{value}{[}float{]}
\sphinxAtStartPar
Valeur numérique finale (signe et décimale appliqués).

\sphinxlineitem{sign}{[}int{]}
\sphinxAtStartPar
Flag de signe (0 = positif, 1 = négatif).

\end{description}
\index{raw\_data (attribut hardware.devices.owon\_multimeter.OwonMultimeterData)@\spxentry{raw\_data}\spxextra{attribut hardware.devices.owon\_multimeter.OwonMultimeterData}}

\begin{fulllineitems}
\phantomsection\label{\detokenize{api/hardware.devices:hardware.devices.owon_multimeter.OwonMultimeterData.raw_data}}
\pysigstartsignatures
\pysigline
{\sphinxbfcode{\sphinxupquote{raw\_data}}\sphinxbfcode{\sphinxupquote{\DUrole{p}{:}\DUrole{w}{ }List\DUrole{p}{{[}}int\DUrole{p}{{]}}}}}
\pysigstopsignatures
\end{fulllineitems}

\index{decimal\_places (attribut hardware.devices.owon\_multimeter.OwonMultimeterData)@\spxentry{decimal\_places}\spxextra{attribut hardware.devices.owon\_multimeter.OwonMultimeterData}}

\begin{fulllineitems}
\phantomsection\label{\detokenize{api/hardware.devices:hardware.devices.owon_multimeter.OwonMultimeterData.decimal_places}}
\pysigstartsignatures
\pysigline
{\sphinxbfcode{\sphinxupquote{decimal\_places}}\sphinxbfcode{\sphinxupquote{\DUrole{p}{:}\DUrole{w}{ }int}}}
\pysigstopsignatures
\end{fulllineitems}

\index{overflow (attribut hardware.devices.owon\_multimeter.OwonMultimeterData)@\spxentry{overflow}\spxextra{attribut hardware.devices.owon\_multimeter.OwonMultimeterData}}

\begin{fulllineitems}
\phantomsection\label{\detokenize{api/hardware.devices:hardware.devices.owon_multimeter.OwonMultimeterData.overflow}}
\pysigstartsignatures
\pysigline
{\sphinxbfcode{\sphinxupquote{overflow}}\sphinxbfcode{\sphinxupquote{\DUrole{p}{:}\DUrole{w}{ }int}}}
\pysigstopsignatures
\end{fulllineitems}

\index{unit (attribut hardware.devices.owon\_multimeter.OwonMultimeterData)@\spxentry{unit}\spxextra{attribut hardware.devices.owon\_multimeter.OwonMultimeterData}}

\begin{fulllineitems}
\phantomsection\label{\detokenize{api/hardware.devices:hardware.devices.owon_multimeter.OwonMultimeterData.unit}}
\pysigstartsignatures
\pysigline
{\sphinxbfcode{\sphinxupquote{unit}}\sphinxbfcode{\sphinxupquote{\DUrole{p}{:}\DUrole{w}{ }int}}}
\pysigstopsignatures
\end{fulllineitems}

\index{unit\_name (attribut hardware.devices.owon\_multimeter.OwonMultimeterData)@\spxentry{unit\_name}\spxextra{attribut hardware.devices.owon\_multimeter.OwonMultimeterData}}

\begin{fulllineitems}
\phantomsection\label{\detokenize{api/hardware.devices:hardware.devices.owon_multimeter.OwonMultimeterData.unit_name}}
\pysigstartsignatures
\pysigline
{\sphinxbfcode{\sphinxupquote{unit\_name}}\sphinxbfcode{\sphinxupquote{\DUrole{p}{:}\DUrole{w}{ }str}}}
\pysigstopsignatures
\end{fulllineitems}

\index{data\_hold\_mode (attribut hardware.devices.owon\_multimeter.OwonMultimeterData)@\spxentry{data\_hold\_mode}\spxextra{attribut hardware.devices.owon\_multimeter.OwonMultimeterData}}

\begin{fulllineitems}
\phantomsection\label{\detokenize{api/hardware.devices:hardware.devices.owon_multimeter.OwonMultimeterData.data_hold_mode}}
\pysigstartsignatures
\pysigline
{\sphinxbfcode{\sphinxupquote{data\_hold\_mode}}\sphinxbfcode{\sphinxupquote{\DUrole{p}{:}\DUrole{w}{ }int}}}
\pysigstopsignatures
\end{fulllineitems}

\index{relative\_mode (attribut hardware.devices.owon\_multimeter.OwonMultimeterData)@\spxentry{relative\_mode}\spxextra{attribut hardware.devices.owon\_multimeter.OwonMultimeterData}}

\begin{fulllineitems}
\phantomsection\label{\detokenize{api/hardware.devices:hardware.devices.owon_multimeter.OwonMultimeterData.relative_mode}}
\pysigstartsignatures
\pysigline
{\sphinxbfcode{\sphinxupquote{relative\_mode}}\sphinxbfcode{\sphinxupquote{\DUrole{p}{:}\DUrole{w}{ }int}}}
\pysigstopsignatures
\end{fulllineitems}

\index{auto\_ranging (attribut hardware.devices.owon\_multimeter.OwonMultimeterData)@\spxentry{auto\_ranging}\spxextra{attribut hardware.devices.owon\_multimeter.OwonMultimeterData}}

\begin{fulllineitems}
\phantomsection\label{\detokenize{api/hardware.devices:hardware.devices.owon_multimeter.OwonMultimeterData.auto_ranging}}
\pysigstartsignatures
\pysigline
{\sphinxbfcode{\sphinxupquote{auto\_ranging}}\sphinxbfcode{\sphinxupquote{\DUrole{p}{:}\DUrole{w}{ }int}}}
\pysigstopsignatures
\end{fulllineitems}

\index{low\_battery (attribut hardware.devices.owon\_multimeter.OwonMultimeterData)@\spxentry{low\_battery}\spxextra{attribut hardware.devices.owon\_multimeter.OwonMultimeterData}}

\begin{fulllineitems}
\phantomsection\label{\detokenize{api/hardware.devices:hardware.devices.owon_multimeter.OwonMultimeterData.low_battery}}
\pysigstartsignatures
\pysigline
{\sphinxbfcode{\sphinxupquote{low\_battery}}\sphinxbfcode{\sphinxupquote{\DUrole{p}{:}\DUrole{w}{ }int}}}
\pysigstopsignatures
\end{fulllineitems}

\index{value (attribut hardware.devices.owon\_multimeter.OwonMultimeterData)@\spxentry{value}\spxextra{attribut hardware.devices.owon\_multimeter.OwonMultimeterData}}

\begin{fulllineitems}
\phantomsection\label{\detokenize{api/hardware.devices:hardware.devices.owon_multimeter.OwonMultimeterData.value}}
\pysigstartsignatures
\pysigline
{\sphinxbfcode{\sphinxupquote{value}}\sphinxbfcode{\sphinxupquote{\DUrole{p}{:}\DUrole{w}{ }float}}}
\pysigstopsignatures
\end{fulllineitems}

\index{sign (attribut hardware.devices.owon\_multimeter.OwonMultimeterData)@\spxentry{sign}\spxextra{attribut hardware.devices.owon\_multimeter.OwonMultimeterData}}

\begin{fulllineitems}
\phantomsection\label{\detokenize{api/hardware.devices:hardware.devices.owon_multimeter.OwonMultimeterData.sign}}
\pysigstartsignatures
\pysigline
{\sphinxbfcode{\sphinxupquote{sign}}\sphinxbfcode{\sphinxupquote{\DUrole{p}{:}\DUrole{w}{ }int}}}
\pysigstopsignatures
\end{fulllineitems}

\index{from\_raw() (méthode de la classe hardware.devices.owon\_multimeter.OwonMultimeterData)@\spxentry{from\_raw()}\spxextra{méthode de la classe hardware.devices.owon\_multimeter.OwonMultimeterData}}

\begin{fulllineitems}
\phantomsection\label{\detokenize{api/hardware.devices:hardware.devices.owon_multimeter.OwonMultimeterData.from_raw}}
\pysigstartsignatures
\pysiglinewithargsret
{\sphinxbfcode{\sphinxupquote{\DUrole{k}{classmethod}\DUrole{w}{ }}}\sphinxbfcode{\sphinxupquote{from\_raw}}}
{\sphinxparam{\DUrole{n}{raw\_data}\DUrole{p}{:}\DUrole{w}{ }\DUrole{n}{List\DUrole{p}{{[}}int\DUrole{p}{{]}}}}}
{{ $\rightarrow$ {\hyperref[\detokenize{api/hardware.devices:hardware.devices.owon_multimeter.OwonMultimeterData}]{\sphinxcrossref{OwonMultimeterData}}}}}
\pysigstopsignatures
\sphinxAtStartPar
Construit une instance à partir d’une trame brute (6 octets).


\subparagraph{Paramètres}
\label{\detokenize{api/hardware.devices:id3}}\begin{description}
\sphinxlineitem{raw\_data}{[}List{[}int{]}{]}
\sphinxAtStartPar
Liste de 6 octets (0\textendash{}255) reçus depuis le multimètre.

\end{description}
\begin{quote}\begin{description}
\sphinxlineitem{returns}
\sphinxAtStartPar
Instance décodée, prête à être utilisée (value, flags, etc.).

\sphinxlineitem{rtype}
\sphinxAtStartPar
OwonMultimeterData

\sphinxlineitem{raises ValueError}
\sphinxAtStartPar
Si la longueur de \sphinxtitleref{raw\_data} n’est pas égale à 6.

\end{description}\end{quote}

\end{fulllineitems}


\end{fulllineitems}

\index{OwonMultimeterDriver (classe dans hardware.devices.owon\_multimeter)@\spxentry{OwonMultimeterDriver}\spxextra{classe dans hardware.devices.owon\_multimeter}}

\begin{fulllineitems}
\phantomsection\label{\detokenize{api/hardware.devices:hardware.devices.owon_multimeter.OwonMultimeterDriver}}
\pysigstartsignatures
\pysiglinewithargsret
{\sphinxbfcode{\sphinxupquote{\DUrole{k}{class}\DUrole{w}{ }}}\sphinxcode{\sphinxupquote{hardware.devices.owon\_multimeter.}}\sphinxbfcode{\sphinxupquote{OwonMultimeterDriver}}}
{\sphinxparam{\DUrole{keyword-only-separator}{\DUrole{o}{\sphinxstyleabbreviation{*}}}}\sphinxparamcomma \sphinxparam{\DUrole{n}{on\_status}\DUrole{p}{:}\DUrole{w}{ }\DUrole{n}{Callable\DUrole{p}{{[}}\DUrole{p}{{[}}str\DUrole{p}{{]}}\DUrole{p}{,}\DUrole{w}{ }None\DUrole{p}{{]}}\DUrole{w}{ }\DUrole{p}{|}\DUrole{w}{ }None}\DUrole{w}{ }\DUrole{o}{=}\DUrole{w}{ }\DUrole{default_value}{None}}\sphinxparamcomma \sphinxparam{\DUrole{n}{on\_new\_measure}\DUrole{p}{:}\DUrole{w}{ }\DUrole{n}{Callable\DUrole{p}{{[}}\DUrole{p}{{[}}{\hyperref[\detokenize{api/hardware.devices:hardware.devices.owon_multimeter.OwonMultimeterData}]{\sphinxcrossref{OwonMultimeterData}}}\DUrole{p}{{]}}\DUrole{p}{,}\DUrole{w}{ }None\DUrole{p}{{]}}\DUrole{w}{ }\DUrole{p}{|}\DUrole{w}{ }None}\DUrole{w}{ }\DUrole{o}{=}\DUrole{w}{ }\DUrole{default_value}{None}}\sphinxparamcomma \sphinxparam{\DUrole{n}{name\_filter}\DUrole{p}{:}\DUrole{w}{ }\DUrole{n}{str}\DUrole{w}{ }\DUrole{o}{=}\DUrole{w}{ }\DUrole{default_value}{\textquotesingle{}BDM\textquotesingle{}}}\sphinxparamcomma \sphinxparam{\DUrole{n}{address\_filter}\DUrole{p}{:}\DUrole{w}{ }\DUrole{n}{str\DUrole{w}{ }\DUrole{p}{|}\DUrole{w}{ }None}\DUrole{w}{ }\DUrole{o}{=}\DUrole{w}{ }\DUrole{default_value}{None}}\sphinxparamcomma \sphinxparam{\DUrole{n}{characteristic\_uuid}\DUrole{p}{:}\DUrole{w}{ }\DUrole{n}{str}\DUrole{w}{ }\DUrole{o}{=}\DUrole{w}{ }\DUrole{default_value}{\textquotesingle{}0000fff4\sphinxhyphen{}0000\sphinxhyphen{}1000\sphinxhyphen{}8000\sphinxhyphen{}00805f9b34fb\textquotesingle{}}}\sphinxparamcomma \sphinxparam{\DUrole{n}{scan\_timeout}\DUrole{p}{:}\DUrole{w}{ }\DUrole{n}{float}\DUrole{w}{ }\DUrole{o}{=}\DUrole{w}{ }\DUrole{default_value}{4.0}}\sphinxparamcomma \sphinxparam{\DUrole{n}{retry\_delay}\DUrole{p}{:}\DUrole{w}{ }\DUrole{n}{float}\DUrole{w}{ }\DUrole{o}{=}\DUrole{w}{ }\DUrole{default_value}{15.0}}\sphinxparamcomma \sphinxparam{\DUrole{n}{reconnect\_delay}\DUrole{p}{:}\DUrole{w}{ }\DUrole{n}{float}\DUrole{w}{ }\DUrole{o}{=}\DUrole{w}{ }\DUrole{default_value}{2.0}}\sphinxparamcomma \sphinxparam{\DUrole{n}{listen\_poll\_interval\_sec}\DUrole{p}{:}\DUrole{w}{ }\DUrole{n}{float}\DUrole{w}{ }\DUrole{o}{=}\DUrole{w}{ }\DUrole{default_value}{0.5}}\sphinxparamcomma \sphinxparam{\DUrole{n}{mode\_announce\_delay}\DUrole{p}{:}\DUrole{w}{ }\DUrole{n}{float}\DUrole{w}{ }\DUrole{o}{=}\DUrole{w}{ }\DUrole{default_value}{0.7}}}
{}
\pysigstopsignatures
\sphinxAtStartPar
Bases : \sphinxcode{\sphinxupquote{object}}

\sphinxAtStartPar
Driver haut niveau pour le multimètre OWON 16.


\subparagraph{Rôle}
\label{\detokenize{api/hardware.devices:role}}\begin{itemize}
\item {} 
\sphinxAtStartPar
Utiliser \sphinxtitleref{BleClient} comme \sphinxstylestrong{transport BLE générique}.

\item {} 
\sphinxAtStartPar
Décoder les trames brutes reçues via BLE en objets \sphinxtitleref{OwonMultimeterData}.

\item {} 
\sphinxAtStartPar
Mémoriser la dernière mesure décodée (\sphinxtitleref{\_last\_data}).

\item {} 
\sphinxAtStartPar
Annoncer les \sphinxstylestrong{changements de mode de mesure} (tension, courant, ohms, …)
via une logique d’hystérésis (temps minimum de stabilité).

\item {} \begin{description}
\sphinxlineitem{Exposer une API simple pour les Modes :}\begin{itemize}
\item {} 
\sphinxAtStartPar
\sphinxtitleref{start()} / \sphinxtitleref{stop()} pour lancer/arrêter la surveillance BLE.

\item {} 
\sphinxAtStartPar
\sphinxtitleref{is\_connected} pour vérifier l’état de connexion.

\item {} 
\sphinxAtStartPar
\sphinxtitleref{get\_last\_measure()} pour récupérer la dernière mesure sous forme de dict.

\end{itemize}

\end{description}

\end{itemize}


\subparagraph{Couche d’abstraction}
\label{\detokenize{api/hardware.devices:couche-dabstraction}}\begin{itemize}
\item {} 
\sphinxAtStartPar
Ce driver ne fait pas de synthèse vocale.
Il pousse simplement des \sphinxstylestrong{messages texte} dans le callback \sphinxtitleref{on\_status}.

\item {} 
\sphinxAtStartPar
Le Mode (ex: \sphinxtitleref{ModeMultimetre}) se charge ensuite de passer ces messages
au \sphinxtitleref{Speaker} central.

\end{itemize}
\index{\_\_init\_\_() (méthode hardware.devices.owon\_multimeter.OwonMultimeterDriver)@\spxentry{\_\_init\_\_()}\spxextra{méthode hardware.devices.owon\_multimeter.OwonMultimeterDriver}}

\begin{fulllineitems}
\phantomsection\label{\detokenize{api/hardware.devices:hardware.devices.owon_multimeter.OwonMultimeterDriver.__init__}}
\pysigstartsignatures
\pysiglinewithargsret
{\sphinxbfcode{\sphinxupquote{\_\_init\_\_}}}
{\sphinxparam{\DUrole{keyword-only-separator}{\DUrole{o}{\sphinxstyleabbreviation{*}}}}\sphinxparamcomma \sphinxparam{\DUrole{n}{on\_status}\DUrole{p}{:}\DUrole{w}{ }\DUrole{n}{Callable\DUrole{p}{{[}}\DUrole{p}{{[}}str\DUrole{p}{{]}}\DUrole{p}{,}\DUrole{w}{ }None\DUrole{p}{{]}}\DUrole{w}{ }\DUrole{p}{|}\DUrole{w}{ }None}\DUrole{w}{ }\DUrole{o}{=}\DUrole{w}{ }\DUrole{default_value}{None}}\sphinxparamcomma \sphinxparam{\DUrole{n}{on\_new\_measure}\DUrole{p}{:}\DUrole{w}{ }\DUrole{n}{Callable\DUrole{p}{{[}}\DUrole{p}{{[}}{\hyperref[\detokenize{api/hardware.devices:hardware.devices.owon_multimeter.OwonMultimeterData}]{\sphinxcrossref{OwonMultimeterData}}}\DUrole{p}{{]}}\DUrole{p}{,}\DUrole{w}{ }None\DUrole{p}{{]}}\DUrole{w}{ }\DUrole{p}{|}\DUrole{w}{ }None}\DUrole{w}{ }\DUrole{o}{=}\DUrole{w}{ }\DUrole{default_value}{None}}\sphinxparamcomma \sphinxparam{\DUrole{n}{name\_filter}\DUrole{p}{:}\DUrole{w}{ }\DUrole{n}{str}\DUrole{w}{ }\DUrole{o}{=}\DUrole{w}{ }\DUrole{default_value}{\textquotesingle{}BDM\textquotesingle{}}}\sphinxparamcomma \sphinxparam{\DUrole{n}{address\_filter}\DUrole{p}{:}\DUrole{w}{ }\DUrole{n}{str\DUrole{w}{ }\DUrole{p}{|}\DUrole{w}{ }None}\DUrole{w}{ }\DUrole{o}{=}\DUrole{w}{ }\DUrole{default_value}{None}}\sphinxparamcomma \sphinxparam{\DUrole{n}{characteristic\_uuid}\DUrole{p}{:}\DUrole{w}{ }\DUrole{n}{str}\DUrole{w}{ }\DUrole{o}{=}\DUrole{w}{ }\DUrole{default_value}{\textquotesingle{}0000fff4\sphinxhyphen{}0000\sphinxhyphen{}1000\sphinxhyphen{}8000\sphinxhyphen{}00805f9b34fb\textquotesingle{}}}\sphinxparamcomma \sphinxparam{\DUrole{n}{scan\_timeout}\DUrole{p}{:}\DUrole{w}{ }\DUrole{n}{float}\DUrole{w}{ }\DUrole{o}{=}\DUrole{w}{ }\DUrole{default_value}{4.0}}\sphinxparamcomma \sphinxparam{\DUrole{n}{retry\_delay}\DUrole{p}{:}\DUrole{w}{ }\DUrole{n}{float}\DUrole{w}{ }\DUrole{o}{=}\DUrole{w}{ }\DUrole{default_value}{15.0}}\sphinxparamcomma \sphinxparam{\DUrole{n}{reconnect\_delay}\DUrole{p}{:}\DUrole{w}{ }\DUrole{n}{float}\DUrole{w}{ }\DUrole{o}{=}\DUrole{w}{ }\DUrole{default_value}{2.0}}\sphinxparamcomma \sphinxparam{\DUrole{n}{listen\_poll\_interval\_sec}\DUrole{p}{:}\DUrole{w}{ }\DUrole{n}{float}\DUrole{w}{ }\DUrole{o}{=}\DUrole{w}{ }\DUrole{default_value}{0.5}}\sphinxparamcomma \sphinxparam{\DUrole{n}{mode\_announce\_delay}\DUrole{p}{:}\DUrole{w}{ }\DUrole{n}{float}\DUrole{w}{ }\DUrole{o}{=}\DUrole{w}{ }\DUrole{default_value}{0.7}}}
{{ $\rightarrow$ None}}
\pysigstopsignatures
\sphinxAtStartPar
Initialise le driver OWON.


\subparagraph{Paramètres}
\label{\detokenize{api/hardware.devices:id4}}\begin{description}
\sphinxlineitem{on\_status}{[}Optional{[}Callable{[}{[}str{]}, None{]}{]}{]}
\sphinxAtStartPar
Callback pour annoncer des messages utilisateur (via Speaker).
Exemple de messages :
\begin{itemize}
\item {} 
\sphinxAtStartPar
« Multimètre connecté. Vous pouvez effectuer vos mesures. »

\item {} 
\sphinxAtStartPar
« Mode voltmètre courant continu. »

\item {} 
\sphinxAtStartPar
« Multimètre non trouvé. Assurez\sphinxhyphen{}vous qu’il est allumé… »

\end{itemize}

\sphinxlineitem{on\_new\_measure}{[}Optional{[}Callable{[}{[}OwonMultimeterData{]}, None{]}{]}{]}
\sphinxAtStartPar
Callback appelé à chaque \sphinxstylestrong{nouvelle trame décodée}.
Permet au Mode d’enregistrer directement la dernière mesure brute.

\sphinxlineitem{retry\_delay}{[}float{]}
\sphinxAtStartPar
Délai (en secondes) entre deux tentatives lorsqu’aucun multimètre
n’est trouvé lors du scan BLE.

\end{description}

\end{fulllineitems}

\index{start() (méthode hardware.devices.owon\_multimeter.OwonMultimeterDriver)@\spxentry{start()}\spxextra{méthode hardware.devices.owon\_multimeter.OwonMultimeterDriver}}

\begin{fulllineitems}
\phantomsection\label{\detokenize{api/hardware.devices:hardware.devices.owon_multimeter.OwonMultimeterDriver.start}}
\pysigstartsignatures
\pysiglinewithargsret
{\sphinxbfcode{\sphinxupquote{start}}}
{}
{{ $\rightarrow$ None}}
\pysigstopsignatures
\sphinxAtStartPar
Démarre la recherche / connexion au multimètre OWON.


\subparagraph{Effets}
\label{\detokenize{api/hardware.devices:effets}}\begin{itemize}
\item {} 
\sphinxAtStartPar
Lance le \sphinxtitleref{BleClient} dans un thread séparé.

\item {} 
\sphinxAtStartPar
Le driver se met à l’écoute des trames et des événements d’état.

\end{itemize}

\end{fulllineitems}

\index{stop() (méthode hardware.devices.owon\_multimeter.OwonMultimeterDriver)@\spxentry{stop()}\spxextra{méthode hardware.devices.owon\_multimeter.OwonMultimeterDriver}}

\begin{fulllineitems}
\phantomsection\label{\detokenize{api/hardware.devices:hardware.devices.owon_multimeter.OwonMultimeterDriver.stop}}
\pysigstartsignatures
\pysiglinewithargsret
{\sphinxbfcode{\sphinxupquote{stop}}}
{}
{{ $\rightarrow$ None}}
\pysigstopsignatures
\sphinxAtStartPar
Arrête la surveillance du multimètre.


\subparagraph{Effets}
\label{\detokenize{api/hardware.devices:id5}}\begin{itemize}
\item {} 
\sphinxAtStartPar
Demande l’arrêt du \sphinxtitleref{BleClient} (stop de la boucle de monitoring).

\item {} 
\sphinxAtStartPar
À terme, la connexion BLE est fermée et l’état repasse à « stopped ».

\end{itemize}

\end{fulllineitems}

\index{is\_connected (propriété hardware.devices.owon\_multimeter.OwonMultimeterDriver)@\spxentry{is\_connected}\spxextra{propriété hardware.devices.owon\_multimeter.OwonMultimeterDriver}}

\begin{fulllineitems}
\phantomsection\label{\detokenize{api/hardware.devices:hardware.devices.owon_multimeter.OwonMultimeterDriver.is_connected}}
\pysigstartsignatures
\pysigline
{\sphinxbfcode{\sphinxupquote{\DUrole{k}{property}\DUrole{w}{ }}}\sphinxbfcode{\sphinxupquote{is\_connected}}\sphinxbfcode{\sphinxupquote{\DUrole{p}{:}\DUrole{w}{ }bool}}}
\pysigstopsignatures
\sphinxAtStartPar
Indique si le multimètre est actuellement connecté au niveau BLE.
\begin{quote}\begin{description}
\sphinxlineitem{Renvoie}
\sphinxAtStartPar
\sphinxtitleref{True} si \sphinxtitleref{BleClient} est connecté au périphérique OWON,
\sphinxtitleref{False} sinon.

\sphinxlineitem{Type renvoyé}
\sphinxAtStartPar
bool

\end{description}\end{quote}

\end{fulllineitems}

\index{get\_last\_measure() (méthode hardware.devices.owon\_multimeter.OwonMultimeterDriver)@\spxentry{get\_last\_measure()}\spxextra{méthode hardware.devices.owon\_multimeter.OwonMultimeterDriver}}

\begin{fulllineitems}
\phantomsection\label{\detokenize{api/hardware.devices:hardware.devices.owon_multimeter.OwonMultimeterDriver.get_last_measure}}
\pysigstartsignatures
\pysiglinewithargsret
{\sphinxbfcode{\sphinxupquote{get\_last\_measure}}}
{}
{{ $\rightarrow$ Dict\DUrole{p}{{[}}str\DUrole{p}{,}\DUrole{w}{ }object\DUrole{w}{ }\DUrole{p}{|}\DUrole{w}{ }None\DUrole{p}{{]}}}}
\pysigstopsignatures
\sphinxAtStartPar
Retourne la dernière mesure sous forme de dict simple, prêt à consommer.

\sphinxAtStartPar
Ce format est volontairement \sphinxstylestrong{générique} pour être facile à utiliser
côté Modes / IHM / logs, sans exposer directement l’objet \sphinxtitleref{OwonMultimeterData}.


\subparagraph{Clés du dictionnaire}
\label{\detokenize{api/hardware.devices:cles-du-dictionnaire}}\begin{itemize}
\item {} 
\sphinxAtStartPar
« value »         : float ou None

\item {} 
\sphinxAtStartPar
« unit\_name »     : str ou None (ex: « MILLI VOLT DC », « OHM NORMAL »)

\item {} 
\sphinxAtStartPar
« overflow »      : bool ou None

\item {} 
\sphinxAtStartPar
« auto\_ranging »  : bool ou None

\item {} 
\sphinxAtStartPar
« data\_hold\_mode »: bool ou None

\item {} 
\sphinxAtStartPar
« relative\_mode » : bool ou None

\end{itemize}
\begin{quote}\begin{description}
\sphinxlineitem{returns}
\sphinxAtStartPar
Dictionnaire avec les champs ci\sphinxhyphen{}dessus. Si aucune mesure n’a
encore été reçue, toutes les valeurs sont à None.

\sphinxlineitem{rtype}
\sphinxAtStartPar
dict

\end{description}\end{quote}

\end{fulllineitems}


\end{fulllineitems}



\paragraph{hardware.devices.text\_reader\_device module}
\label{\detokenize{api/hardware.devices:module-hardware.devices.text_reader_device}}\label{\detokenize{api/hardware.devices:hardware-devices-text-reader-device-module}}\index{module@\spxentry{module}!hardware.devices.text\_reader\_device@\spxentry{hardware.devices.text\_reader\_device}}\index{hardware.devices.text\_reader\_device@\spxentry{hardware.devices.text\_reader\_device}!module@\spxentry{module}}

\subparagraph{hardware.devices.text\_reader\_device}
\label{\detokenize{api/hardware.devices:hardware-devices-text-reader-device}}
\sphinxAtStartPar
Device « machine à lire » (couche au\sphinxhyphen{}dessus du hardware brut).


\subparagraph{Responsabilité}
\label{\detokenize{api/hardware.devices:responsabilite}}
\sphinxAtStartPar
Orchestrer les trois premières étapes de la chaîne de lecture :
\begin{enumerate}
\sphinxsetlistlabels{\arabic}{enumi}{enumii}{}{.}%
\item {} 
\sphinxAtStartPar
Capture caméra       \(\rightarrow\) \textless{}basename\textgreater{}.jpg          (via PiCamera)

\item {} 
\sphinxAtStartPar
OCR                  \(\rightarrow\) \textless{}basename\textgreater{}\_raw.txt       (via pytesseract)

\item {} 
\sphinxAtStartPar
Nettoyage du texte   \(\rightarrow\) \textless{}basename\textgreater{}.txt           (logique propre)

\end{enumerate}
\begin{description}
\sphinxlineitem{Ce module NE gère PAS :}\begin{itemize}
\item {} 
\sphinxAtStartPar
La synthèse vocale (TTS) : déléguée au Speaker central via
Speaker.synthesize\_to\_file(), ce qui évite de charger le modèle
Piper une deuxième fois en mémoire.

\item {} 
\sphinxAtStartPar
Le keypad / les touches : gérés par le mode (ModeReadingMachine).

\item {} 
\sphinxAtStartPar
Le playback / pause / seek : gérés par Speaker/mplayer côté core.

\item {} 
\sphinxAtStartPar
La logique de mode (on\_enter/on\_exit, mapping touches, etc.).

\end{itemize}

\end{description}


\subparagraph{Fichiers produits}
\label{\detokenize{api/hardware.devices:fichiers-produits}}\begin{description}
\sphinxlineitem{Pour un basename donné (ex}{[}« /tmp/readforme/scan ») :{]}\begin{itemize}
\item {} 
\sphinxAtStartPar
Image brute  : /tmp/readforme/scan.jpg

\item {} 
\sphinxAtStartPar
Texte OCR    : /tmp/readforme/scan\_raw.txt

\item {} 
\sphinxAtStartPar
Texte propre : /tmp/readforme/scan.txt

\end{itemize}

\end{description}

\sphinxAtStartPar
Le fichier WAV est produit EN DEHORS de ce module, par le mode qui
appelle Speaker.synthesize\_to\_file() puis Speaker.play().


\subparagraph{Pourquoi cette séparation ?}
\label{\detokenize{api/hardware.devices:pourquoi-cette-separation}}
\sphinxAtStartPar
L’ancien code embarquait le TTS (pico2wave, puis Piper) directement ici.
Cela impliquait de charger le modèle Piper une seconde fois, consommant
inutilement de la RAM sur le Raspberry Pi.
Désormais, un seul modèle Piper est chargé (dans Speaker.\_\_init\_\_) et
partagé via Speaker.synthesize\_to\_file().
\index{OCRConfig (classe dans hardware.devices.text\_reader\_device)@\spxentry{OCRConfig}\spxextra{classe dans hardware.devices.text\_reader\_device}}

\begin{fulllineitems}
\phantomsection\label{\detokenize{api/hardware.devices:hardware.devices.text_reader_device.OCRConfig}}
\pysigstartsignatures
\pysiglinewithargsret
{\sphinxbfcode{\sphinxupquote{\DUrole{k}{class}\DUrole{w}{ }}}\sphinxcode{\sphinxupquote{hardware.devices.text\_reader\_device.}}\sphinxbfcode{\sphinxupquote{OCRConfig}}}
{\sphinxparam{\DUrole{n}{lang}\DUrole{p}{:}\DUrole{w}{ }\DUrole{n}{str}\DUrole{w}{ }\DUrole{o}{=}\DUrole{w}{ }\DUrole{default_value}{\textquotesingle{}fra\textquotesingle{}}}\sphinxparamcomma \sphinxparam{\DUrole{n}{psm}\DUrole{p}{:}\DUrole{w}{ }\DUrole{n}{int}\DUrole{w}{ }\DUrole{o}{=}\DUrole{w}{ }\DUrole{default_value}{6}}\sphinxparamcomma \sphinxparam{\DUrole{n}{oem}\DUrole{p}{:}\DUrole{w}{ }\DUrole{n}{int}\DUrole{w}{ }\DUrole{o}{=}\DUrole{w}{ }\DUrole{default_value}{1}}\sphinxparamcomma \sphinxparam{\DUrole{n}{preprocess\_enabled}\DUrole{p}{:}\DUrole{w}{ }\DUrole{n}{bool}\DUrole{w}{ }\DUrole{o}{=}\DUrole{w}{ }\DUrole{default_value}{True}}\sphinxparamcomma \sphinxparam{\DUrole{n}{autocontrast\_cutoff}\DUrole{p}{:}\DUrole{w}{ }\DUrole{n}{int}\DUrole{w}{ }\DUrole{o}{=}\DUrole{w}{ }\DUrole{default_value}{2}}\sphinxparamcomma \sphinxparam{\DUrole{n}{sharpen\_factor}\DUrole{p}{:}\DUrole{w}{ }\DUrole{n}{float}\DUrole{w}{ }\DUrole{o}{=}\DUrole{w}{ }\DUrole{default_value}{1.2}}\sphinxparamcomma \sphinxparam{\DUrole{n}{threshold}\DUrole{p}{:}\DUrole{w}{ }\DUrole{n}{int}\DUrole{w}{ }\DUrole{o}{=}\DUrole{w}{ }\DUrole{default_value}{0}}\sphinxparamcomma \sphinxparam{\DUrole{n}{timeout\_sec}\DUrole{p}{:}\DUrole{w}{ }\DUrole{n}{float}\DUrole{w}{ }\DUrole{o}{=}\DUrole{w}{ }\DUrole{default_value}{20.0}}\sphinxparamcomma \sphinxparam{\DUrole{n}{fallback\_psm}\DUrole{p}{:}\DUrole{w}{ }\DUrole{n}{tuple\DUrole{p}{{[}}int\DUrole{p}{,}\DUrole{w}{ }\DUrole{p}{...}\DUrole{p}{{]}}}\DUrole{w}{ }\DUrole{o}{=}\DUrole{w}{ }\DUrole{default_value}{(7, 11, 3)}}\sphinxparamcomma \sphinxparam{\DUrole{n}{save\_debug\_ocr\_input}\DUrole{p}{:}\DUrole{w}{ }\DUrole{n}{bool}\DUrole{w}{ }\DUrole{o}{=}\DUrole{w}{ }\DUrole{default_value}{True}}}
{}
\pysigstopsignatures
\sphinxAtStartPar
Bases : \sphinxcode{\sphinxupquote{object}}

\sphinxAtStartPar
Paramètres de la reconnaissance optique de caractères (OCR).


\subparagraph{Attributs}
\label{\detokenize{api/hardware.devices:id6}}\begin{description}
\sphinxlineitem{lang :}
\sphinxAtStartPar
Code langue tesseract (ex : « fra » pour le français, « eng » pour
l’anglais). Doit correspondre à un pack de langues installé.

\sphinxlineitem{psm :}
\sphinxAtStartPar
Page Segmentation Mode tesseract (\textendash{}psm).
3 = détection automatique (valeur recommandée pour documents).
Voir : tesseract \textendash{}help\sphinxhyphen{}psm

\end{description}
\index{lang (attribut hardware.devices.text\_reader\_device.OCRConfig)@\spxentry{lang}\spxextra{attribut hardware.devices.text\_reader\_device.OCRConfig}}

\begin{fulllineitems}
\phantomsection\label{\detokenize{api/hardware.devices:hardware.devices.text_reader_device.OCRConfig.lang}}
\pysigstartsignatures
\pysigline
{\sphinxbfcode{\sphinxupquote{lang}}\sphinxbfcode{\sphinxupquote{\DUrole{p}{:}\DUrole{w}{ }str}}\sphinxbfcode{\sphinxupquote{\DUrole{w}{ }\DUrole{p}{=}\DUrole{w}{ }\textquotesingle{}fra\textquotesingle{}}}}
\pysigstopsignatures
\end{fulllineitems}

\index{psm (attribut hardware.devices.text\_reader\_device.OCRConfig)@\spxentry{psm}\spxextra{attribut hardware.devices.text\_reader\_device.OCRConfig}}

\begin{fulllineitems}
\phantomsection\label{\detokenize{api/hardware.devices:hardware.devices.text_reader_device.OCRConfig.psm}}
\pysigstartsignatures
\pysigline
{\sphinxbfcode{\sphinxupquote{psm}}\sphinxbfcode{\sphinxupquote{\DUrole{p}{:}\DUrole{w}{ }int}}\sphinxbfcode{\sphinxupquote{\DUrole{w}{ }\DUrole{p}{=}\DUrole{w}{ }6}}}
\pysigstopsignatures
\end{fulllineitems}

\index{oem (attribut hardware.devices.text\_reader\_device.OCRConfig)@\spxentry{oem}\spxextra{attribut hardware.devices.text\_reader\_device.OCRConfig}}

\begin{fulllineitems}
\phantomsection\label{\detokenize{api/hardware.devices:hardware.devices.text_reader_device.OCRConfig.oem}}
\pysigstartsignatures
\pysigline
{\sphinxbfcode{\sphinxupquote{oem}}\sphinxbfcode{\sphinxupquote{\DUrole{p}{:}\DUrole{w}{ }int}}\sphinxbfcode{\sphinxupquote{\DUrole{w}{ }\DUrole{p}{=}\DUrole{w}{ }1}}}
\pysigstopsignatures
\end{fulllineitems}

\index{preprocess\_enabled (attribut hardware.devices.text\_reader\_device.OCRConfig)@\spxentry{preprocess\_enabled}\spxextra{attribut hardware.devices.text\_reader\_device.OCRConfig}}

\begin{fulllineitems}
\phantomsection\label{\detokenize{api/hardware.devices:hardware.devices.text_reader_device.OCRConfig.preprocess_enabled}}
\pysigstartsignatures
\pysigline
{\sphinxbfcode{\sphinxupquote{preprocess\_enabled}}\sphinxbfcode{\sphinxupquote{\DUrole{p}{:}\DUrole{w}{ }bool}}\sphinxbfcode{\sphinxupquote{\DUrole{w}{ }\DUrole{p}{=}\DUrole{w}{ }True}}}
\pysigstopsignatures
\end{fulllineitems}

\index{autocontrast\_cutoff (attribut hardware.devices.text\_reader\_device.OCRConfig)@\spxentry{autocontrast\_cutoff}\spxextra{attribut hardware.devices.text\_reader\_device.OCRConfig}}

\begin{fulllineitems}
\phantomsection\label{\detokenize{api/hardware.devices:hardware.devices.text_reader_device.OCRConfig.autocontrast_cutoff}}
\pysigstartsignatures
\pysigline
{\sphinxbfcode{\sphinxupquote{autocontrast\_cutoff}}\sphinxbfcode{\sphinxupquote{\DUrole{p}{:}\DUrole{w}{ }int}}\sphinxbfcode{\sphinxupquote{\DUrole{w}{ }\DUrole{p}{=}\DUrole{w}{ }2}}}
\pysigstopsignatures
\end{fulllineitems}

\index{sharpen\_factor (attribut hardware.devices.text\_reader\_device.OCRConfig)@\spxentry{sharpen\_factor}\spxextra{attribut hardware.devices.text\_reader\_device.OCRConfig}}

\begin{fulllineitems}
\phantomsection\label{\detokenize{api/hardware.devices:hardware.devices.text_reader_device.OCRConfig.sharpen_factor}}
\pysigstartsignatures
\pysigline
{\sphinxbfcode{\sphinxupquote{sharpen\_factor}}\sphinxbfcode{\sphinxupquote{\DUrole{p}{:}\DUrole{w}{ }float}}\sphinxbfcode{\sphinxupquote{\DUrole{w}{ }\DUrole{p}{=}\DUrole{w}{ }1.2}}}
\pysigstopsignatures
\end{fulllineitems}

\index{threshold (attribut hardware.devices.text\_reader\_device.OCRConfig)@\spxentry{threshold}\spxextra{attribut hardware.devices.text\_reader\_device.OCRConfig}}

\begin{fulllineitems}
\phantomsection\label{\detokenize{api/hardware.devices:hardware.devices.text_reader_device.OCRConfig.threshold}}
\pysigstartsignatures
\pysigline
{\sphinxbfcode{\sphinxupquote{threshold}}\sphinxbfcode{\sphinxupquote{\DUrole{p}{:}\DUrole{w}{ }int}}\sphinxbfcode{\sphinxupquote{\DUrole{w}{ }\DUrole{p}{=}\DUrole{w}{ }0}}}
\pysigstopsignatures
\end{fulllineitems}

\index{timeout\_sec (attribut hardware.devices.text\_reader\_device.OCRConfig)@\spxentry{timeout\_sec}\spxextra{attribut hardware.devices.text\_reader\_device.OCRConfig}}

\begin{fulllineitems}
\phantomsection\label{\detokenize{api/hardware.devices:hardware.devices.text_reader_device.OCRConfig.timeout_sec}}
\pysigstartsignatures
\pysigline
{\sphinxbfcode{\sphinxupquote{timeout\_sec}}\sphinxbfcode{\sphinxupquote{\DUrole{p}{:}\DUrole{w}{ }float}}\sphinxbfcode{\sphinxupquote{\DUrole{w}{ }\DUrole{p}{=}\DUrole{w}{ }20.0}}}
\pysigstopsignatures
\end{fulllineitems}

\index{fallback\_psm (attribut hardware.devices.text\_reader\_device.OCRConfig)@\spxentry{fallback\_psm}\spxextra{attribut hardware.devices.text\_reader\_device.OCRConfig}}

\begin{fulllineitems}
\phantomsection\label{\detokenize{api/hardware.devices:hardware.devices.text_reader_device.OCRConfig.fallback_psm}}
\pysigstartsignatures
\pysigline
{\sphinxbfcode{\sphinxupquote{fallback\_psm}}\sphinxbfcode{\sphinxupquote{\DUrole{p}{:}\DUrole{w}{ }tuple\DUrole{p}{{[}}int\DUrole{p}{,}\DUrole{w}{ }\DUrole{p}{...}\DUrole{p}{{]}}}}\sphinxbfcode{\sphinxupquote{\DUrole{w}{ }\DUrole{p}{=}\DUrole{w}{ }(7, 11, 3)}}}
\pysigstopsignatures
\end{fulllineitems}

\index{save\_debug\_ocr\_input (attribut hardware.devices.text\_reader\_device.OCRConfig)@\spxentry{save\_debug\_ocr\_input}\spxextra{attribut hardware.devices.text\_reader\_device.OCRConfig}}

\begin{fulllineitems}
\phantomsection\label{\detokenize{api/hardware.devices:hardware.devices.text_reader_device.OCRConfig.save_debug_ocr_input}}
\pysigstartsignatures
\pysigline
{\sphinxbfcode{\sphinxupquote{save\_debug\_ocr\_input}}\sphinxbfcode{\sphinxupquote{\DUrole{p}{:}\DUrole{w}{ }bool}}\sphinxbfcode{\sphinxupquote{\DUrole{w}{ }\DUrole{p}{=}\DUrole{w}{ }True}}}
\pysigstopsignatures
\end{fulllineitems}


\end{fulllineitems}

\index{TextReaderDevice (classe dans hardware.devices.text\_reader\_device)@\spxentry{TextReaderDevice}\spxextra{classe dans hardware.devices.text\_reader\_device}}

\begin{fulllineitems}
\phantomsection\label{\detokenize{api/hardware.devices:hardware.devices.text_reader_device.TextReaderDevice}}
\pysigstartsignatures
\pysiglinewithargsret
{\sphinxbfcode{\sphinxupquote{\DUrole{k}{class}\DUrole{w}{ }}}\sphinxcode{\sphinxupquote{hardware.devices.text\_reader\_device.}}\sphinxbfcode{\sphinxupquote{TextReaderDevice}}}
{\sphinxparam{\DUrole{n}{camera}\DUrole{p}{:}\DUrole{w}{ }\DUrole{n}{{\hyperref[\detokenize{api/hardware.platform:hardware.platform.pi_camera.PiCamera}]{\sphinxcrossref{PiCamera}}}}}\sphinxparamcomma \sphinxparam{\DUrole{keyword-only-separator}{\DUrole{o}{\sphinxstyleabbreviation{*}}}}\sphinxparamcomma \sphinxparam{\DUrole{n}{ocr}\DUrole{p}{:}\DUrole{w}{ }\DUrole{n}{{\hyperref[\detokenize{api/hardware.devices:hardware.devices.text_reader_device.OCRConfig}]{\sphinxcrossref{OCRConfig}}}\DUrole{w}{ }\DUrole{p}{|}\DUrole{w}{ }None}\DUrole{w}{ }\DUrole{o}{=}\DUrole{w}{ }\DUrole{default_value}{None}}}
{}
\pysigstopsignatures
\sphinxAtStartPar
Bases : \sphinxcode{\sphinxupquote{object}}

\sphinxAtStartPar
Device « machine à lire » — pipeline capture \(\rightarrow\) OCR \(\rightarrow\) texte propre.

\sphinxAtStartPar
Ce device s’arrête à la production du texte nettoyé.
La synthèse vocale est gérée par le Speaker central.


\subparagraph{API principale}
\label{\detokenize{api/hardware.devices:api-principale}}
\sphinxAtStartPar
capture(basename)         \(\rightarrow\) str   chemin de l’image
ocr\_to\_text(basename)     \(\rightarrow\) str   chemin du texte OCR brut
clean\_text(basename)      \(\rightarrow\) str   chemin du texte nettoyé
run\_full\_pipeline(basename) \(\rightarrow\) str   chemin du texte nettoyé (enchaîne les 3)
\begin{description}
\sphinxlineitem{Convention de nommage (basename = « /tmp/readforme/scan ») :}
\sphinxAtStartPar
Image brute  : /tmp/readforme/scan.jpg
Texte OCR    : /tmp/readforme/scan\_raw.txt
Texte propre : /tmp/readforme/scan.txt

\end{description}
\index{\_\_init\_\_() (méthode hardware.devices.text\_reader\_device.TextReaderDevice)@\spxentry{\_\_init\_\_()}\spxextra{méthode hardware.devices.text\_reader\_device.TextReaderDevice}}

\begin{fulllineitems}
\phantomsection\label{\detokenize{api/hardware.devices:hardware.devices.text_reader_device.TextReaderDevice.__init__}}
\pysigstartsignatures
\pysiglinewithargsret
{\sphinxbfcode{\sphinxupquote{\_\_init\_\_}}}
{\sphinxparam{\DUrole{n}{camera}\DUrole{p}{:}\DUrole{w}{ }\DUrole{n}{{\hyperref[\detokenize{api/hardware.platform:hardware.platform.pi_camera.PiCamera}]{\sphinxcrossref{PiCamera}}}}}\sphinxparamcomma \sphinxparam{\DUrole{keyword-only-separator}{\DUrole{o}{\sphinxstyleabbreviation{*}}}}\sphinxparamcomma \sphinxparam{\DUrole{n}{ocr}\DUrole{p}{:}\DUrole{w}{ }\DUrole{n}{{\hyperref[\detokenize{api/hardware.devices:hardware.devices.text_reader_device.OCRConfig}]{\sphinxcrossref{OCRConfig}}}\DUrole{w}{ }\DUrole{p}{|}\DUrole{w}{ }None}\DUrole{w}{ }\DUrole{o}{=}\DUrole{w}{ }\DUrole{default_value}{None}}}
{{ $\rightarrow$ None}}
\pysigstopsignatures
\sphinxAtStartPar
Initialise le device.


\subparagraph{Paramètres}
\label{\detokenize{api/hardware.devices:id7}}\begin{description}
\sphinxlineitem{camera :}
\sphinxAtStartPar
Instance de PiCamera configurée (rotation, résolution…).

\sphinxlineitem{ocr :}
\sphinxAtStartPar
Configuration OCR. Si None, utilise les valeurs par défaut
(lang= »fra », psm=3).

\end{description}

\end{fulllineitems}

\index{capture() (méthode hardware.devices.text\_reader\_device.TextReaderDevice)@\spxentry{capture()}\spxextra{méthode hardware.devices.text\_reader\_device.TextReaderDevice}}

\begin{fulllineitems}
\phantomsection\label{\detokenize{api/hardware.devices:hardware.devices.text_reader_device.TextReaderDevice.capture}}
\pysigstartsignatures
\pysiglinewithargsret
{\sphinxbfcode{\sphinxupquote{capture}}}
{\sphinxparam{\DUrole{n}{basename}\DUrole{p}{:}\DUrole{w}{ }\DUrole{n}{str}}}
{{ $\rightarrow$ str}}
\pysigstopsignatures
\sphinxAtStartPar
Capture une photo et écrit \textless{}basename\textgreater{}.jpg.


\subparagraph{Paramètres}
\label{\detokenize{api/hardware.devices:id8}}\begin{description}
\sphinxlineitem{basename :}
\sphinxAtStartPar
Chemin sans extension (ex : « /tmp/readforme/scan »).

\end{description}


\subparagraph{Retourne}
\label{\detokenize{api/hardware.devices:retourne}}\begin{description}
\sphinxlineitem{str}
\sphinxAtStartPar
Chemin complet du fichier image créé.

\end{description}

\end{fulllineitems}

\index{ocr\_to\_text() (méthode hardware.devices.text\_reader\_device.TextReaderDevice)@\spxentry{ocr\_to\_text()}\spxextra{méthode hardware.devices.text\_reader\_device.TextReaderDevice}}

\begin{fulllineitems}
\phantomsection\label{\detokenize{api/hardware.devices:hardware.devices.text_reader_device.TextReaderDevice.ocr_to_text}}
\pysigstartsignatures
\pysiglinewithargsret
{\sphinxbfcode{\sphinxupquote{ocr\_to\_text}}}
{\sphinxparam{\DUrole{n}{basename}\DUrole{p}{:}\DUrole{w}{ }\DUrole{n}{str}}}
{{ $\rightarrow$ str}}
\pysigstopsignatures
\sphinxAtStartPar
Effectue l’OCR sur \textless{}basename\textgreater{}.jpg et écrit \textless{}basename\textgreater{}\_raw.txt.

\sphinxAtStartPar
Utilise pytesseract avec la langue et le PSM configurés.


\subparagraph{Paramètres}
\label{\detokenize{api/hardware.devices:id9}}\begin{description}
\sphinxlineitem{basename :}
\sphinxAtStartPar
Chemin sans extension.

\end{description}


\subparagraph{Retourne}
\label{\detokenize{api/hardware.devices:id10}}\begin{description}
\sphinxlineitem{str}
\sphinxAtStartPar
Chemin du fichier texte brut produit.

\end{description}
\begin{quote}\begin{description}
\sphinxlineitem{raises FileNotFoundError}
\sphinxAtStartPar
Si l’image n’existe pas (étape capture non réalisée).

\end{description}\end{quote}

\end{fulllineitems}

\index{clean\_text() (méthode hardware.devices.text\_reader\_device.TextReaderDevice)@\spxentry{clean\_text()}\spxextra{méthode hardware.devices.text\_reader\_device.TextReaderDevice}}

\begin{fulllineitems}
\phantomsection\label{\detokenize{api/hardware.devices:hardware.devices.text_reader_device.TextReaderDevice.clean_text}}
\pysigstartsignatures
\pysiglinewithargsret
{\sphinxbfcode{\sphinxupquote{clean\_text}}}
{\sphinxparam{\DUrole{n}{basename}\DUrole{p}{:}\DUrole{w}{ }\DUrole{n}{str}}}
{{ $\rightarrow$ str}}
\pysigstopsignatures
\sphinxAtStartPar
Nettoie le texte OCR brut et écrit \textless{}basename\textgreater{}.txt.

\sphinxAtStartPar
Stratégie de nettoyage (identique à l’ancien projet) :
\sphinxhyphen{} Les lignes non vides sont concaténées avec des espaces.
\sphinxhyphen{} Une ligne vide introduit un saut de paragraphe (double newline).
\sphinxhyphen{} Les traits d’union de fin de ligne ( »\sphinxhyphen{} n ») sont supprimés
\begin{quote}

\sphinxAtStartPar
(gestion des mots coupés par l’OCR).
\end{quote}


\subparagraph{Paramètres}
\label{\detokenize{api/hardware.devices:id11}}\begin{description}
\sphinxlineitem{basename :}
\sphinxAtStartPar
Chemin sans extension.

\end{description}


\subparagraph{Retourne}
\label{\detokenize{api/hardware.devices:id12}}\begin{description}
\sphinxlineitem{str}
\sphinxAtStartPar
Chemin du fichier texte nettoyé.

\end{description}
\begin{quote}\begin{description}
\sphinxlineitem{raises FileNotFoundError}
\sphinxAtStartPar
Si le texte brut n’existe pas (étape ocr\_to\_text non réalisée).

\sphinxlineitem{raises RuntimeError}
\sphinxAtStartPar
Si le texte nettoyé est vide (document illisible ou page vide).

\end{description}\end{quote}

\end{fulllineitems}

\index{run\_full\_pipeline() (méthode hardware.devices.text\_reader\_device.TextReaderDevice)@\spxentry{run\_full\_pipeline()}\spxextra{méthode hardware.devices.text\_reader\_device.TextReaderDevice}}

\begin{fulllineitems}
\phantomsection\label{\detokenize{api/hardware.devices:hardware.devices.text_reader_device.TextReaderDevice.run_full_pipeline}}
\pysigstartsignatures
\pysiglinewithargsret
{\sphinxbfcode{\sphinxupquote{run\_full\_pipeline}}}
{\sphinxparam{\DUrole{n}{basename}\DUrole{p}{:}\DUrole{w}{ }\DUrole{n}{str}}}
{{ $\rightarrow$ str}}
\pysigstopsignatures
\sphinxAtStartPar
Exécute la chaîne complète : capture \(\rightarrow\) OCR \(\rightarrow\) nettoyage.

\sphinxAtStartPar
Équivalent à appeler capture(), ocr\_to\_text(), clean\_text()
dans l’ordre, avec gestion d’erreurs centralisée.


\subparagraph{Paramètres}
\label{\detokenize{api/hardware.devices:id13}}\begin{description}
\sphinxlineitem{basename :}
\sphinxAtStartPar
Chemin sans extension (ex : « /tmp/readforme/scan »).

\end{description}


\subparagraph{Retourne}
\label{\detokenize{api/hardware.devices:id14}}\begin{description}
\sphinxlineitem{str}
\sphinxAtStartPar
Chemin du fichier texte nettoyé (\textless{}basename\textgreater{}.txt).
Le contenu est prêt à être synthétisé via
Speaker.synthesize\_to\_file().

\end{description}
\begin{quote}\begin{description}
\sphinxlineitem{raises FileNotFoundError}
\sphinxAtStartPar
Si un fichier intermédiaire n’existe pas.

\sphinxlineitem{raises RuntimeError}
\sphinxAtStartPar
Si le texte résultant est vide.

\end{description}\end{quote}

\end{fulllineitems}


\end{fulllineitems}



\paragraph{Module contents}
\label{\detokenize{api/hardware.devices:module-hardware.devices}}\label{\detokenize{api/hardware.devices:module-contents}}\index{module@\spxentry{module}!hardware.devices@\spxentry{hardware.devices}}\index{hardware.devices@\spxentry{hardware.devices}!module@\spxentry{module}}
\sphinxstepscope


\subsubsection{hardware.platform package}
\label{\detokenize{api/hardware.platform:hardware-platform-package}}\label{\detokenize{api/hardware.platform::doc}}

\paragraph{Submodules}
\label{\detokenize{api/hardware.platform:submodules}}

\paragraph{hardware.platform.keypad\_4x4 module}
\label{\detokenize{api/hardware.platform:module-hardware.platform.keypad_4x4}}\label{\detokenize{api/hardware.platform:hardware-platform-keypad-4x4-module}}\index{module@\spxentry{module}!hardware.platform.keypad\_4x4@\spxentry{hardware.platform.keypad\_4x4}}\index{hardware.platform.keypad\_4x4@\spxentry{hardware.platform.keypad\_4x4}!module@\spxentry{module}}
\sphinxAtStartPar
Driver clavier matriciel 4x4 (GPIO) pour Raspberry Pi.


\subparagraph{Responsabilité}
\label{\detokenize{api/hardware.platform:responsabilite}}\begin{itemize}
\item {} 
\sphinxAtStartPar
Scanner un clavier 4x4 câblé en matrice (4 lignes / 4 colonnes)

\item {} 
\sphinxAtStartPar
Déclencher un callback quand une touche est pressée

\item {} 
\sphinxAtStartPar
Fournir start/stop pour activer ce périphérique uniquement dans certains modes

\end{itemize}


\subparagraph{Choix d’API (propre pour l’archi)}
\label{\detokenize{api/hardware.platform:choix-d-api-propre-pour-l-archi}}\begin{itemize}
\item {} 
\sphinxAtStartPar
Keypad4x4.start()  : démarre un thread de scan (non\sphinxhyphen{}bloquant)

\item {} 
\sphinxAtStartPar
Keypad4x4.stop()   : stoppe le thread proprement

\item {} 
\sphinxAtStartPar
Keypad4x4.on\_key   : callback optionnel appelé quand une touche est détectée

\end{itemize}
\index{KeypadPins (classe dans hardware.platform.keypad\_4x4)@\spxentry{KeypadPins}\spxextra{classe dans hardware.platform.keypad\_4x4}}

\begin{fulllineitems}
\phantomsection\label{\detokenize{api/hardware.platform:hardware.platform.keypad_4x4.KeypadPins}}
\pysigstartsignatures
\pysiglinewithargsret
{\sphinxbfcode{\sphinxupquote{\DUrole{k}{class}\DUrole{w}{ }}}\sphinxcode{\sphinxupquote{hardware.platform.keypad\_4x4.}}\sphinxbfcode{\sphinxupquote{KeypadPins}}}
{\sphinxparam{\DUrole{n}{row1}\DUrole{p}{:}\DUrole{w}{ }\DUrole{n}{\DUrole{s}{\textquotesingle{}int\textquotesingle{}}}}\sphinxparamcomma \sphinxparam{\DUrole{n}{row2}\DUrole{p}{:}\DUrole{w}{ }\DUrole{n}{\DUrole{s}{\textquotesingle{}int\textquotesingle{}}}}\sphinxparamcomma \sphinxparam{\DUrole{n}{row3}\DUrole{p}{:}\DUrole{w}{ }\DUrole{n}{\DUrole{s}{\textquotesingle{}int\textquotesingle{}}}}\sphinxparamcomma \sphinxparam{\DUrole{n}{row4}\DUrole{p}{:}\DUrole{w}{ }\DUrole{n}{\DUrole{s}{\textquotesingle{}int\textquotesingle{}}}}\sphinxparamcomma \sphinxparam{\DUrole{n}{col1}\DUrole{p}{:}\DUrole{w}{ }\DUrole{n}{\DUrole{s}{\textquotesingle{}int\textquotesingle{}}}}\sphinxparamcomma \sphinxparam{\DUrole{n}{col2}\DUrole{p}{:}\DUrole{w}{ }\DUrole{n}{\DUrole{s}{\textquotesingle{}int\textquotesingle{}}}}\sphinxparamcomma \sphinxparam{\DUrole{n}{col3}\DUrole{p}{:}\DUrole{w}{ }\DUrole{n}{\DUrole{s}{\textquotesingle{}int\textquotesingle{}}}}\sphinxparamcomma \sphinxparam{\DUrole{n}{col4}\DUrole{p}{:}\DUrole{w}{ }\DUrole{n}{\DUrole{s}{\textquotesingle{}int\textquotesingle{}}}}}
{}
\pysigstopsignatures
\sphinxAtStartPar
Bases : \sphinxcode{\sphinxupquote{object}}
\index{row1 (attribut hardware.platform.keypad\_4x4.KeypadPins)@\spxentry{row1}\spxextra{attribut hardware.platform.keypad\_4x4.KeypadPins}}

\begin{fulllineitems}
\phantomsection\label{\detokenize{api/hardware.platform:hardware.platform.keypad_4x4.KeypadPins.row1}}
\pysigstartsignatures
\pysigline
{\sphinxbfcode{\sphinxupquote{row1}}\sphinxbfcode{\sphinxupquote{\DUrole{p}{:}\DUrole{w}{ }int}}}
\pysigstopsignatures
\end{fulllineitems}

\index{row2 (attribut hardware.platform.keypad\_4x4.KeypadPins)@\spxentry{row2}\spxextra{attribut hardware.platform.keypad\_4x4.KeypadPins}}

\begin{fulllineitems}
\phantomsection\label{\detokenize{api/hardware.platform:hardware.platform.keypad_4x4.KeypadPins.row2}}
\pysigstartsignatures
\pysigline
{\sphinxbfcode{\sphinxupquote{row2}}\sphinxbfcode{\sphinxupquote{\DUrole{p}{:}\DUrole{w}{ }int}}}
\pysigstopsignatures
\end{fulllineitems}

\index{row3 (attribut hardware.platform.keypad\_4x4.KeypadPins)@\spxentry{row3}\spxextra{attribut hardware.platform.keypad\_4x4.KeypadPins}}

\begin{fulllineitems}
\phantomsection\label{\detokenize{api/hardware.platform:hardware.platform.keypad_4x4.KeypadPins.row3}}
\pysigstartsignatures
\pysigline
{\sphinxbfcode{\sphinxupquote{row3}}\sphinxbfcode{\sphinxupquote{\DUrole{p}{:}\DUrole{w}{ }int}}}
\pysigstopsignatures
\end{fulllineitems}

\index{row4 (attribut hardware.platform.keypad\_4x4.KeypadPins)@\spxentry{row4}\spxextra{attribut hardware.platform.keypad\_4x4.KeypadPins}}

\begin{fulllineitems}
\phantomsection\label{\detokenize{api/hardware.platform:hardware.platform.keypad_4x4.KeypadPins.row4}}
\pysigstartsignatures
\pysigline
{\sphinxbfcode{\sphinxupquote{row4}}\sphinxbfcode{\sphinxupquote{\DUrole{p}{:}\DUrole{w}{ }int}}}
\pysigstopsignatures
\end{fulllineitems}

\index{col1 (attribut hardware.platform.keypad\_4x4.KeypadPins)@\spxentry{col1}\spxextra{attribut hardware.platform.keypad\_4x4.KeypadPins}}

\begin{fulllineitems}
\phantomsection\label{\detokenize{api/hardware.platform:hardware.platform.keypad_4x4.KeypadPins.col1}}
\pysigstartsignatures
\pysigline
{\sphinxbfcode{\sphinxupquote{col1}}\sphinxbfcode{\sphinxupquote{\DUrole{p}{:}\DUrole{w}{ }int}}}
\pysigstopsignatures
\end{fulllineitems}

\index{col2 (attribut hardware.platform.keypad\_4x4.KeypadPins)@\spxentry{col2}\spxextra{attribut hardware.platform.keypad\_4x4.KeypadPins}}

\begin{fulllineitems}
\phantomsection\label{\detokenize{api/hardware.platform:hardware.platform.keypad_4x4.KeypadPins.col2}}
\pysigstartsignatures
\pysigline
{\sphinxbfcode{\sphinxupquote{col2}}\sphinxbfcode{\sphinxupquote{\DUrole{p}{:}\DUrole{w}{ }int}}}
\pysigstopsignatures
\end{fulllineitems}

\index{col3 (attribut hardware.platform.keypad\_4x4.KeypadPins)@\spxentry{col3}\spxextra{attribut hardware.platform.keypad\_4x4.KeypadPins}}

\begin{fulllineitems}
\phantomsection\label{\detokenize{api/hardware.platform:hardware.platform.keypad_4x4.KeypadPins.col3}}
\pysigstartsignatures
\pysigline
{\sphinxbfcode{\sphinxupquote{col3}}\sphinxbfcode{\sphinxupquote{\DUrole{p}{:}\DUrole{w}{ }int}}}
\pysigstopsignatures
\end{fulllineitems}

\index{col4 (attribut hardware.platform.keypad\_4x4.KeypadPins)@\spxentry{col4}\spxextra{attribut hardware.platform.keypad\_4x4.KeypadPins}}

\begin{fulllineitems}
\phantomsection\label{\detokenize{api/hardware.platform:hardware.platform.keypad_4x4.KeypadPins.col4}}
\pysigstartsignatures
\pysigline
{\sphinxbfcode{\sphinxupquote{col4}}\sphinxbfcode{\sphinxupquote{\DUrole{p}{:}\DUrole{w}{ }int}}}
\pysigstopsignatures
\end{fulllineitems}


\end{fulllineitems}

\index{Keypad4x4 (classe dans hardware.platform.keypad\_4x4)@\spxentry{Keypad4x4}\spxextra{classe dans hardware.platform.keypad\_4x4}}

\begin{fulllineitems}
\phantomsection\label{\detokenize{api/hardware.platform:hardware.platform.keypad_4x4.Keypad4x4}}
\pysigstartsignatures
\pysiglinewithargsret
{\sphinxbfcode{\sphinxupquote{\DUrole{k}{class}\DUrole{w}{ }}}\sphinxcode{\sphinxupquote{hardware.platform.keypad\_4x4.}}\sphinxbfcode{\sphinxupquote{Keypad4x4}}}
{\sphinxparam{\DUrole{n}{pins}\DUrole{p}{:}\DUrole{w}{ }\DUrole{n}{{\hyperref[\detokenize{api/hardware.platform:hardware.platform.keypad_4x4.KeypadPins}]{\sphinxcrossref{KeypadPins}}}}}\sphinxparamcomma \sphinxparam{\DUrole{keyword-only-separator}{\DUrole{o}{\sphinxstyleabbreviation{*}}}}\sphinxparamcomma \sphinxparam{\DUrole{n}{layout}\DUrole{p}{:}\DUrole{w}{ }\DUrole{n}{List\DUrole{p}{{[}}List\DUrole{p}{{[}}str\DUrole{p}{{]}}\DUrole{p}{{]}}\DUrole{w}{ }\DUrole{p}{|}\DUrole{w}{ }None}\DUrole{w}{ }\DUrole{o}{=}\DUrole{w}{ }\DUrole{default_value}{None}}\sphinxparamcomma \sphinxparam{\DUrole{n}{scan\_interval}\DUrole{p}{:}\DUrole{w}{ }\DUrole{n}{float}\DUrole{w}{ }\DUrole{o}{=}\DUrole{w}{ }\DUrole{default_value}{0.02}}\sphinxparamcomma \sphinxparam{\DUrole{n}{debounce\_time}\DUrole{p}{:}\DUrole{w}{ }\DUrole{n}{float}\DUrole{w}{ }\DUrole{o}{=}\DUrole{w}{ }\DUrole{default_value}{0.2}}}
{}
\pysigstopsignatures
\sphinxAtStartPar
Bases : \sphinxcode{\sphinxupquote{object}}

\sphinxAtStartPar
Driver de scan clavier 4x4.
\begin{itemize}
\item {} 
\sphinxAtStartPar
Scan cyclique des lignes (row) \sphinxhyphen{}\textgreater{} lecture colonnes (col)

\item {} 
\sphinxAtStartPar
Anti\sphinxhyphen{}rebond logiciel simple

\item {} 
\sphinxAtStartPar
Détection « press » (pas maintien)

\end{itemize}
\index{DEFAULT\_LAYOUT (attribut hardware.platform.keypad\_4x4.Keypad4x4)@\spxentry{DEFAULT\_LAYOUT}\spxextra{attribut hardware.platform.keypad\_4x4.Keypad4x4}}

\begin{fulllineitems}
\phantomsection\label{\detokenize{api/hardware.platform:hardware.platform.keypad_4x4.Keypad4x4.DEFAULT_LAYOUT}}
\pysigstartsignatures
\pysigline
{\sphinxbfcode{\sphinxupquote{DEFAULT\_LAYOUT}}\sphinxbfcode{\sphinxupquote{\DUrole{p}{:}\DUrole{w}{ }List\DUrole{p}{{[}}List\DUrole{p}{{[}}str\DUrole{p}{{]}}\DUrole{p}{{]}}}}\sphinxbfcode{\sphinxupquote{\DUrole{w}{ }\DUrole{p}{=}\DUrole{w}{ }{[}{[}\textquotesingle{}1\textquotesingle{}, \textquotesingle{}2\textquotesingle{}, \textquotesingle{}3\textquotesingle{}, \textquotesingle{}A\textquotesingle{}{]}, {[}\textquotesingle{}4\textquotesingle{}, \textquotesingle{}5\textquotesingle{}, \textquotesingle{}6\textquotesingle{}, \textquotesingle{}B\textquotesingle{}{]}, {[}\textquotesingle{}7\textquotesingle{}, \textquotesingle{}8\textquotesingle{}, \textquotesingle{}9\textquotesingle{}, \textquotesingle{}C\textquotesingle{}{]}, {[}\textquotesingle{}*\textquotesingle{}, \textquotesingle{}0\textquotesingle{}, \textquotesingle{}\#\textquotesingle{}, \textquotesingle{}D\textquotesingle{}{]}{]}}}}
\pysigstopsignatures
\end{fulllineitems}

\index{on\_key (attribut hardware.platform.keypad\_4x4.Keypad4x4)@\spxentry{on\_key}\spxextra{attribut hardware.platform.keypad\_4x4.Keypad4x4}}

\begin{fulllineitems}
\phantomsection\label{\detokenize{api/hardware.platform:hardware.platform.keypad_4x4.Keypad4x4.on_key}}
\pysigstartsignatures
\pysigline
{\sphinxbfcode{\sphinxupquote{on\_key}}\sphinxbfcode{\sphinxupquote{\DUrole{p}{:}\DUrole{w}{ }Callable\DUrole{p}{{[}}\DUrole{p}{{[}}str\DUrole{p}{{]}}\DUrole{p}{,}\DUrole{w}{ }None\DUrole{p}{{]}}\DUrole{w}{ }\DUrole{p}{|}\DUrole{w}{ }None}}}
\pysigstopsignatures
\end{fulllineitems}

\index{start() (méthode hardware.platform.keypad\_4x4.Keypad4x4)@\spxentry{start()}\spxextra{méthode hardware.platform.keypad\_4x4.Keypad4x4}}

\begin{fulllineitems}
\phantomsection\label{\detokenize{api/hardware.platform:hardware.platform.keypad_4x4.Keypad4x4.start}}
\pysigstartsignatures
\pysiglinewithargsret
{\sphinxbfcode{\sphinxupquote{start}}}
{}
{{ $\rightarrow$ None}}
\pysigstopsignatures
\sphinxAtStartPar
Démarre le scan clavier dans un thread.

\end{fulllineitems}

\index{stop() (méthode hardware.platform.keypad\_4x4.Keypad4x4)@\spxentry{stop()}\spxextra{méthode hardware.platform.keypad\_4x4.Keypad4x4}}

\begin{fulllineitems}
\phantomsection\label{\detokenize{api/hardware.platform:hardware.platform.keypad_4x4.Keypad4x4.stop}}
\pysigstartsignatures
\pysiglinewithargsret
{\sphinxbfcode{\sphinxupquote{stop}}}
{}
{{ $\rightarrow$ None}}
\pysigstopsignatures
\sphinxAtStartPar
Stoppe le scan clavier.

\end{fulllineitems}


\end{fulllineitems}



\paragraph{hardware.platform.pi\_camera module}
\label{\detokenize{api/hardware.platform:module-hardware.platform.pi_camera}}\label{\detokenize{api/hardware.platform:hardware-platform-pi-camera-module}}\index{module@\spxentry{module}!hardware.platform.pi\_camera@\spxentry{hardware.platform.pi\_camera}}\index{hardware.platform.pi\_camera@\spxentry{hardware.platform.pi\_camera}!module@\spxentry{module}}

\subparagraph{hardware.platform.pi\_camera}
\label{\detokenize{api/hardware.platform:hardware-platform-pi-camera}}
\sphinxAtStartPar
Driver bas niveau pour la caméra Raspberry Pi (libcamera / rpicam).


\subparagraph{Responsabilité}
\label{\detokenize{api/hardware.platform:id1}}\begin{itemize}
\item {} 
\sphinxAtStartPar
Capturer une image via \sphinxtitleref{rpicam\sphinxhyphen{}still} (Bookworm) ou \sphinxtitleref{libcamera\sphinxhyphen{}still} (legacy)

\item {} 
\sphinxAtStartPar
Ne contient AUCUNE logique métier (pas d’OCR, pas de TTS, pas de mode)

\end{itemize}


\subparagraph{API}
\label{\detokenize{api/hardware.platform:api}}\begin{itemize}
\item {} \begin{description}
\sphinxlineitem{PiCamera.capture(output\_path) \sphinxhyphen{}\textgreater{} str}
\sphinxAtStartPar
Capture une photo et l’écrit dans output\_path.

\end{description}

\end{itemize}
\subsubsection*{Notes}
\begin{itemize}
\item {} 
\sphinxAtStartPar
Utilise \sphinxtitleref{subprocess.run} (plus fiable que os.system).

\item {} 
\sphinxAtStartPar
Lève des exceptions explicites en cas d’erreur.

\end{itemize}
\index{CameraConfig (classe dans hardware.platform.pi\_camera)@\spxentry{CameraConfig}\spxextra{classe dans hardware.platform.pi\_camera}}

\begin{fulllineitems}
\phantomsection\label{\detokenize{api/hardware.platform:hardware.platform.pi_camera.CameraConfig}}
\pysigstartsignatures
\pysiglinewithargsret
{\sphinxbfcode{\sphinxupquote{\DUrole{k}{class}\DUrole{w}{ }}}\sphinxcode{\sphinxupquote{hardware.platform.pi\_camera.}}\sphinxbfcode{\sphinxupquote{CameraConfig}}}
{\sphinxparam{\DUrole{n}{rotation}\DUrole{p}{:}\DUrole{w}{ }\DUrole{n}{int}\DUrole{w}{ }\DUrole{o}{=}\DUrole{w}{ }\DUrole{default_value}{180}}\sphinxparamcomma \sphinxparam{\DUrole{n}{timeout\_ms}\DUrole{p}{:}\DUrole{w}{ }\DUrole{n}{int}\DUrole{w}{ }\DUrole{o}{=}\DUrole{w}{ }\DUrole{default_value}{500}}\sphinxparamcomma \sphinxparam{\DUrole{n}{width}\DUrole{p}{:}\DUrole{w}{ }\DUrole{n}{int\DUrole{w}{ }\DUrole{p}{|}\DUrole{w}{ }None}\DUrole{w}{ }\DUrole{o}{=}\DUrole{w}{ }\DUrole{default_value}{None}}\sphinxparamcomma \sphinxparam{\DUrole{n}{height}\DUrole{p}{:}\DUrole{w}{ }\DUrole{n}{int\DUrole{w}{ }\DUrole{p}{|}\DUrole{w}{ }None}\DUrole{w}{ }\DUrole{o}{=}\DUrole{w}{ }\DUrole{default_value}{None}}\sphinxparamcomma \sphinxparam{\DUrole{n}{command\_timeout\_sec}\DUrole{p}{:}\DUrole{w}{ }\DUrole{n}{float}\DUrole{w}{ }\DUrole{o}{=}\DUrole{w}{ }\DUrole{default_value}{10.0}}}
{}
\pysigstopsignatures
\sphinxAtStartPar
Bases : \sphinxcode{\sphinxupquote{object}}

\sphinxAtStartPar
Configuration de capture.
\begin{description}
\sphinxlineitem{rotation :}
\sphinxAtStartPar
Rotation appliquée par l’outil de capture (0/90/180/270)

\sphinxlineitem{timeout\_ms :}
\sphinxAtStartPar
Délai avant capture (ms). Permet auto\sphinxhyphen{}exposure stable.

\sphinxlineitem{width / height :}
\sphinxAtStartPar
Optionnels. Si None, l’outil choisit par défaut.

\end{description}
\index{rotation (attribut hardware.platform.pi\_camera.CameraConfig)@\spxentry{rotation}\spxextra{attribut hardware.platform.pi\_camera.CameraConfig}}

\begin{fulllineitems}
\phantomsection\label{\detokenize{api/hardware.platform:hardware.platform.pi_camera.CameraConfig.rotation}}
\pysigstartsignatures
\pysigline
{\sphinxbfcode{\sphinxupquote{rotation}}\sphinxbfcode{\sphinxupquote{\DUrole{p}{:}\DUrole{w}{ }int}}\sphinxbfcode{\sphinxupquote{\DUrole{w}{ }\DUrole{p}{=}\DUrole{w}{ }180}}}
\pysigstopsignatures
\end{fulllineitems}

\index{timeout\_ms (attribut hardware.platform.pi\_camera.CameraConfig)@\spxentry{timeout\_ms}\spxextra{attribut hardware.platform.pi\_camera.CameraConfig}}

\begin{fulllineitems}
\phantomsection\label{\detokenize{api/hardware.platform:hardware.platform.pi_camera.CameraConfig.timeout_ms}}
\pysigstartsignatures
\pysigline
{\sphinxbfcode{\sphinxupquote{timeout\_ms}}\sphinxbfcode{\sphinxupquote{\DUrole{p}{:}\DUrole{w}{ }int}}\sphinxbfcode{\sphinxupquote{\DUrole{w}{ }\DUrole{p}{=}\DUrole{w}{ }500}}}
\pysigstopsignatures
\end{fulllineitems}

\index{width (attribut hardware.platform.pi\_camera.CameraConfig)@\spxentry{width}\spxextra{attribut hardware.platform.pi\_camera.CameraConfig}}

\begin{fulllineitems}
\phantomsection\label{\detokenize{api/hardware.platform:hardware.platform.pi_camera.CameraConfig.width}}
\pysigstartsignatures
\pysigline
{\sphinxbfcode{\sphinxupquote{width}}\sphinxbfcode{\sphinxupquote{\DUrole{p}{:}\DUrole{w}{ }int\DUrole{w}{ }\DUrole{p}{|}\DUrole{w}{ }None}}\sphinxbfcode{\sphinxupquote{\DUrole{w}{ }\DUrole{p}{=}\DUrole{w}{ }None}}}
\pysigstopsignatures
\end{fulllineitems}

\index{height (attribut hardware.platform.pi\_camera.CameraConfig)@\spxentry{height}\spxextra{attribut hardware.platform.pi\_camera.CameraConfig}}

\begin{fulllineitems}
\phantomsection\label{\detokenize{api/hardware.platform:hardware.platform.pi_camera.CameraConfig.height}}
\pysigstartsignatures
\pysigline
{\sphinxbfcode{\sphinxupquote{height}}\sphinxbfcode{\sphinxupquote{\DUrole{p}{:}\DUrole{w}{ }int\DUrole{w}{ }\DUrole{p}{|}\DUrole{w}{ }None}}\sphinxbfcode{\sphinxupquote{\DUrole{w}{ }\DUrole{p}{=}\DUrole{w}{ }None}}}
\pysigstopsignatures
\end{fulllineitems}

\index{command\_timeout\_sec (attribut hardware.platform.pi\_camera.CameraConfig)@\spxentry{command\_timeout\_sec}\spxextra{attribut hardware.platform.pi\_camera.CameraConfig}}

\begin{fulllineitems}
\phantomsection\label{\detokenize{api/hardware.platform:hardware.platform.pi_camera.CameraConfig.command_timeout_sec}}
\pysigstartsignatures
\pysigline
{\sphinxbfcode{\sphinxupquote{command\_timeout\_sec}}\sphinxbfcode{\sphinxupquote{\DUrole{p}{:}\DUrole{w}{ }float}}\sphinxbfcode{\sphinxupquote{\DUrole{w}{ }\DUrole{p}{=}\DUrole{w}{ }10.0}}}
\pysigstopsignatures
\end{fulllineitems}


\end{fulllineitems}

\index{PiCamera (classe dans hardware.platform.pi\_camera)@\spxentry{PiCamera}\spxextra{classe dans hardware.platform.pi\_camera}}

\begin{fulllineitems}
\phantomsection\label{\detokenize{api/hardware.platform:hardware.platform.pi_camera.PiCamera}}
\pysigstartsignatures
\pysiglinewithargsret
{\sphinxbfcode{\sphinxupquote{\DUrole{k}{class}\DUrole{w}{ }}}\sphinxcode{\sphinxupquote{hardware.platform.pi\_camera.}}\sphinxbfcode{\sphinxupquote{PiCamera}}}
{\sphinxparam{\DUrole{n}{config}\DUrole{p}{:}\DUrole{w}{ }\DUrole{n}{{\hyperref[\detokenize{api/hardware.platform:hardware.platform.pi_camera.CameraConfig}]{\sphinxcrossref{CameraConfig}}}\DUrole{w}{ }\DUrole{p}{|}\DUrole{w}{ }None}\DUrole{w}{ }\DUrole{o}{=}\DUrole{w}{ }\DUrole{default_value}{None}}}
{}
\pysigstopsignatures
\sphinxAtStartPar
Bases : \sphinxcode{\sphinxupquote{object}}

\sphinxAtStartPar
Caméra Raspberry Pi via rpicam\sphinxhyphen{}still (recommandé) ou libcamera\sphinxhyphen{}still (legacy).
\index{probe() (méthode statique hardware.platform.pi\_camera.PiCamera)@\spxentry{probe()}\spxextra{méthode statique hardware.platform.pi\_camera.PiCamera}}

\begin{fulllineitems}
\phantomsection\label{\detokenize{api/hardware.platform:hardware.platform.pi_camera.PiCamera.probe}}
\pysigstartsignatures
\pysiglinewithargsret
{\sphinxbfcode{\sphinxupquote{\DUrole{k}{static}\DUrole{w}{ }}}\sphinxbfcode{\sphinxupquote{probe}}}
{}
{{ $\rightarrow$ tuple\DUrole{p}{{[}}bool\DUrole{p}{,}\DUrole{w}{ }str\DUrole{p}{{]}}}}
\pysigstopsignatures
\sphinxAtStartPar
Vérifie si une caméra est détectée et retourne (available, diagnostic).

\end{fulllineitems}

\index{is\_available() (méthode statique hardware.platform.pi\_camera.PiCamera)@\spxentry{is\_available()}\spxextra{méthode statique hardware.platform.pi\_camera.PiCamera}}

\begin{fulllineitems}
\phantomsection\label{\detokenize{api/hardware.platform:hardware.platform.pi_camera.PiCamera.is_available}}
\pysigstartsignatures
\pysiglinewithargsret
{\sphinxbfcode{\sphinxupquote{\DUrole{k}{static}\DUrole{w}{ }}}\sphinxbfcode{\sphinxupquote{is\_available}}}
{}
{{ $\rightarrow$ bool}}
\pysigstopsignatures
\sphinxAtStartPar
Vérifie si l’outil libcamera/rpicam voit au moins une caméra.

\end{fulllineitems}

\index{detect\_blocking\_processes() (méthode statique hardware.platform.pi\_camera.PiCamera)@\spxentry{detect\_blocking\_processes()}\spxextra{méthode statique hardware.platform.pi\_camera.PiCamera}}

\begin{fulllineitems}
\phantomsection\label{\detokenize{api/hardware.platform:hardware.platform.pi_camera.PiCamera.detect_blocking_processes}}
\pysigstartsignatures
\pysiglinewithargsret
{\sphinxbfcode{\sphinxupquote{\DUrole{k}{static}\DUrole{w}{ }}}\sphinxbfcode{\sphinxupquote{detect\_blocking\_processes}}}
{}
{{ $\rightarrow$ List\DUrole{p}{{[}}str\DUrole{p}{{]}}}}
\pysigstopsignatures
\sphinxAtStartPar
Retourne best\sphinxhyphen{}effort la liste des processus utilisant /dev/video*.

\end{fulllineitems}

\index{capture() (méthode hardware.platform.pi\_camera.PiCamera)@\spxentry{capture()}\spxextra{méthode hardware.platform.pi\_camera.PiCamera}}

\begin{fulllineitems}
\phantomsection\label{\detokenize{api/hardware.platform:hardware.platform.pi_camera.PiCamera.capture}}
\pysigstartsignatures
\pysiglinewithargsret
{\sphinxbfcode{\sphinxupquote{capture}}}
{\sphinxparam{\DUrole{n}{output\_path}\DUrole{p}{:}\DUrole{w}{ }\DUrole{n}{str}}}
{{ $\rightarrow$ str}}
\pysigstopsignatures
\sphinxAtStartPar
Capture une image et l’écrit dans \sphinxtitleref{output\_path}.
\begin{quote}\begin{description}
\sphinxlineitem{Renvoie}
\sphinxAtStartPar
output\_path

\sphinxlineitem{Type renvoyé}
\sphinxAtStartPar
str

\sphinxlineitem{Lève}\begin{itemize}
\item {} 
\sphinxAtStartPar
\sphinxstyleliteralstrong{\sphinxupquote{RuntimeError}} \textendash{} si l’outil de capture échoue

\item {} 
\sphinxAtStartPar
\sphinxstyleliteralstrong{\sphinxupquote{FileNotFoundError}} \textendash{} si le fichier n’est pas créé

\end{itemize}

\end{description}\end{quote}

\end{fulllineitems}


\end{fulllineitems}



\paragraph{hardware.platform.rotary\_selector module}
\label{\detokenize{api/hardware.platform:module-hardware.platform.rotary_selector}}\label{\detokenize{api/hardware.platform:hardware-platform-rotary-selector-module}}\index{module@\spxentry{module}!hardware.platform.rotary\_selector@\spxentry{hardware.platform.rotary\_selector}}\index{hardware.platform.rotary\_selector@\spxentry{hardware.platform.rotary\_selector}!module@\spxentry{module}}

\subparagraph{hardware.platform.rotary\_selector}
\label{\detokenize{api/hardware.platform:hardware-platform-rotary-selector}}
\sphinxAtStartPar
Driver bas niveau pour le \sphinxstylestrong{sélecteur rotatif + bouton poussoir}.


\subparagraph{Objectif}
\label{\detokenize{api/hardware.platform:objectif}}
\sphinxAtStartPar
Ce module encapsule la logique « physique » du rotateur :
\begin{itemize}
\item {} 
\sphinxAtStartPar
lecture des signaux \sphinxstylestrong{quadrature} (canaux A/B) pour la rotation,

\item {} 
\sphinxAtStartPar
calcul du \sphinxstylestrong{sens} du geste (« sens horaire » / « sens inverse »),

\item {} 
\sphinxAtStartPar
agrégation de plusieurs impulsions en un seul \sphinxstylestrong{geste stable},

\item {} 
\sphinxAtStartPar
protection contre les \sphinxstylestrong{rebonds} et gestes trop rapides,

\item {} \begin{description}
\sphinxlineitem{gestion des \sphinxstylestrong{appuis sur le bouton} :}\begin{itemize}
\item {} 
\sphinxAtStartPar
appui court,

\item {} 
\sphinxAtStartPar
appui long,

\item {} 
\sphinxAtStartPar
double\sphinxhyphen{}clic (double appui court rapide).

\end{itemize}

\end{description}

\end{itemize}

\sphinxAtStartPar
Important : ce driver est \sphinxstylestrong{agnostique} de la logique métier.
\begin{itemize}
\item {} 
\sphinxAtStartPar
Il \sphinxstylestrong{ne connaît pas} les modes, ni le speaker, ni le multimètre.

\item {} \begin{description}
\sphinxlineitem{Il expose simplement des \sphinxstylestrong{callbacks} que le reste du système branche :}\begin{itemize}
\item {} 
\sphinxAtStartPar
\sphinxtitleref{on\_position\_changed(index, direction)}

\item {} 
\sphinxAtStartPar
\sphinxtitleref{on\_short\_press()}

\item {} 
\sphinxAtStartPar
\sphinxtitleref{on\_long\_press()}

\item {} 
\sphinxAtStartPar
\sphinxtitleref{on\_double\_press()}

\end{itemize}

\end{description}

\end{itemize}

\sphinxAtStartPar
L’objectif est d’avoir un composant réutilisable, proprement isolé, qui
s’occupe uniquement de transformer des signaux GPIO bruts en événements
haut niveau exploitables par l’assistant.
\index{RotarySelector (classe dans hardware.platform.rotary\_selector)@\spxentry{RotarySelector}\spxextra{classe dans hardware.platform.rotary\_selector}}

\begin{fulllineitems}
\phantomsection\label{\detokenize{api/hardware.platform:hardware.platform.rotary_selector.RotarySelector}}
\pysigstartsignatures
\pysiglinewithargsret
{\sphinxbfcode{\sphinxupquote{\DUrole{k}{class}\DUrole{w}{ }}}\sphinxcode{\sphinxupquote{hardware.platform.rotary\_selector.}}\sphinxbfcode{\sphinxupquote{RotarySelector}}}
{\sphinxparam{\DUrole{n}{pin\_a}\DUrole{p}{:}\DUrole{w}{ }\DUrole{n}{int}\DUrole{w}{ }\DUrole{o}{=}\DUrole{w}{ }\DUrole{default_value}{17}}\sphinxparamcomma \sphinxparam{\DUrole{n}{pin\_b}\DUrole{p}{:}\DUrole{w}{ }\DUrole{n}{int}\DUrole{w}{ }\DUrole{o}{=}\DUrole{w}{ }\DUrole{default_value}{27}}\sphinxparamcomma \sphinxparam{\DUrole{n}{pin\_sw}\DUrole{p}{:}\DUrole{w}{ }\DUrole{n}{int}\DUrole{w}{ }\DUrole{o}{=}\DUrole{w}{ }\DUrole{default_value}{22}}\sphinxparamcomma \sphinxparam{\DUrole{n}{positions\_count}\DUrole{p}{:}\DUrole{w}{ }\DUrole{n}{int}\DUrole{w}{ }\DUrole{o}{=}\DUrole{w}{ }\DUrole{default_value}{4}}\sphinxparamcomma \sphinxparam{\DUrole{n}{long\_press\_threshold}\DUrole{p}{:}\DUrole{w}{ }\DUrole{n}{float}\DUrole{w}{ }\DUrole{o}{=}\DUrole{w}{ }\DUrole{default_value}{1.0}}\sphinxparamcomma \sphinxparam{\DUrole{n}{gesture\_settle\_delay}\DUrole{p}{:}\DUrole{w}{ }\DUrole{n}{float}\DUrole{w}{ }\DUrole{o}{=}\DUrole{w}{ }\DUrole{default_value}{0.2}}\sphinxparamcomma \sphinxparam{\DUrole{n}{rotate\_button\_deadzone}\DUrole{p}{:}\DUrole{w}{ }\DUrole{n}{float}\DUrole{w}{ }\DUrole{o}{=}\DUrole{w}{ }\DUrole{default_value}{0.15}}\sphinxparamcomma \sphinxparam{\DUrole{n}{min\_step\_interval}\DUrole{p}{:}\DUrole{w}{ }\DUrole{n}{float}\DUrole{w}{ }\DUrole{o}{=}\DUrole{w}{ }\DUrole{default_value}{0.3}}\sphinxparamcomma \sphinxparam{\DUrole{n}{double\_click\_window}\DUrole{p}{:}\DUrole{w}{ }\DUrole{n}{float}\DUrole{w}{ }\DUrole{o}{=}\DUrole{w}{ }\DUrole{default_value}{0.4}}\sphinxparamcomma \sphinxparam{\DUrole{n}{button\_bounce\_time}\DUrole{p}{:}\DUrole{w}{ }\DUrole{n}{float}\DUrole{w}{ }\DUrole{o}{=}\DUrole{w}{ }\DUrole{default_value}{0.05}}}
{}
\pysigstopsignatures
\sphinxAtStartPar
Bases : \sphinxcode{\sphinxupquote{object}}

\sphinxAtStartPar
Driver pour un encodeur rotatif incrémental avec bouton poussoir intégré.


\subparagraph{Rôle}
\label{\detokenize{api/hardware.platform:role}}\begin{itemize}
\item {} 
\sphinxAtStartPar
Maintient une \sphinxstylestrong{position logique de mode} : \sphinxtitleref{0 .. positions\_count\sphinxhyphen{}1}.

\item {} 
\sphinxAtStartPar
Détecte le \sphinxstylestrong{sens de rotation} et produit des gestes stables.

\item {} \begin{description}
\sphinxlineitem{Gère les appuis :}\begin{itemize}
\item {} 
\sphinxAtStartPar
court,

\item {} 
\sphinxAtStartPar
long,

\item {} 
\sphinxAtStartPar
double appui rapide (double\sphinxhyphen{}clic).

\end{itemize}

\end{description}

\end{itemize}


\subparagraph{Stratégie UX}
\label{\detokenize{api/hardware.platform:strategie-ux}}\begin{itemize}
\item {} 
\sphinxAtStartPar
Plusieurs impulsions rapprochées dans le \sphinxstylestrong{même sens} sont agrégées
en \sphinxstylestrong{un seul geste} :
\begin{quote}

\sphinxAtStartPar
=\textgreater{} un seul changement de position / mode.
\end{quote}

\item {} 
\sphinxAtStartPar
Les allers\sphinxhyphen{}retours très rapides (bruit, gestes hésitants) sont
filtrés via :
\begin{itemize}
\item {} 
\sphinxAtStartPar
un délai minimum entre deux changements de mode
(\sphinxtitleref{min\_step\_interval}),

\item {} 
\sphinxAtStartPar
un délai de \sphinxstylestrong{stabilisation de geste} (\sphinxtitleref{gesture\_settle\_delay}).

\end{itemize}

\item {} 
\sphinxAtStartPar
Une \sphinxstylestrong{zone morte} autour des appuis bouton (\sphinxtitleref{rotate\_button\_deadzone})
évite que des rebonds mécaniques du bouton soient interprétés comme
des rotations.

\item {} 
\sphinxAtStartPar
La gestion du \sphinxstylestrong{double\sphinxhyphen{}clic} repose sur une petite fenêtre de temps
(\sphinxtitleref{double\_click\_window}) : si deux appuis courts se succèdent dans ce
délai, on considère qu’il s’agit d’un double appui.

\end{itemize}


\subparagraph{Paramètres}
\label{\detokenize{api/hardware.platform:parametres}}\begin{description}
\sphinxlineitem{pin\_a}{[}int{]}
\sphinxAtStartPar
Broche GPIO du canal A de l’encodeur.

\sphinxlineitem{pin\_b}{[}int{]}
\sphinxAtStartPar
Broche GPIO du canal B de l’encodeur.

\sphinxlineitem{pin\_sw}{[}int{]}
\sphinxAtStartPar
Broche GPIO du bouton poussoir.

\sphinxlineitem{positions\_count}{[}int{]}
\sphinxAtStartPar
Nombre total de positions logiques (nombre de modes disponibles).

\sphinxlineitem{long\_press\_threshold}{[}float{]}
\sphinxAtStartPar
Durée (en secondes) à partir de laquelle un appui est considéré comme
« long » plutôt que « court ».

\sphinxlineitem{gesture\_settle\_delay}{[}float{]}
\sphinxAtStartPar
Temps sans nouvelle impulsion avant de considérer qu’un geste de
rotation est terminé et de changer effectivement de position.

\sphinxlineitem{rotate\_button\_deadzone}{[}float{]}
\sphinxAtStartPar
Intervalle de temps (en secondes) autour d’un événement bouton
pendant lequel on ignore les impulsions de rotation (anti\sphinxhyphen{}parasites).

\sphinxlineitem{min\_step\_interval}{[}float{]}
\sphinxAtStartPar
Délai minimum (en secondes) entre deux changements de position
consécutifs (limite les changements trop rapides).

\sphinxlineitem{double\_click\_window}{[}float{]}
\sphinxAtStartPar
Temps maximum (en secondes) entre deux appuis courts consécutifs pour
être interprétés comme un \sphinxstylestrong{double appui}.

\end{description}

\end{fulllineitems}



\paragraph{Module contents}
\label{\detokenize{api/hardware.platform:module-hardware.platform}}\label{\detokenize{api/hardware.platform:module-contents}}\index{module@\spxentry{module}!hardware.platform@\spxentry{hardware.platform}}\index{hardware.platform@\spxentry{hardware.platform}!module@\spxentry{module}}
\sphinxstepscope


\subsubsection{hardware.transports package}
\label{\detokenize{api/hardware.transports:hardware-transports-package}}\label{\detokenize{api/hardware.transports::doc}}

\paragraph{Submodules}
\label{\detokenize{api/hardware.transports:submodules}}

\paragraph{hardware.transports.ble\_client module}
\label{\detokenize{api/hardware.transports:module-hardware.transports.ble_client}}\label{\detokenize{api/hardware.transports:hardware-transports-ble-client-module}}\index{module@\spxentry{module}!hardware.transports.ble\_client@\spxentry{hardware.transports.ble\_client}}\index{hardware.transports.ble\_client@\spxentry{hardware.transports.ble\_client}!module@\spxentry{module}}

\subparagraph{hardware.transports.ble\_client}
\label{\detokenize{api/hardware.transports:hardware-transports-ble-client}}
\sphinxAtStartPar
Client BLE générique basé sur \sphinxstylestrong{bleak}, utilisé comme couche de transport.


\subparagraph{Rôle du module}
\label{\detokenize{api/hardware.transports:role-du-module}}
\sphinxAtStartPar
Ce module fournit un petit \sphinxstylestrong{service BLE réutilisable} qui :
\begin{itemize}
\item {} 
\sphinxAtStartPar
scanne l’environnement pour trouver un périphérique BLE cible,

\item {} 
\sphinxAtStartPar
se connecte à ce périphérique,

\item {} 
\sphinxAtStartPar
s’abonne à une \sphinxstylestrong{caractéristique} pour recevoir des notifications,

\item {} 
\sphinxAtStartPar
gère les \sphinxstylestrong{reconnexions automatiques} en cas de perte de connexion
ou d’erreur,

\item {} \begin{description}
\sphinxlineitem{expose deux callbacks indépendants de la logique métier :}\begin{itemize}
\item {} \begin{description}
\sphinxlineitem{\sphinxtitleref{on\_status(event: str)}  :}\begin{itemize}
\item {} 
\sphinxAtStartPar
informe l’appelant de l’état du client (recherche, trouvé,
connecté, déconnecté, erreurs, etc.) ;

\end{itemize}

\end{description}

\item {} \begin{description}
\sphinxlineitem{\sphinxtitleref{on\_packet(data: bytes)} :}\begin{itemize}
\item {} 
\sphinxAtStartPar
transmet les \sphinxstylestrong{trames brutes} reçues sur la caractéristique.

\end{itemize}

\end{description}

\end{itemize}

\end{description}

\end{itemize}

\begin{sphinxadmonition}{important}{Important:}
\sphinxAtStartPar
Ce client \sphinxstylestrong{ne connaît pas} la notion de « multimètre », « capteur X », etc.
\begin{itemize}
\item {} 
\sphinxAtStartPar
Il ne fait que gérer le \sphinxstylestrong{transport BLE} (scan, connexion, notifications).

\item {} 
\sphinxAtStartPar
La logique métier (décodage des trames, phrases vocales, etc.)
est déportée dans des drivers plus haut niveau (ex : \sphinxtitleref{OwonMultimeterDriver}).

\end{itemize}

\sphinxAtStartPar
Il est donc possible de réutiliser cette classe pour \sphinxstylestrong{n’importe quel
périphérique BLE}, à condition de fournir :
\sphinxhyphen{} un filtre de nom ou d’adresse,
\sphinxhyphen{} l’UUID de la caractéristique à écouter,
\sphinxhyphen{} des callbacks adaptés.
\end{sphinxadmonition}
\index{BleClient (classe dans hardware.transports.ble\_client)@\spxentry{BleClient}\spxextra{classe dans hardware.transports.ble\_client}}

\begin{fulllineitems}
\phantomsection\label{\detokenize{api/hardware.transports:hardware.transports.ble_client.BleClient}}
\pysigstartsignatures
\pysiglinewithargsret
{\sphinxbfcode{\sphinxupquote{\DUrole{k}{class}\DUrole{w}{ }}}\sphinxcode{\sphinxupquote{hardware.transports.ble\_client.}}\sphinxbfcode{\sphinxupquote{BleClient}}}
{\sphinxparam{\DUrole{keyword-only-separator}{\DUrole{o}{\sphinxstyleabbreviation{*}}}}\sphinxparamcomma \sphinxparam{\DUrole{n}{name\_filter}\DUrole{p}{:}\DUrole{w}{ }\DUrole{n}{str\DUrole{w}{ }\DUrole{p}{|}\DUrole{w}{ }None}}\sphinxparamcomma \sphinxparam{\DUrole{n}{characteristic\_uuid}\DUrole{p}{:}\DUrole{w}{ }\DUrole{n}{str}}\sphinxparamcomma \sphinxparam{\DUrole{n}{address\_filter}\DUrole{p}{:}\DUrole{w}{ }\DUrole{n}{str\DUrole{w}{ }\DUrole{p}{|}\DUrole{w}{ }None}\DUrole{w}{ }\DUrole{o}{=}\DUrole{w}{ }\DUrole{default_value}{None}}\sphinxparamcomma \sphinxparam{\DUrole{n}{on\_status}\DUrole{p}{:}\DUrole{w}{ }\DUrole{n}{Callable\DUrole{p}{{[}}\DUrole{p}{{[}}str\DUrole{p}{{]}}\DUrole{p}{,}\DUrole{w}{ }None\DUrole{p}{{]}}\DUrole{w}{ }\DUrole{p}{|}\DUrole{w}{ }None}\DUrole{w}{ }\DUrole{o}{=}\DUrole{w}{ }\DUrole{default_value}{None}}\sphinxparamcomma \sphinxparam{\DUrole{n}{on\_packet}\DUrole{p}{:}\DUrole{w}{ }\DUrole{n}{Callable\DUrole{p}{{[}}\DUrole{p}{{[}}bytes\DUrole{p}{{]}}\DUrole{p}{,}\DUrole{w}{ }None\DUrole{p}{{]}}\DUrole{w}{ }\DUrole{p}{|}\DUrole{w}{ }None}\DUrole{w}{ }\DUrole{o}{=}\DUrole{w}{ }\DUrole{default_value}{None}}\sphinxparamcomma \sphinxparam{\DUrole{n}{scan\_timeout}\DUrole{p}{:}\DUrole{w}{ }\DUrole{n}{float}\DUrole{w}{ }\DUrole{o}{=}\DUrole{w}{ }\DUrole{default_value}{4.0}}\sphinxparamcomma \sphinxparam{\DUrole{n}{retry\_delay}\DUrole{p}{:}\DUrole{w}{ }\DUrole{n}{float}\DUrole{w}{ }\DUrole{o}{=}\DUrole{w}{ }\DUrole{default_value}{15.0}}\sphinxparamcomma \sphinxparam{\DUrole{n}{reconnect\_delay}\DUrole{p}{:}\DUrole{w}{ }\DUrole{n}{float}\DUrole{w}{ }\DUrole{o}{=}\DUrole{w}{ }\DUrole{default_value}{2.0}}\sphinxparamcomma \sphinxparam{\DUrole{n}{listen\_poll\_interval\_sec}\DUrole{p}{:}\DUrole{w}{ }\DUrole{n}{float}\DUrole{w}{ }\DUrole{o}{=}\DUrole{w}{ }\DUrole{default_value}{0.5}}}
{}
\pysigstopsignatures
\sphinxAtStartPar
Bases : \sphinxcode{\sphinxupquote{object}}

\sphinxAtStartPar
Client BLE générique avec boucle de reconnexion automatique.

\sphinxAtStartPar
Ce composant tourne dans un \sphinxstylestrong{thread dédié} pour ne pas bloquer
le thread principal de l’application. La logique interne utilise
\sphinxtitleref{asyncio} (scan / connexion / notifications), mais l’API publique
(\sphinxtitleref{start()}, \sphinxtitleref{stop()}, propriétés) reste \sphinxstylestrong{synchrone} et simple à utiliser.


\subparagraph{Paramètres}
\label{\detokenize{api/hardware.transports:parametres}}\begin{description}
\sphinxlineitem{name\_filter}{[}Optional{[}str{]}{]}
\sphinxAtStartPar
Nom BLE à rechercher (exact, insensible à la casse).
Exemple : \sphinxtitleref{« BDM »} pour le multimètre OWON 16.

\sphinxlineitem{address\_filter}{[}Optional{[}str{]}{]}
\sphinxAtStartPar
Adresse MAC BLE à rechercher (si fournie, elle est prioritaire
sur \sphinxtitleref{name\_filter}).

\sphinxlineitem{characteristic\_uuid}{[}str{]}
\sphinxAtStartPar
UUID de la caractéristique sur laquelle on reçoit les notifications.
Exemple : \sphinxtitleref{« 0000fff4\sphinxhyphen{}0000\sphinxhyphen{}1000\sphinxhyphen{}8000\sphinxhyphen{}00805f9b34fb »}.

\sphinxlineitem{on\_status}{[}Optional{[}Callable{[}{[}str{]}, None{]}{]}{]}
\sphinxAtStartPar
Callback pour les événements d’état. Les valeurs possibles sont :
\sphinxhyphen{} \sphinxtitleref{« search\_start »}  : début d’une phase de recherche.
\sphinxhyphen{} \sphinxtitleref{« not\_found »}     : aucun device cible trouvé pour cette tentative.
\sphinxhyphen{} \sphinxtitleref{« found »}         : device cible détecté (avant connexion).
\sphinxhyphen{} \sphinxtitleref{« connecting »}    : tentative de connexion en cours.
\sphinxhyphen{} \sphinxtitleref{« connected »}     : connexion établie.
\sphinxhyphen{} \sphinxtitleref{« disconnected »}  : connexion perdue alors que le client tournait.
\sphinxhyphen{} \sphinxtitleref{« error »}         : erreur générique (scan / boucle principale).
\sphinxhyphen{} \sphinxtitleref{« connect\_error »} : erreur lors de la connexion ou de l’écoute.
\sphinxhyphen{} \sphinxtitleref{« stopped »}       : la boucle principale s’est arrêtée (après stop()).

\sphinxlineitem{on\_packet}{[}Optional{[}Callable{[}{[}bytes{]}, None{]}{]}{]}
\sphinxAtStartPar
Callback appelé à chaque paquet reçu sur la caractéristique.
Le paramètre est la trame brute \sphinxtitleref{bytes} (6 octets pour l’OWON, par ex.).

\sphinxlineitem{scan\_timeout}{[}float{]}
\sphinxAtStartPar
Durée (en secondes) d’une phase de scan BLE (\sphinxtitleref{BleakScanner.discover}).

\sphinxlineitem{retry\_delay}{[}float{]}
\sphinxAtStartPar
Délai (en secondes) entre deux tentatives de scan lorsqu’aucun device
n’a été trouvé.

\sphinxlineitem{reconnect\_delay}{[}float{]}
\sphinxAtStartPar
Délai (en secondes) avant de retenter la connexion après une erreur.

\end{description}
\index{is\_running (propriété hardware.transports.ble\_client.BleClient)@\spxentry{is\_running}\spxextra{propriété hardware.transports.ble\_client.BleClient}}

\begin{fulllineitems}
\phantomsection\label{\detokenize{api/hardware.transports:hardware.transports.ble_client.BleClient.is_running}}
\pysigstartsignatures
\pysigline
{\sphinxbfcode{\sphinxupquote{\DUrole{k}{property}\DUrole{w}{ }}}\sphinxbfcode{\sphinxupquote{is\_running}}\sphinxbfcode{\sphinxupquote{\DUrole{p}{:}\DUrole{w}{ }bool}}}
\pysigstopsignatures
\sphinxAtStartPar
Indique si la boucle de monitoring BLE est actuellement active.
\begin{quote}\begin{description}
\sphinxlineitem{Renvoie}
\sphinxAtStartPar
\sphinxtitleref{True} si le thread BLE est démarré (\sphinxtitleref{start()} appelé et non stoppé),
\sphinxtitleref{False} sinon.

\sphinxlineitem{Type renvoyé}
\sphinxAtStartPar
bool

\end{description}\end{quote}

\end{fulllineitems}

\index{is\_connected (propriété hardware.transports.ble\_client.BleClient)@\spxentry{is\_connected}\spxextra{propriété hardware.transports.ble\_client.BleClient}}

\begin{fulllineitems}
\phantomsection\label{\detokenize{api/hardware.transports:hardware.transports.ble_client.BleClient.is_connected}}
\pysigstartsignatures
\pysigline
{\sphinxbfcode{\sphinxupquote{\DUrole{k}{property}\DUrole{w}{ }}}\sphinxbfcode{\sphinxupquote{is\_connected}}\sphinxbfcode{\sphinxupquote{\DUrole{p}{:}\DUrole{w}{ }bool}}}
\pysigstopsignatures
\sphinxAtStartPar
Indique si le client est \sphinxstylestrong{actuellement connecté} à un périphérique BLE.
\begin{quote}\begin{description}
\sphinxlineitem{Renvoie}
\sphinxAtStartPar
\sphinxtitleref{True} si une connexion est active et que les notifications
sont en cours, \sphinxtitleref{False} sinon.

\sphinxlineitem{Type renvoyé}
\sphinxAtStartPar
bool

\end{description}\end{quote}

\end{fulllineitems}

\index{start() (méthode hardware.transports.ble\_client.BleClient)@\spxentry{start()}\spxextra{méthode hardware.transports.ble\_client.BleClient}}

\begin{fulllineitems}
\phantomsection\label{\detokenize{api/hardware.transports:hardware.transports.ble_client.BleClient.start}}
\pysigstartsignatures
\pysiglinewithargsret
{\sphinxbfcode{\sphinxupquote{start}}}
{}
{{ $\rightarrow$ None}}
\pysigstopsignatures
\sphinxAtStartPar
Démarre la boucle de monitoring dans un \sphinxstylestrong{thread séparé}.


\subparagraph{Comportement}
\label{\detokenize{api/hardware.transports:comportement}}\begin{itemize}
\item {} 
\sphinxAtStartPar
Si le client est déjà en cours d’exécution (\sphinxtitleref{is\_running == True}),
la méthode ne fait rien (idempotent).

\item {} \begin{description}
\sphinxlineitem{Sinon :}\begin{itemize}
\item {} 
\sphinxAtStartPar
crée un nouveau thread (\sphinxtitleref{BleClientThread}),

\item {} 
\sphinxAtStartPar
lance \sphinxtitleref{\_run\_asyncio\_loop()} dans ce thread,

\item {} 
\sphinxAtStartPar
active la reconnexion automatique tant que \sphinxtitleref{stop()} n’est pas appelé.

\end{itemize}

\end{description}

\end{itemize}

\end{fulllineitems}

\index{stop() (méthode hardware.transports.ble\_client.BleClient)@\spxentry{stop()}\spxextra{méthode hardware.transports.ble\_client.BleClient}}

\begin{fulllineitems}
\phantomsection\label{\detokenize{api/hardware.transports:hardware.transports.ble_client.BleClient.stop}}
\pysigstartsignatures
\pysiglinewithargsret
{\sphinxbfcode{\sphinxupquote{stop}}}
{}
{{ $\rightarrow$ None}}
\pysigstopsignatures
\sphinxAtStartPar
Demande l’arrêt de la boucle de monitoring BLE.


\subparagraph{Effets}
\label{\detokenize{api/hardware.transports:effets}}\begin{itemize}
\item {} 
\sphinxAtStartPar
Met \sphinxtitleref{\_running} à \sphinxtitleref{False} : la boucle principale (\sphinxtitleref{\_main\_loop}) va
se terminer proprement à la prochaine itération.

\item {} 
\sphinxAtStartPar
À la fin, un événement \sphinxtitleref{« stopped »} sera émis via \sphinxtitleref{on\_status}, et
\sphinxtitleref{is\_connected} passera à \sphinxtitleref{False}.

\end{itemize}

\end{fulllineitems}


\end{fulllineitems}



\paragraph{hardware.transports.nrf24\_transport module}
\label{\detokenize{api/hardware.transports:module-hardware.transports.nrf24_transport}}\label{\detokenize{api/hardware.transports:hardware-transports-nrf24-transport-module}}\index{module@\spxentry{module}!hardware.transports.nrf24\_transport@\spxentry{hardware.transports.nrf24\_transport}}\index{hardware.transports.nrf24\_transport@\spxentry{hardware.transports.nrf24\_transport}!module@\spxentry{module}}

\subparagraph{hardware.transports.nrf24\_transport}
\label{\detokenize{api/hardware.transports:hardware-transports-nrf24-transport}}
\sphinxAtStartPar
Ce module définit la classe \sphinxtitleref{Nrf24Transport}, responsable de la communication
bas niveau avec le module radio NRF24L01+.


\subparagraph{Vue d’ensemble}
\label{\detokenize{api/hardware.transports:vue-d-ensemble}}
\sphinxAtStartPar
Le transport NRF24 agit comme une couche d’abstraction au\sphinxhyphen{}dessus du bus SPI :
\sphinxhyphen{} Il initialise le matériel radio (CE, CSN, Canal, Débit).
\sphinxhyphen{} Il gère la réception des données en arrière\sphinxhyphen{}plan via un \sphinxstylestrong{thread dédié}.
\sphinxhyphen{} Il propage les paquets bruts reçus vers les drivers de périphériques
\begin{quote}

\sphinxAtStartPar
(ex. : pied à coulisse) via un système de callback.
\end{quote}


\subparagraph{Objectifs du design}
\label{\detokenize{api/hardware.transports:objectifs-du-design}}\begin{itemize}
\item {} 
\sphinxAtStartPar
\sphinxstylestrong{Robustesse} : Utilise des blocs try/except pour éviter le crash de l’application
si la librairie \sphinxtitleref{pyrf24} est absente ou si le module radio est débranché.

\item {} 
\sphinxAtStartPar
\sphinxstylestrong{Réactivité} : Le thread \sphinxtitleref{Nrf24Worker} écoute en permanence le bus SPI sans
bloquer la boucle principale du \sphinxtitleref{ModeManager}.

\item {} 
\sphinxAtStartPar
\sphinxstylestrong{Découplage} : Le transport ne connaît pas le contenu des messages ; il se
contente de livrer des octets bruts (\sphinxtitleref{bytes}).

\end{itemize}


\subparagraph{Câblage type (Raspberry Pi 5)}
\label{\detokenize{api/hardware.transports:cablage-type-raspberry-pi-5}}
\sphinxAtStartPar
VCC \sphinxhyphen{}\textgreater{} 3.3V | GND \sphinxhyphen{}\textgreater{} GND | CE \sphinxhyphen{}\textgreater{} GPIO 25 | CSN \sphinxhyphen{}\textgreater{} GPIO 8 (CE0)
SCK \sphinxhyphen{}\textgreater{} GPIO 11 | MOSI \sphinxhyphen{}\textgreater{} GPIO 10 | MISO \sphinxhyphen{}\textgreater{} GPIO 9
\index{Nrf24Transport (classe dans hardware.transports.nrf24\_transport)@\spxentry{Nrf24Transport}\spxextra{classe dans hardware.transports.nrf24\_transport}}

\begin{fulllineitems}
\phantomsection\label{\detokenize{api/hardware.transports:hardware.transports.nrf24_transport.Nrf24Transport}}
\pysigstartsignatures
\pysiglinewithargsret
{\sphinxbfcode{\sphinxupquote{\DUrole{k}{class}\DUrole{w}{ }}}\sphinxcode{\sphinxupquote{hardware.transports.nrf24\_transport.}}\sphinxbfcode{\sphinxupquote{Nrf24Transport}}}
{\sphinxparam{\DUrole{n}{ce\_pin}\DUrole{p}{:}\DUrole{w}{ }\DUrole{n}{int}\DUrole{w}{ }\DUrole{o}{=}\DUrole{w}{ }\DUrole{default_value}{25}}\sphinxparamcomma \sphinxparam{\DUrole{n}{csn\_pin}\DUrole{p}{:}\DUrole{w}{ }\DUrole{n}{int}\DUrole{w}{ }\DUrole{o}{=}\DUrole{w}{ }\DUrole{default_value}{0}}\sphinxparamcomma \sphinxparam{\DUrole{n}{channel}\DUrole{p}{:}\DUrole{w}{ }\DUrole{n}{int}\DUrole{w}{ }\DUrole{o}{=}\DUrole{w}{ }\DUrole{default_value}{120}}\sphinxparamcomma \sphinxparam{\DUrole{n}{address}\DUrole{p}{:}\DUrole{w}{ }\DUrole{n}{bytes}\DUrole{w}{ }\DUrole{o}{=}\DUrole{w}{ }\DUrole{default_value}{b\textquotesingle{}00001\textquotesingle{}}}}
{}
\pysigstopsignatures
\sphinxAtStartPar
Bases : \sphinxcode{\sphinxupquote{object}}

\sphinxAtStartPar
Gère la couche basse radio NRF24L01+ sur bus SPI.


\subparagraph{Rôle}
\label{\detokenize{api/hardware.transports:role}}
\sphinxAtStartPar
Cette classe encapsule la librairie \sphinxtitleref{pyrf24} pour fournir une interface
de lecture simplifiée et threadée. Elle est conçue pour être utilisée
par un driver de périphérique (comme le Pied à coulisse).
\index{\_\_init\_\_() (méthode hardware.transports.nrf24\_transport.Nrf24Transport)@\spxentry{\_\_init\_\_()}\spxextra{méthode hardware.transports.nrf24\_transport.Nrf24Transport}}

\begin{fulllineitems}
\phantomsection\label{\detokenize{api/hardware.transports:hardware.transports.nrf24_transport.Nrf24Transport.__init__}}
\pysigstartsignatures
\pysiglinewithargsret
{\sphinxbfcode{\sphinxupquote{\_\_init\_\_}}}
{\sphinxparam{\DUrole{n}{ce\_pin}\DUrole{p}{:}\DUrole{w}{ }\DUrole{n}{int}\DUrole{w}{ }\DUrole{o}{=}\DUrole{w}{ }\DUrole{default_value}{25}}\sphinxparamcomma \sphinxparam{\DUrole{n}{csn\_pin}\DUrole{p}{:}\DUrole{w}{ }\DUrole{n}{int}\DUrole{w}{ }\DUrole{o}{=}\DUrole{w}{ }\DUrole{default_value}{0}}\sphinxparamcomma \sphinxparam{\DUrole{n}{channel}\DUrole{p}{:}\DUrole{w}{ }\DUrole{n}{int}\DUrole{w}{ }\DUrole{o}{=}\DUrole{w}{ }\DUrole{default_value}{120}}\sphinxparamcomma \sphinxparam{\DUrole{n}{address}\DUrole{p}{:}\DUrole{w}{ }\DUrole{n}{bytes}\DUrole{w}{ }\DUrole{o}{=}\DUrole{w}{ }\DUrole{default_value}{b\textquotesingle{}00001\textquotesingle{}}}}
{}
\pysigstopsignatures
\sphinxAtStartPar
Initialise les paramètres du transport radio.


\subparagraph{Paramètres}
\label{\detokenize{api/hardware.transports:id1}}\begin{description}
\sphinxlineitem{ce\_pin}{[}int{]}
\sphinxAtStartPar
Broche GPIO pour le signal Chip Enable (défaut : 25).

\sphinxlineitem{csn\_pin}{[}int{]}
\sphinxAtStartPar
Broche pour le signal Chip Select Not (0 pour SPI0 CE0 / GPIO 8).

\sphinxlineitem{channel}{[}int{]}
\sphinxAtStartPar
Canal radio utilisé (0 à 125). Le pied à coulisse MHK utilise le 120.

\sphinxlineitem{address}{[}bytes{]}
\sphinxAtStartPar
Adresse unique du tuyau de lecture (pipe), format 5 octets.

\end{description}

\end{fulllineitems}

\index{start() (méthode hardware.transports.nrf24\_transport.Nrf24Transport)@\spxentry{start()}\spxextra{méthode hardware.transports.nrf24\_transport.Nrf24Transport}}

\begin{fulllineitems}
\phantomsection\label{\detokenize{api/hardware.transports:hardware.transports.nrf24_transport.Nrf24Transport.start}}
\pysigstartsignatures
\pysiglinewithargsret
{\sphinxbfcode{\sphinxupquote{start}}}
{\sphinxparam{\DUrole{n}{on\_packet\_callback}\DUrole{p}{:}\DUrole{w}{ }\DUrole{n}{Callable\DUrole{p}{{[}}\DUrole{p}{{[}}bytes\DUrole{p}{{]}}\DUrole{p}{,}\DUrole{w}{ }None\DUrole{p}{{]}}}}}
{{ $\rightarrow$ Tuple\DUrole{p}{{[}}bool\DUrole{p}{,}\DUrole{w}{ }str\DUrole{p}{{]}}}}
\pysigstopsignatures
\sphinxAtStartPar
Initialise le module radio et démarre le thread d’écoute.


\subparagraph{Paramètres}
\label{\detokenize{api/hardware.transports:id2}}\begin{description}
\sphinxlineitem{on\_packet\_callback}{[}Callable{]}
\sphinxAtStartPar
Fonction appelée à chaque réception d’un nouveau paquet radio.
Prend un argument de type \sphinxtitleref{bytes}.

\end{description}
\begin{quote}\begin{description}
\sphinxlineitem{returns}
\sphinxAtStartPar
Un tuple contenant le succès (True/False) et un message de statut.

\sphinxlineitem{rtype}
\sphinxAtStartPar
(bool, str)

\end{description}\end{quote}

\end{fulllineitems}

\index{stop() (méthode hardware.transports.nrf24\_transport.Nrf24Transport)@\spxentry{stop()}\spxextra{méthode hardware.transports.nrf24\_transport.Nrf24Transport}}

\begin{fulllineitems}
\phantomsection\label{\detokenize{api/hardware.transports:hardware.transports.nrf24_transport.Nrf24Transport.stop}}
\pysigstartsignatures
\pysiglinewithargsret
{\sphinxbfcode{\sphinxupquote{stop}}}
{}
{{ $\rightarrow$ None}}
\pysigstopsignatures
\sphinxAtStartPar
Arrête proprement l’écoute radio et le thread associé.

\end{fulllineitems}


\end{fulllineitems}



\paragraph{hardware.transports.wifi\_client module}
\label{\detokenize{api/hardware.transports:module-hardware.transports.wifi_client}}\label{\detokenize{api/hardware.transports:hardware-transports-wifi-client-module}}\index{module@\spxentry{module}!hardware.transports.wifi\_client@\spxentry{hardware.transports.wifi\_client}}\index{hardware.transports.wifi\_client@\spxentry{hardware.transports.wifi\_client}!module@\spxentry{module}}

\subparagraph{hardware.transports.wifi\_client}
\label{\detokenize{api/hardware.transports:hardware-transports-wifi-client}}
\sphinxAtStartPar
Client HTTP générique pour le polling de capteurs sur réseau local.
Architecture inspirée de \sphinxtitleref{ble\_client.py}.


\subparagraph{Rôle du module}
\label{\detokenize{api/hardware.transports:id3}}
\sphinxAtStartPar
Ce module fournit un \sphinxstylestrong{service de transport HTTP réutilisable} qui :
\begin{itemize}
\item {} 
\sphinxAtStartPar
Tente de contacter une API HTTP cible (URL) à intervalles réguliers (polling).

\item {} 
\sphinxAtStartPar
Détecte automatiquement la disponibilité du périphérique (connexion/déconnexion).

\item {} 
\sphinxAtStartPar
Gère les timeouts et erreurs réseaux de manière transparente.

\item {} \begin{description}
\sphinxlineitem{Expose deux callbacks indépendants de la logique métier :}\begin{itemize}
\item {} 
\sphinxAtStartPar
\sphinxtitleref{on\_status(event: str)} : informe de l’état (connecting, connected, error…).

\item {} 
\sphinxAtStartPar
\sphinxtitleref{on\_data(data: str)}    : transmet le corps de la réponse HTTP brute.

\end{itemize}

\end{description}

\end{itemize}

\begin{sphinxadmonition}{important}{Important:}
\sphinxAtStartPar
Ce client ne gère \sphinxstylestrong{aucun protocole spécifique} au\sphinxhyphen{}delà du HTTP standard.
Il ne fait que transporter des chaînes de caractères reçues via HTTP.

\sphinxAtStartPar
Il est réutilisable pour tout capteur exposant une API REST simple,
à condition de fournir l’URL cible.
\end{sphinxadmonition}
\index{WifiClient (classe dans hardware.transports.wifi\_client)@\spxentry{WifiClient}\spxextra{classe dans hardware.transports.wifi\_client}}

\begin{fulllineitems}
\phantomsection\label{\detokenize{api/hardware.transports:hardware.transports.wifi_client.WifiClient}}
\pysigstartsignatures
\pysiglinewithargsret
{\sphinxbfcode{\sphinxupquote{\DUrole{k}{class}\DUrole{w}{ }}}\sphinxcode{\sphinxupquote{hardware.transports.wifi\_client.}}\sphinxbfcode{\sphinxupquote{WifiClient}}}
{\sphinxparam{\DUrole{keyword-only-separator}{\DUrole{o}{\sphinxstyleabbreviation{*}}}}\sphinxparamcomma \sphinxparam{\DUrole{n}{target\_url}\DUrole{p}{:}\DUrole{w}{ }\DUrole{n}{str}}\sphinxparamcomma \sphinxparam{\DUrole{n}{on\_status}\DUrole{p}{:}\DUrole{w}{ }\DUrole{n}{Callable\DUrole{p}{{[}}\DUrole{p}{{[}}str\DUrole{p}{{]}}\DUrole{p}{,}\DUrole{w}{ }None\DUrole{p}{{]}}\DUrole{w}{ }\DUrole{p}{|}\DUrole{w}{ }None}\DUrole{w}{ }\DUrole{o}{=}\DUrole{w}{ }\DUrole{default_value}{None}}\sphinxparamcomma \sphinxparam{\DUrole{n}{on\_data}\DUrole{p}{:}\DUrole{w}{ }\DUrole{n}{Callable\DUrole{p}{{[}}\DUrole{p}{{[}}str\DUrole{p}{{]}}\DUrole{p}{,}\DUrole{w}{ }None\DUrole{p}{{]}}\DUrole{w}{ }\DUrole{p}{|}\DUrole{w}{ }None}\DUrole{w}{ }\DUrole{o}{=}\DUrole{w}{ }\DUrole{default_value}{None}}\sphinxparamcomma \sphinxparam{\DUrole{n}{poll\_interval}\DUrole{p}{:}\DUrole{w}{ }\DUrole{n}{float}\DUrole{w}{ }\DUrole{o}{=}\DUrole{w}{ }\DUrole{default_value}{2.0}}\sphinxparamcomma \sphinxparam{\DUrole{n}{request\_timeout}\DUrole{p}{:}\DUrole{w}{ }\DUrole{n}{float}\DUrole{w}{ }\DUrole{o}{=}\DUrole{w}{ }\DUrole{default_value}{3.0}}}
{}
\pysigstopsignatures
\sphinxAtStartPar
Bases : \sphinxcode{\sphinxupquote{object}}

\sphinxAtStartPar
Client de Polling HTTP générique.

\sphinxAtStartPar
Ce composant tourne dans un \sphinxstylestrong{thread dédié} pour ne pas bloquer le thread principal.


\subparagraph{Paramètres}
\label{\detokenize{api/hardware.transports:id4}}\begin{description}
\sphinxlineitem{target\_url}{[}str{]}
\sphinxAtStartPar
L’URL complète de l’API à interroger (ex: « \sphinxurl{http://thermo.local/api} »).

\sphinxlineitem{on\_status}{[}Optional{[}Callable{[}{[}str{]}, None{]}{]}{]}
\sphinxAtStartPar
Callback pour les événements d’état. Valeurs possibles :
\sphinxhyphen{} « connecting »     : tentative de contact avec l’URL (début ou reconnexion).
\sphinxhyphen{} « connected »      : périphérique joint avec succès (HTTP 200).
\sphinxhyphen{} « disconnected »   : perte de liaison (timeout ou erreur réseau).
\sphinxhyphen{} « http\_error »     : le périphérique répond mais avec une erreur (404, 500…).
\sphinxhyphen{} « stopped »        : arrêt volontaire du client.

\sphinxlineitem{on\_data}{[}Optional{[}Callable{[}{[}str{]}, None{]}{]}{]}
\sphinxAtStartPar
Callback appelé à chaque réponse HTTP valide.
Le paramètre est le texte brut (str) de la réponse.

\sphinxlineitem{poll\_interval}{[}float{]}
\sphinxAtStartPar
Temps (en secondes) entre deux requêtes HTTP.

\sphinxlineitem{request\_timeout}{[}float{]}
\sphinxAtStartPar
Temps (en secondes) avant de considérer le périphérique comme absent.

\end{description}
\index{is\_running (propriété hardware.transports.wifi\_client.WifiClient)@\spxentry{is\_running}\spxextra{propriété hardware.transports.wifi\_client.WifiClient}}

\begin{fulllineitems}
\phantomsection\label{\detokenize{api/hardware.transports:hardware.transports.wifi_client.WifiClient.is_running}}
\pysigstartsignatures
\pysigline
{\sphinxbfcode{\sphinxupquote{\DUrole{k}{property}\DUrole{w}{ }}}\sphinxbfcode{\sphinxupquote{is\_running}}\sphinxbfcode{\sphinxupquote{\DUrole{p}{:}\DUrole{w}{ }bool}}}
\pysigstopsignatures
\sphinxAtStartPar
Indique si la boucle de polling est active.

\end{fulllineitems}

\index{is\_connected (propriété hardware.transports.wifi\_client.WifiClient)@\spxentry{is\_connected}\spxextra{propriété hardware.transports.wifi\_client.WifiClient}}

\begin{fulllineitems}
\phantomsection\label{\detokenize{api/hardware.transports:hardware.transports.wifi_client.WifiClient.is_connected}}
\pysigstartsignatures
\pysigline
{\sphinxbfcode{\sphinxupquote{\DUrole{k}{property}\DUrole{w}{ }}}\sphinxbfcode{\sphinxupquote{is\_connected}}\sphinxbfcode{\sphinxupquote{\DUrole{p}{:}\DUrole{w}{ }bool}}}
\pysigstopsignatures
\sphinxAtStartPar
Indique si la dernière communication HTTP a réussi.

\end{fulllineitems}

\index{start() (méthode hardware.transports.wifi\_client.WifiClient)@\spxentry{start()}\spxextra{méthode hardware.transports.wifi\_client.WifiClient}}

\begin{fulllineitems}
\phantomsection\label{\detokenize{api/hardware.transports:hardware.transports.wifi_client.WifiClient.start}}
\pysigstartsignatures
\pysiglinewithargsret
{\sphinxbfcode{\sphinxupquote{start}}}
{}
{{ $\rightarrow$ None}}
\pysigstopsignatures
\sphinxAtStartPar
Démarre la boucle de polling dans un thread séparé.

\end{fulllineitems}

\index{stop() (méthode hardware.transports.wifi\_client.WifiClient)@\spxentry{stop()}\spxextra{méthode hardware.transports.wifi\_client.WifiClient}}

\begin{fulllineitems}
\phantomsection\label{\detokenize{api/hardware.transports:hardware.transports.wifi_client.WifiClient.stop}}
\pysigstartsignatures
\pysiglinewithargsret
{\sphinxbfcode{\sphinxupquote{stop}}}
{}
{{ $\rightarrow$ None}}
\pysigstopsignatures
\sphinxAtStartPar
Arrête la boucle de polling.

\end{fulllineitems}


\end{fulllineitems}



\paragraph{Module contents}
\label{\detokenize{api/hardware.transports:module-hardware.transports}}\label{\detokenize{api/hardware.transports:module-contents}}\index{module@\spxentry{module}!hardware.transports@\spxentry{hardware.transports}}\index{hardware.transports@\spxentry{hardware.transports}!module@\spxentry{module}}

\subsection{Module contents}
\label{\detokenize{api/hardware:module-hardware}}\label{\detokenize{api/hardware:module-contents}}\index{module@\spxentry{module}!hardware@\spxentry{hardware}}\index{hardware@\spxentry{hardware}!module@\spxentry{module}}
\sphinxstepscope


\section{main module}
\label{\detokenize{api/main:main-module}}\label{\detokenize{api/main::doc}}
\sphinxstepscope


\section{modes package}
\label{\detokenize{api/modes:modes-package}}\label{\detokenize{api/modes::doc}}

\subsection{Submodules}
\label{\detokenize{api/modes:submodules}}

\subsection{modes.mode\_caliper module}
\label{\detokenize{api/modes:module-modes.mode_caliper}}\label{\detokenize{api/modes:modes-mode-caliper-module}}\index{module@\spxentry{module}!modes.mode\_caliper@\spxentry{modes.mode\_caliper}}\index{modes.mode\_caliper@\spxentry{modes.mode\_caliper}!module@\spxentry{module}}
\sphinxAtStartPar
Mode « Pied à coulisse » pour l’assistant ReadForMe.

\sphinxAtStartPar
Ce module encapsule toute la logique « métier » et l’interface utilisateur pour
les mesures de distance sans fil :
\begin{itemize}
\item {} \begin{description}
\sphinxlineitem{Gestion de la voix :}\begin{itemize}
\item {} 
\sphinxAtStartPar
Toutes les annonces passent par le \sphinxtitleref{Speaker} central.

\end{itemize}

\end{description}

\item {} \begin{description}
\sphinxlineitem{Gestion du matériel :}\begin{itemize}
\item {} 
\sphinxAtStartPar
Dialogue avec le récepteur radio via \sphinxtitleref{CaliperDriver}.

\end{itemize}

\end{description}

\item {} \begin{description}
\sphinxlineitem{Logique UX (Bouton poussoir) :}\begin{itemize}
\item {} 
\sphinxAtStartPar
Appui court   : Annoncer la dernière mesure reçue.

\item {} 
\sphinxAtStartPar
Appui long    : Annoncer l’état de connexion radio (SPI/NRF24).

\item {} 
\sphinxAtStartPar
Double appui  : Activer ou désactiver la lecture automatique périodique.

\end{itemize}

\end{description}

\end{itemize}

\sphinxAtStartPar
Ce mode suit le cycle de vie défini par \sphinxtitleref{core.mode\_base.Mode}.
\index{ModeCaliper (classe dans modes.mode\_caliper)@\spxentry{ModeCaliper}\spxextra{classe dans modes.mode\_caliper}}

\begin{fulllineitems}
\phantomsection\label{\detokenize{api/modes:modes.mode_caliper.ModeCaliper}}
\pysigstartsignatures
\pysiglinewithargsret
{\sphinxbfcode{\sphinxupquote{\DUrole{k}{class}\DUrole{w}{ }}}\sphinxcode{\sphinxupquote{modes.mode\_caliper.}}\sphinxbfcode{\sphinxupquote{ModeCaliper}}}
{\sphinxparam{\DUrole{n}{speaker}\DUrole{p}{:}\DUrole{w}{ }\DUrole{n}{{\hyperref[\detokenize{api/core:core.speaker.Speaker}]{\sphinxcrossref{Speaker}}}}}}
{}
\pysigstopsignatures
\sphinxAtStartPar
Bases : {\hyperref[\detokenize{api/core:core.mode_base.Mode}]{\sphinxcrossref{\sphinxcode{\sphinxupquote{Mode}}}}}

\sphinxAtStartPar
Mode « Pied à coulisse » de l’assistant.


\subsubsection{Rôle}
\label{\detokenize{api/modes:role}}
\sphinxAtStartPar
Fournit une interface vocale pour le pied à coulisse numérique. Il traduit
les paquets radio reçus en annonces audibles et gère les interactions
utilisateur via le bouton du sélecteur rotatif.
\index{\_\_init\_\_() (méthode modes.mode\_caliper.ModeCaliper)@\spxentry{\_\_init\_\_()}\spxextra{méthode modes.mode\_caliper.ModeCaliper}}

\begin{fulllineitems}
\phantomsection\label{\detokenize{api/modes:modes.mode_caliper.ModeCaliper.__init__}}
\pysigstartsignatures
\pysiglinewithargsret
{\sphinxbfcode{\sphinxupquote{\_\_init\_\_}}}
{\sphinxparam{\DUrole{n}{speaker}\DUrole{p}{:}\DUrole{w}{ }\DUrole{n}{{\hyperref[\detokenize{api/core:core.speaker.Speaker}]{\sphinxcrossref{Speaker}}}}}}
{}
\pysigstopsignatures
\sphinxAtStartPar
Initialise le mode Pied à coulisse.


\paragraph{Paramètres}
\label{\detokenize{api/modes:parametres}}\begin{description}
\sphinxlineitem{speaker}{[}Speaker{]}
\sphinxAtStartPar
Instance du gestionnaire vocal central pour les annonces.

\end{description}

\end{fulllineitems}

\index{on\_enter() (méthode modes.mode\_caliper.ModeCaliper)@\spxentry{on\_enter()}\spxextra{méthode modes.mode\_caliper.ModeCaliper}}

\begin{fulllineitems}
\phantomsection\label{\detokenize{api/modes:modes.mode_caliper.ModeCaliper.on_enter}}
\pysigstartsignatures
\pysiglinewithargsret
{\sphinxbfcode{\sphinxupquote{on\_enter}}}
{}
{{ $\rightarrow$ None}}
\pysigstopsignatures
\sphinxAtStartPar
Hook appelé lors de la sélection du mode.
\sphinxhyphen{} Désactive la lecture automatique par sécurité.
\sphinxhyphen{} Démarre le thread de lecture radio.

\end{fulllineitems}

\index{on\_exit() (méthode modes.mode\_caliper.ModeCaliper)@\spxentry{on\_exit()}\spxextra{méthode modes.mode\_caliper.ModeCaliper}}

\begin{fulllineitems}
\phantomsection\label{\detokenize{api/modes:modes.mode_caliper.ModeCaliper.on_exit}}
\pysigstartsignatures
\pysiglinewithargsret
{\sphinxbfcode{\sphinxupquote{on\_exit}}}
{}
{{ $\rightarrow$ None}}
\pysigstopsignatures
\sphinxAtStartPar
Hook appelé quand on change de mode.
\sphinxhyphen{} Arrête les annonces automatiques.
\sphinxhyphen{} Coupe proprement le driver radio.

\end{fulllineitems}

\index{on\_short\_press() (méthode modes.mode\_caliper.ModeCaliper)@\spxentry{on\_short\_press()}\spxextra{méthode modes.mode\_caliper.ModeCaliper}}

\begin{fulllineitems}
\phantomsection\label{\detokenize{api/modes:modes.mode_caliper.ModeCaliper.on_short_press}}
\pysigstartsignatures
\pysiglinewithargsret
{\sphinxbfcode{\sphinxupquote{on\_short\_press}}}
{}
{{ $\rightarrow$ None}}
\pysigstopsignatures
\sphinxAtStartPar
Appui court : Annonce vocalement la dernière mesure mémorisée.

\sphinxAtStartPar
Gère le formatage du signe et de l’unité (millimètres ou pouces).
Exemple : « moins 12.45 millimètres »

\end{fulllineitems}

\index{on\_long\_press() (méthode modes.mode\_caliper.ModeCaliper)@\spxentry{on\_long\_press()}\spxextra{méthode modes.mode\_caliper.ModeCaliper}}

\begin{fulllineitems}
\phantomsection\label{\detokenize{api/modes:modes.mode_caliper.ModeCaliper.on_long_press}}
\pysigstartsignatures
\pysiglinewithargsret
{\sphinxbfcode{\sphinxupquote{on\_long\_press}}}
{}
{{ $\rightarrow$ None}}
\pysigstopsignatures
\sphinxAtStartPar
Appui long : Annonce l’état du matériel.
Utile pour diagnostiquer si le module NRF24 est bien branché au Pi.

\end{fulllineitems}

\index{on\_double\_press() (méthode modes.mode\_caliper.ModeCaliper)@\spxentry{on\_double\_press()}\spxextra{méthode modes.mode\_caliper.ModeCaliper}}

\begin{fulllineitems}
\phantomsection\label{\detokenize{api/modes:modes.mode_caliper.ModeCaliper.on_double_press}}
\pysigstartsignatures
\pysiglinewithargsret
{\sphinxbfcode{\sphinxupquote{on\_double\_press}}}
{}
{{ $\rightarrow$ None}}
\pysigstopsignatures
\sphinxAtStartPar
Double appui : Bascule l’état de la lecture automatique.

\end{fulllineitems}


\end{fulllineitems}



\subsection{modes.mode\_datetime module}
\label{\detokenize{api/modes:module-modes.mode_datetime}}\label{\detokenize{api/modes:modes-mode-datetime-module}}\index{module@\spxentry{module}!modes.mode\_datetime@\spxentry{modes.mode\_datetime}}\index{modes.mode\_datetime@\spxentry{modes.mode\_datetime}!module@\spxentry{module}}
\sphinxAtStartPar
Mode « Date et heure ».

\sphinxAtStartPar
Ce module définit un mode très simple qui exploite l’horloge système
pour annoncer :
\begin{itemize}
\item {} 
\sphinxAtStartPar
l’heure actuelle (appui court),

\item {} 
\sphinxAtStartPar
la date du jour (appui long).

\end{itemize}

\sphinxAtStartPar
Il ne gère aucune logique matérielle directement : il s’appuie sur
le \sphinxtitleref{ModeManager} (pour la sélection du mode) et sur le \sphinxtitleref{Speaker}
(pour la synthèse vocale).
\index{ModeDateHeure (classe dans modes.mode\_datetime)@\spxentry{ModeDateHeure}\spxextra{classe dans modes.mode\_datetime}}

\begin{fulllineitems}
\phantomsection\label{\detokenize{api/modes:modes.mode_datetime.ModeDateHeure}}
\pysigstartsignatures
\pysiglinewithargsret
{\sphinxbfcode{\sphinxupquote{\DUrole{k}{class}\DUrole{w}{ }}}\sphinxcode{\sphinxupquote{modes.mode\_datetime.}}\sphinxbfcode{\sphinxupquote{ModeDateHeure}}}
{\sphinxparam{\DUrole{n}{speaker}\DUrole{p}{:}\DUrole{w}{ }\DUrole{n}{{\hyperref[\detokenize{api/core:core.speaker.Speaker}]{\sphinxcrossref{Speaker}}}}}}
{}
\pysigstopsignatures
\sphinxAtStartPar
Bases : {\hyperref[\detokenize{api/core:core.mode_base.Mode}]{\sphinxcrossref{\sphinxcode{\sphinxupquote{Mode}}}}}

\sphinxAtStartPar
Mode « Date et heure » de l’assistant.


\subsubsection{Comportement UX}
\label{\detokenize{api/modes:comportement-ux}}\begin{itemize}
\item {} \begin{description}
\sphinxlineitem{À l’entrée du mode (\sphinxtitleref{on\_enter}) :}\begin{itemize}
\item {} 
\sphinxAtStartPar
rien n’est annoncé ici, le \sphinxtitleref{ModeManager} annonce déjà
« Mode X : Date et heure ».

\end{itemize}

\end{description}

\item {} \begin{description}
\sphinxlineitem{Appui court (\sphinxtitleref{on\_short\_press}) :}\begin{itemize}
\item {} 
\sphinxAtStartPar
annonce l’heure système actuelle (ex : « Il est 14 heures 32 »).

\end{itemize}

\end{description}

\item {} \begin{description}
\sphinxlineitem{Appui long (\sphinxtitleref{on\_long\_press}) :}\begin{itemize}
\item {} 
\sphinxAtStartPar
annonce la date du jour (ex : « Nous sommes le mardi 5 juin 2025 »).

\end{itemize}

\end{description}

\end{itemize}


\subsubsection{Remarques}
\label{\detokenize{api/modes:remarques}}\begin{itemize}
\item {} 
\sphinxAtStartPar
Les libellés de jour et de mois (\sphinxtitleref{\%A}, \sphinxtitleref{\%B}) dépendent de la locale
système (ex : \sphinxtitleref{fr\_FR} recommandé pour avoir les noms en français).

\end{itemize}
\index{\_\_init\_\_() (méthode modes.mode\_datetime.ModeDateHeure)@\spxentry{\_\_init\_\_()}\spxextra{méthode modes.mode\_datetime.ModeDateHeure}}

\begin{fulllineitems}
\phantomsection\label{\detokenize{api/modes:modes.mode_datetime.ModeDateHeure.__init__}}
\pysigstartsignatures
\pysiglinewithargsret
{\sphinxbfcode{\sphinxupquote{\_\_init\_\_}}}
{\sphinxparam{\DUrole{n}{speaker}\DUrole{p}{:}\DUrole{w}{ }\DUrole{n}{{\hyperref[\detokenize{api/core:core.speaker.Speaker}]{\sphinxcrossref{Speaker}}}}}}
{}
\pysigstopsignatures
\sphinxAtStartPar
Initialise le mode « Date et heure ».


\paragraph{Paramètres}
\label{\detokenize{api/modes:id1}}\begin{description}
\sphinxlineitem{speaker}{[}Speaker{]}
\sphinxAtStartPar
Instance du \sphinxtitleref{Speaker} central utilisée pour prononcer les messages.

\end{description}

\end{fulllineitems}

\index{on\_enter() (méthode modes.mode\_datetime.ModeDateHeure)@\spxentry{on\_enter()}\spxextra{méthode modes.mode\_datetime.ModeDateHeure}}

\begin{fulllineitems}
\phantomsection\label{\detokenize{api/modes:modes.mode_datetime.ModeDateHeure.on_enter}}
\pysigstartsignatures
\pysiglinewithargsret
{\sphinxbfcode{\sphinxupquote{on\_enter}}}
{}
{{ $\rightarrow$ None}}
\pysigstopsignatures
\sphinxAtStartPar
Hook appelé lorsque l’on ARRIVE sur ce mode.

\sphinxAtStartPar
Ici, on ne fait rien :
\sphinxhyphen{} le \sphinxtitleref{ModeManager} se charge déjà d’annoncer le nom du mode,
\sphinxhyphen{} cela évite d’ajouter du bruit vocal à chaque changement de mode.

\sphinxAtStartPar
Ce hook existe principalement pour symétrie et pour d’éventuelles
évolutions (par ex. « rafraîchir » un cache, etc.).

\end{fulllineitems}

\index{on\_exit() (méthode modes.mode\_datetime.ModeDateHeure)@\spxentry{on\_exit()}\spxextra{méthode modes.mode\_datetime.ModeDateHeure}}

\begin{fulllineitems}
\phantomsection\label{\detokenize{api/modes:modes.mode_datetime.ModeDateHeure.on_exit}}
\pysigstartsignatures
\pysiglinewithargsret
{\sphinxbfcode{\sphinxupquote{on\_exit}}}
{}
{{ $\rightarrow$ None}}
\pysigstopsignatures
\sphinxAtStartPar
Hook appelé lorsque l’on QUITTE ce mode.

\sphinxAtStartPar
Aucun nettoyage spécifique n’est nécessaire pour ce mode, donc
l’implémentation est vide pour le moment.

\end{fulllineitems}

\index{on\_short\_press() (méthode modes.mode\_datetime.ModeDateHeure)@\spxentry{on\_short\_press()}\spxextra{méthode modes.mode\_datetime.ModeDateHeure}}

\begin{fulllineitems}
\phantomsection\label{\detokenize{api/modes:modes.mode_datetime.ModeDateHeure.on_short_press}}
\pysigstartsignatures
\pysiglinewithargsret
{\sphinxbfcode{\sphinxupquote{on\_short\_press}}}
{}
{{ $\rightarrow$ None}}
\pysigstopsignatures
\sphinxAtStartPar
Appui court : annonce l’heure système actuelle.
\begin{description}
\sphinxlineitem{Exemple de rendu vocal :}
\sphinxAtStartPar
« Il est 14 heures 32 »

\end{description}


\paragraph{Implémentation}
\label{\detokenize{api/modes:implementation}}\begin{itemize}
\item {} 
\sphinxAtStartPar
Récupère l’heure courante via \sphinxtitleref{datetime.now()}.

\item {} 
\sphinxAtStartPar
Formate l’heure avec \sphinxtitleref{strftime(« \%H heures \%M »)}.

\item {} 
\sphinxAtStartPar
Envoie la phrase complète au \sphinxtitleref{Speaker}.

\end{itemize}

\end{fulllineitems}

\index{on\_long\_press() (méthode modes.mode\_datetime.ModeDateHeure)@\spxentry{on\_long\_press()}\spxextra{méthode modes.mode\_datetime.ModeDateHeure}}

\begin{fulllineitems}
\phantomsection\label{\detokenize{api/modes:modes.mode_datetime.ModeDateHeure.on_long_press}}
\pysigstartsignatures
\pysiglinewithargsret
{\sphinxbfcode{\sphinxupquote{on\_long\_press}}}
{}
{{ $\rightarrow$ None}}
\pysigstopsignatures
\sphinxAtStartPar
Appui long : annonce la date du jour.
\begin{description}
\sphinxlineitem{Exemple de rendu vocal :}
\sphinxAtStartPar
« Nous sommes le mardi 5 juin 2025 »

\end{description}


\paragraph{Implémentation}
\label{\detokenize{api/modes:id2}}\begin{itemize}
\item {} 
\sphinxAtStartPar
Utilise \sphinxtitleref{datetime.now()} pour récupérer la date actuelle.

\item {} \begin{description}
\sphinxlineitem{Décompose :}\begin{itemize}
\item {} 
\sphinxAtStartPar
jour de la semaine : \sphinxtitleref{\%A}

\item {} 
\sphinxAtStartPar
jour du mois       : \sphinxtitleref{\%d} (converti en int pour éviter les zéros en tête)

\item {} 
\sphinxAtStartPar
mois               : \sphinxtitleref{\%B}

\item {} 
\sphinxAtStartPar
année              : \sphinxtitleref{\%Y}

\end{itemize}

\end{description}

\item {} 
\sphinxAtStartPar
Construit une phrase en français et la passe au \sphinxtitleref{Speaker}.

\end{itemize}


\paragraph{Remarque}
\label{\detokenize{api/modes:remarque}}\begin{itemize}
\item {} 
\sphinxAtStartPar
Selon la configuration locale du système (locale),
\sphinxtitleref{\%A} et \sphinxtitleref{\%B} peuvent être en anglais. Pour forcer le français,
il est recommandé de configurer la locale en \sphinxtitleref{fr\_FR} au niveau OS.

\end{itemize}

\end{fulllineitems}


\end{fulllineitems}



\subsection{modes.mode\_dummy module}
\label{\detokenize{api/modes:module-modes.mode_dummy}}\label{\detokenize{api/modes:modes-mode-dummy-module}}\index{module@\spxentry{module}!modes.mode\_dummy@\spxentry{modes.mode\_dummy}}\index{modes.mode\_dummy@\spxentry{modes.mode\_dummy}!module@\spxentry{module}}
\sphinxAtStartPar
Mode de démonstration / test.

\sphinxAtStartPar
Ce mode n’a aucune logique métier réelle : il sert uniquement
à valider les interactions entre :
\begin{itemize}
\item {} 
\sphinxAtStartPar
le sélecteur rotatif (rotations, appui court, long, double),

\item {} 
\sphinxAtStartPar
le ModeManager,

\item {} 
\sphinxAtStartPar
la synthèse vocale (Speaker).

\end{itemize}

\sphinxAtStartPar
Il est utile comme « mode de développement » pour vérifier que
l’enchaînement complet hardware \sphinxhyphen{}\textgreater{} callbacks \sphinxhyphen{}\textgreater{} mode \sphinxhyphen{}\textgreater{} TTS
fonctionne correctement.

\sphinxAtStartPar
Ce mode peut servir de modèle minimal pour créer rapidement
un nouveau mode.
\index{ModeDummy (classe dans modes.mode\_dummy)@\spxentry{ModeDummy}\spxextra{classe dans modes.mode\_dummy}}

\begin{fulllineitems}
\phantomsection\label{\detokenize{api/modes:modes.mode_dummy.ModeDummy}}
\pysigstartsignatures
\pysiglinewithargsret
{\sphinxbfcode{\sphinxupquote{\DUrole{k}{class}\DUrole{w}{ }}}\sphinxcode{\sphinxupquote{modes.mode\_dummy.}}\sphinxbfcode{\sphinxupquote{ModeDummy}}}
{\sphinxparam{\DUrole{n}{speaker}}}
{}
\pysigstopsignatures
\sphinxAtStartPar
Bases : {\hyperref[\detokenize{api/core:core.mode_base.Mode}]{\sphinxcrossref{\sphinxcode{\sphinxupquote{Mode}}}}}

\sphinxAtStartPar
Mode « test » basique.

\sphinxAtStartPar
Comportement :
\sphinxhyphen{} on\_enter       \sphinxhyphen{}\textgreater{} ne fait rien
\sphinxhyphen{} on\_short\_press \sphinxhyphen{}\textgreater{} annonce un message simple
\sphinxhyphen{} on\_long\_press  \sphinxhyphen{}\textgreater{} annonce un message spécifique à l’appui long
\sphinxhyphen{} on\_double\_press \sphinxhyphen{}\textgreater{} confirme le détecteur de double\sphinxhyphen{}clic

\sphinxAtStartPar
Ce mode n’effectue aucune action matérielle.
\index{\_\_init\_\_() (méthode modes.mode\_dummy.ModeDummy)@\spxentry{\_\_init\_\_()}\spxextra{méthode modes.mode\_dummy.ModeDummy}}

\begin{fulllineitems}
\phantomsection\label{\detokenize{api/modes:modes.mode_dummy.ModeDummy.__init__}}
\pysigstartsignatures
\pysiglinewithargsret
{\sphinxbfcode{\sphinxupquote{\_\_init\_\_}}}
{\sphinxparam{\DUrole{n}{speaker}}}
{}
\pysigstopsignatures\begin{quote}\begin{description}
\sphinxlineitem{Paramètres}
\sphinxAtStartPar
\sphinxstyleliteralstrong{\sphinxupquote{speaker}} ({\hyperref[\detokenize{api/core:core.speaker.Speaker}]{\sphinxcrossref{\sphinxstyleliteralemphasis{\sphinxupquote{Speaker}}}}}) \textendash{} Synthèse vocale centrale, utilisée pour les annonces.

\end{description}\end{quote}

\end{fulllineitems}

\index{on\_enter() (méthode modes.mode\_dummy.ModeDummy)@\spxentry{on\_enter()}\spxextra{méthode modes.mode\_dummy.ModeDummy}}

\begin{fulllineitems}
\phantomsection\label{\detokenize{api/modes:modes.mode_dummy.ModeDummy.on_enter}}
\pysigstartsignatures
\pysiglinewithargsret
{\sphinxbfcode{\sphinxupquote{on\_enter}}}
{}
{}
\pysigstopsignatures
\sphinxAtStartPar
Appelé quand on arrive sur ce mode.
Aucun feedback vocal ici pour ne pas surcharger :
le ModeManager annonce déjà « Mode X : Mode test ».

\end{fulllineitems}

\index{on\_exit() (méthode modes.mode\_dummy.ModeDummy)@\spxentry{on\_exit()}\spxextra{méthode modes.mode\_dummy.ModeDummy}}

\begin{fulllineitems}
\phantomsection\label{\detokenize{api/modes:modes.mode_dummy.ModeDummy.on_exit}}
\pysigstartsignatures
\pysiglinewithargsret
{\sphinxbfcode{\sphinxupquote{on\_exit}}}
{}
{}
\pysigstopsignatures
\sphinxAtStartPar
Appelé lorsqu’on quitte ce mode.
Rien à nettoyer.

\end{fulllineitems}

\index{on\_short\_press() (méthode modes.mode\_dummy.ModeDummy)@\spxentry{on\_short\_press()}\spxextra{méthode modes.mode\_dummy.ModeDummy}}

\begin{fulllineitems}
\phantomsection\label{\detokenize{api/modes:modes.mode_dummy.ModeDummy.on_short_press}}
\pysigstartsignatures
\pysiglinewithargsret
{\sphinxbfcode{\sphinxupquote{on\_short\_press}}}
{}
{}
\pysigstopsignatures
\sphinxAtStartPar
Annonce une confirmation d’appui court.

\end{fulllineitems}

\index{on\_long\_press() (méthode modes.mode\_dummy.ModeDummy)@\spxentry{on\_long\_press()}\spxextra{méthode modes.mode\_dummy.ModeDummy}}

\begin{fulllineitems}
\phantomsection\label{\detokenize{api/modes:modes.mode_dummy.ModeDummy.on_long_press}}
\pysigstartsignatures
\pysiglinewithargsret
{\sphinxbfcode{\sphinxupquote{on\_long\_press}}}
{}
{}
\pysigstopsignatures
\sphinxAtStartPar
Annonce une confirmation d’appui long.

\end{fulllineitems}

\index{on\_double\_press() (méthode modes.mode\_dummy.ModeDummy)@\spxentry{on\_double\_press()}\spxextra{méthode modes.mode\_dummy.ModeDummy}}

\begin{fulllineitems}
\phantomsection\label{\detokenize{api/modes:modes.mode_dummy.ModeDummy.on_double_press}}
\pysigstartsignatures
\pysiglinewithargsret
{\sphinxbfcode{\sphinxupquote{on\_double\_press}}}
{}
{}
\pysigstopsignatures
\sphinxAtStartPar
Annonce un double appui détecté.

\end{fulllineitems}


\end{fulllineitems}



\subsection{modes.mode\_multimeter module}
\label{\detokenize{api/modes:module-modes.mode_multimeter}}\label{\detokenize{api/modes:modes-mode-multimeter-module}}\index{module@\spxentry{module}!modes.mode\_multimeter@\spxentry{modes.mode\_multimeter}}\index{modes.mode\_multimeter@\spxentry{modes.mode\_multimeter}!module@\spxentry{module}}
\sphinxAtStartPar
Mode « Multimètre OWON 16 » pour l’assistant.

\sphinxAtStartPar
Ce module encapsule toute la logique « métier » côté utilisateur pour le multimètre :
\begin{itemize}
\item {} \begin{description}
\sphinxlineitem{Gestion de la voix :}\begin{itemize}
\item {} 
\sphinxAtStartPar
Toutes les annonces passent par le \sphinxtitleref{Speaker} central.

\end{itemize}

\end{description}

\item {} \begin{description}
\sphinxlineitem{Gestion du matériel :}\begin{itemize}
\item {} 
\sphinxAtStartPar
Dialogue avec le multimètre via \sphinxtitleref{OwonMultimeterDriver}
(qui lui\sphinxhyphen{}même utilise le transport BLE générique \sphinxtitleref{BleClient}).

\end{itemize}

\end{description}

\item {} \begin{description}
\sphinxlineitem{Logique UX :}\begin{itemize}
\item {} 
\sphinxAtStartPar
Appui court  : lire la dernière mesure.

\item {} 
\sphinxAtStartPar
Appui long   : annoncer l’état de connexion du multimètre.

\item {} 
\sphinxAtStartPar
Double appui : activer / désactiver une lecture automatique périodique
(toutes les X secondes).

\end{itemize}

\end{description}

\end{itemize}

\sphinxAtStartPar
Le mode ne contient AUCUN code BLE direct et ne gère pas les broches GPIO :
il consomme simplement un driver de plus bas niveau.
\index{ModeMultimetre (classe dans modes.mode\_multimeter)@\spxentry{ModeMultimetre}\spxextra{classe dans modes.mode\_multimeter}}

\begin{fulllineitems}
\phantomsection\label{\detokenize{api/modes:modes.mode_multimeter.ModeMultimetre}}
\pysigstartsignatures
\pysiglinewithargsret
{\sphinxbfcode{\sphinxupquote{\DUrole{k}{class}\DUrole{w}{ }}}\sphinxcode{\sphinxupquote{modes.mode\_multimeter.}}\sphinxbfcode{\sphinxupquote{ModeMultimetre}}}
{\sphinxparam{\DUrole{n}{speaker}\DUrole{p}{:}\DUrole{w}{ }\DUrole{n}{{\hyperref[\detokenize{api/core:core.speaker.Speaker}]{\sphinxcrossref{Speaker}}}}}\sphinxparamcomma \sphinxparam{\DUrole{keyword-only-separator}{\DUrole{o}{\sphinxstyleabbreviation{*}}}}\sphinxparamcomma \sphinxparam{\DUrole{n}{auto\_read\_interval}\DUrole{p}{:}\DUrole{w}{ }\DUrole{n}{float}\DUrole{w}{ }\DUrole{o}{=}\DUrole{w}{ }\DUrole{default_value}{5.0}}\sphinxparamcomma \sphinxparam{\DUrole{n}{driver\_config}\DUrole{p}{:}\DUrole{w}{ }\DUrole{n}{Dict\DUrole{p}{{[}}str\DUrole{p}{,}\DUrole{w}{ }Any\DUrole{p}{{]}}\DUrole{w}{ }\DUrole{p}{|}\DUrole{w}{ }None}\DUrole{w}{ }\DUrole{o}{=}\DUrole{w}{ }\DUrole{default_value}{None}}}
{}
\pysigstopsignatures
\sphinxAtStartPar
Bases : {\hyperref[\detokenize{api/core:core.mode_base.Mode}]{\sphinxcrossref{\sphinxcode{\sphinxupquote{Mode}}}}}

\sphinxAtStartPar
Mode « Multimètre » de l’assistant.


\subsubsection{Rôle}
\label{\detokenize{api/modes:id3}}
\sphinxAtStartPar
Ce mode fournit une interface vocale par\sphinxhyphen{}dessus le multimètre OWON 16 :
il traduit les actions utilisateur (appuis sur le bouton, double\sphinxhyphen{}clic,
sélection du mode) en :
\begin{itemize}
\item {} 
\sphinxAtStartPar
commandes haut niveau vers \sphinxtitleref{OwonMultimeterDriver} (start/stop),

\item {} 
\sphinxAtStartPar
annonces vocales lisibles pour un utilisateur non voyant.

\end{itemize}


\subsubsection{Comportement UX}
\label{\detokenize{api/modes:id4}}\begin{itemize}
\item {} \begin{description}
\sphinxlineitem{\sphinxtitleref{on\_enter} :}\begin{itemize}
\item {} 
\sphinxAtStartPar
démarre la recherche / connexion BLE en arrière\sphinxhyphen{}plan via le driver.

\item {} 
\sphinxAtStartPar
remet à zéro le mode de lecture automatique (on revient en manuel).

\item {} 
\sphinxAtStartPar
l’annonce « Mode X : Multimètre » est gérée par le \sphinxtitleref{ModeManager}.

\end{itemize}

\end{description}

\item {} \begin{description}
\sphinxlineitem{\sphinxtitleref{on\_exit} :}\begin{itemize}
\item {} 
\sphinxAtStartPar
arrête proprement le driver OWON (scan + connexion BLE).

\item {} 
\sphinxAtStartPar
désactive la lecture automatique et annule les timers.

\end{itemize}

\end{description}

\item {} \begin{description}
\sphinxlineitem{\sphinxtitleref{on\_short\_press} (clic simple) :}\begin{itemize}
\item {} 
\sphinxAtStartPar
annonce la dernière mesure disponible (ex : « 3.27 VOLT DC »).

\item {} 
\sphinxAtStartPar
si aucune mesure n’est encore disponible : message d’erreur vocal.

\end{itemize}

\end{description}

\item {} \begin{description}
\sphinxlineitem{\sphinxtitleref{on\_long\_press} (appui long) :}\begin{itemize}
\item {} 
\sphinxAtStartPar
annonce l’état de connexion actuel (connecté / en recherche).

\end{itemize}

\end{description}

\item {} \begin{description}
\sphinxlineitem{\sphinxtitleref{on\_double\_press} (double clic) :}\begin{itemize}
\item {} 
\sphinxAtStartPar
active / désactive un mode de lecture automatique :
toutes les \sphinxtitleref{self.\_auto\_read\_interval} secondes, la dernière
mesure est annoncée, tant que le multimètre est connecté.

\end{itemize}

\end{description}

\end{itemize}


\subsubsection{Détails techniques}
\label{\detokenize{api/modes:details-techniques}}\begin{itemize}
\item {} 
\sphinxAtStartPar
La conversion « nombre \sphinxhyphen{}\textgreater{} texte » est centralisée dans
\sphinxtitleref{\_format\_value\_for\_speech()} afin de gérer proprement :
\begin{itemize}
\item {} 
\sphinxAtStartPar
les valeurs négatives (« moins 1.3 »),

\item {} 
\sphinxAtStartPar
le format numérique passé au TTS.

\end{itemize}

\item {} 
\sphinxAtStartPar
La lecture automatique est pilotée par un \sphinxtitleref{threading.Timer}
qui se réarme lui\sphinxhyphen{}même (\sphinxtitleref{\_schedule\_auto\_read} / \sphinxtitleref{\_auto\_read\_tick}).

\item {} 
\sphinxAtStartPar
Quand la connexion BLE est perdue, le mode auto reste activé
mais ne parle pas tant que le multimètre n’est pas reconnecté.

\end{itemize}
\index{\_\_init\_\_() (méthode modes.mode\_multimeter.ModeMultimetre)@\spxentry{\_\_init\_\_()}\spxextra{méthode modes.mode\_multimeter.ModeMultimetre}}

\begin{fulllineitems}
\phantomsection\label{\detokenize{api/modes:modes.mode_multimeter.ModeMultimetre.__init__}}
\pysigstartsignatures
\pysiglinewithargsret
{\sphinxbfcode{\sphinxupquote{\_\_init\_\_}}}
{\sphinxparam{\DUrole{n}{speaker}\DUrole{p}{:}\DUrole{w}{ }\DUrole{n}{{\hyperref[\detokenize{api/core:core.speaker.Speaker}]{\sphinxcrossref{Speaker}}}}}\sphinxparamcomma \sphinxparam{\DUrole{keyword-only-separator}{\DUrole{o}{\sphinxstyleabbreviation{*}}}}\sphinxparamcomma \sphinxparam{\DUrole{n}{auto\_read\_interval}\DUrole{p}{:}\DUrole{w}{ }\DUrole{n}{float}\DUrole{w}{ }\DUrole{o}{=}\DUrole{w}{ }\DUrole{default_value}{5.0}}\sphinxparamcomma \sphinxparam{\DUrole{n}{driver\_config}\DUrole{p}{:}\DUrole{w}{ }\DUrole{n}{Dict\DUrole{p}{{[}}str\DUrole{p}{,}\DUrole{w}{ }Any\DUrole{p}{{]}}\DUrole{w}{ }\DUrole{p}{|}\DUrole{w}{ }None}\DUrole{w}{ }\DUrole{o}{=}\DUrole{w}{ }\DUrole{default_value}{None}}}
{}
\pysigstopsignatures
\sphinxAtStartPar
Initialise le mode Multimètre.


\paragraph{Paramètres}
\label{\detokenize{api/modes:id5}}\begin{description}
\sphinxlineitem{speaker}{[}Speaker{]}
\sphinxAtStartPar
Instance du \sphinxtitleref{Speaker} central utilisée pour toutes les annonces vocales.

\end{description}

\end{fulllineitems}

\index{on\_enter() (méthode modes.mode\_multimeter.ModeMultimetre)@\spxentry{on\_enter()}\spxextra{méthode modes.mode\_multimeter.ModeMultimetre}}

\begin{fulllineitems}
\phantomsection\label{\detokenize{api/modes:modes.mode_multimeter.ModeMultimetre.on_enter}}
\pysigstartsignatures
\pysiglinewithargsret
{\sphinxbfcode{\sphinxupquote{on\_enter}}}
{}
{{ $\rightarrow$ None}}
\pysigstopsignatures
\sphinxAtStartPar
Hook appelé quand on arrive sur ce mode via le sélecteur rotatif.


\paragraph{Actions réalisées}
\label{\detokenize{api/modes:actions-realisees}}\begin{itemize}
\item {} 
\sphinxAtStartPar
Désactive la lecture automatique (on revient en mode manuel).

\item {} 
\sphinxAtStartPar
Annule tout timer éventuellement encore actif.

\item {} \begin{description}
\sphinxlineitem{Démarre le driver OWON :}\begin{itemize}
\item {} 
\sphinxAtStartPar
scan BLE en arrière\sphinxhyphen{}plan,

\item {} 
\sphinxAtStartPar
tentatives de connexion automatiques.

\end{itemize}

\end{description}

\end{itemize}


\paragraph{Remarque}
\label{\detokenize{api/modes:id6}}
\sphinxAtStartPar
Le \sphinxtitleref{ModeManager} se charge déjà d’annoncer « Mode X : Multimètre »,
donc on ne parle pas ici pour éviter les doublons.

\end{fulllineitems}

\index{on\_exit() (méthode modes.mode\_multimeter.ModeMultimetre)@\spxentry{on\_exit()}\spxextra{méthode modes.mode\_multimeter.ModeMultimetre}}

\begin{fulllineitems}
\phantomsection\label{\detokenize{api/modes:modes.mode_multimeter.ModeMultimetre.on_exit}}
\pysigstartsignatures
\pysiglinewithargsret
{\sphinxbfcode{\sphinxupquote{on\_exit}}}
{}
{{ $\rightarrow$ None}}
\pysigstopsignatures
\sphinxAtStartPar
Hook appelé quand on quitte ce mode.


\paragraph{Actions réalisées}
\label{\detokenize{api/modes:id7}}\begin{itemize}
\item {} 
\sphinxAtStartPar
Désactive le mode de lecture automatique.

\item {} 
\sphinxAtStartPar
Annule le timer de lecture automatique.

\item {} 
\sphinxAtStartPar
Arrête le driver OWON (arrêt des scans / connexions BLE).

\end{itemize}

\end{fulllineitems}

\index{on\_short\_press() (méthode modes.mode\_multimeter.ModeMultimetre)@\spxentry{on\_short\_press()}\spxextra{méthode modes.mode\_multimeter.ModeMultimetre}}

\begin{fulllineitems}
\phantomsection\label{\detokenize{api/modes:modes.mode_multimeter.ModeMultimetre.on_short_press}}
\pysigstartsignatures
\pysiglinewithargsret
{\sphinxbfcode{\sphinxupquote{on\_short\_press}}}
{}
{{ $\rightarrow$ None}}
\pysigstopsignatures
\sphinxAtStartPar
Appui court : annonce la dernière mesure reçue.


\paragraph{Comportement}
\label{\detokenize{api/modes:comportement}}\begin{itemize}
\item {} 
\sphinxAtStartPar
Si aucune mesure n’est disponible (pas encore de trame reçue) :
annonce « Aucune mesure disponible pour le moment. ».

\item {} 
\sphinxAtStartPar
Sinon :
* récupère le dictionnaire simplifié via \sphinxtitleref{get\_last\_measure()},
* formate la valeur pour la voix (\sphinxtitleref{\_format\_value\_for\_speech}),
* concatène la valeur + le nom d’unité tel que décodé par le driver
\begin{quote}

\sphinxAtStartPar
(ex : « MILLI VOLT DC », « OHM NORMAL », « CELSIUS »…),
\end{quote}
\begin{itemize}
\item {} 
\sphinxAtStartPar
envoie la phrase au \sphinxtitleref{Speaker}.

\end{itemize}

\end{itemize}


\paragraph{Exemple de rendu}
\label{\detokenize{api/modes:exemple-de-rendu}}\begin{quote}

\sphinxAtStartPar
« moins 1.3 MILLI VOLT DC »
\end{quote}

\end{fulllineitems}

\index{on\_long\_press() (méthode modes.mode\_multimeter.ModeMultimetre)@\spxentry{on\_long\_press()}\spxextra{méthode modes.mode\_multimeter.ModeMultimetre}}

\begin{fulllineitems}
\phantomsection\label{\detokenize{api/modes:modes.mode_multimeter.ModeMultimetre.on_long_press}}
\pysigstartsignatures
\pysiglinewithargsret
{\sphinxbfcode{\sphinxupquote{on\_long\_press}}}
{}
{{ $\rightarrow$ None}}
\pysigstopsignatures
\sphinxAtStartPar
Appui long : annonce l’état de connexion actuel du multimètre.


\paragraph{Cas gérés}
\label{\detokenize{api/modes:cas-geres}}\begin{itemize}
\item {} \begin{description}
\sphinxlineitem{Si \sphinxtitleref{OwonMultimeterDriver.is\_connected} est True :}
\sphinxAtStartPar
\sphinxhyphen{}\textgreater{} « Le multimètre est connecté. »

\end{description}

\item {} \begin{description}
\sphinxlineitem{Sinon :}\begin{description}
\sphinxlineitem{\sphinxhyphen{}\textgreater{} « Recherche ou connexion au multimètre en cours.}
\sphinxAtStartPar
Veuillez patienter ou vérifier qu’il est allumé. »

\end{description}

\end{description}

\end{itemize}

\end{fulllineitems}

\index{on\_double\_press() (méthode modes.mode\_multimeter.ModeMultimetre)@\spxentry{on\_double\_press()}\spxextra{méthode modes.mode\_multimeter.ModeMultimetre}}

\begin{fulllineitems}
\phantomsection\label{\detokenize{api/modes:modes.mode_multimeter.ModeMultimetre.on_double_press}}
\pysigstartsignatures
\pysiglinewithargsret
{\sphinxbfcode{\sphinxupquote{on\_double\_press}}}
{}
{{ $\rightarrow$ None}}
\pysigstopsignatures
\sphinxAtStartPar
Double appui : bascule entre lecture automatique et lecture manuelle.


\paragraph{Comportement}
\label{\detokenize{api/modes:id8}}\begin{itemize}
\item {} \begin{description}
\sphinxlineitem{Si la lecture automatique est désactivée :}\begin{itemize}
\item {} 
\sphinxAtStartPar
l’active,

\item {} 
\sphinxAtStartPar
annonce « Lecture automatique du multimètre activée. »,

\item {} 
\sphinxAtStartPar
planifie le premier tick avec \sphinxtitleref{\_schedule\_auto\_read()}.

\end{itemize}

\end{description}

\item {} \begin{description}
\sphinxlineitem{Si la lecture automatique est déjà active :}\begin{itemize}
\item {} 
\sphinxAtStartPar
la désactive,

\item {} 
\sphinxAtStartPar
annule le timer associé,

\item {} 
\sphinxAtStartPar
annonce « Lecture automatique désactivée. Retour au mode manuel. ».

\end{itemize}

\end{description}

\end{itemize}

\end{fulllineitems}


\end{fulllineitems}



\subsection{modes.mode\_reading\_machine module}
\label{\detokenize{api/modes:module-modes.mode_reading_machine}}\label{\detokenize{api/modes:modes-mode-reading-machine-module}}\index{module@\spxentry{module}!modes.mode\_reading\_machine@\spxentry{modes.mode\_reading\_machine}}\index{modes.mode\_reading\_machine@\spxentry{modes.mode\_reading\_machine}!module@\spxentry{module}}
\sphinxAtStartPar
Mode « Machine à lire ».

\sphinxAtStartPar
Ce module orchestre le flux utilisateur le plus complet du projet :
\begin{enumerate}
\sphinxsetlistlabels{\arabic}{enumi}{enumii}{}{.}%
\item {} 
\sphinxAtStartPar
Détection de la disponibilité de la caméra.

\item {} 
\sphinxAtStartPar
Capture d’une image via \sphinxtitleref{PiCamera}.

\item {} 
\sphinxAtStartPar
OCR et nettoyage du texte via \sphinxtitleref{TextReaderDevice}.

\item {} 
\sphinxAtStartPar
Synthèse du texte en WAV via \sphinxtitleref{Speaker.synthesize\_to\_file}.

\item {} 
\sphinxAtStartPar
Lecture longue contrôlable via \sphinxtitleref{Speaker.play}.

\end{enumerate}

\sphinxAtStartPar
Le mode ne lit pas directement le keypad GPIO. Il reçoit les touches
non\sphinxhyphen{}globales depuis \sphinxtitleref{ModeManager}, ce qui permet de conserver :
\begin{itemize}
\item {} 
\sphinxAtStartPar
les touches globales volume/vitesse dans tous les modes ;

\item {} 
\sphinxAtStartPar
une logique de contrôle métier isolée ici ;

\item {} 
\sphinxAtStartPar
une implémentation testable avec des faux objets.

\end{itemize}
\index{ModeReadingMachine (classe dans modes.mode\_reading\_machine)@\spxentry{ModeReadingMachine}\spxextra{classe dans modes.mode\_reading\_machine}}

\begin{fulllineitems}
\phantomsection\label{\detokenize{api/modes:modes.mode_reading_machine.ModeReadingMachine}}
\pysigstartsignatures
\pysiglinewithargsret
{\sphinxbfcode{\sphinxupquote{\DUrole{k}{class}\DUrole{w}{ }}}\sphinxcode{\sphinxupquote{modes.mode\_reading\_machine.}}\sphinxbfcode{\sphinxupquote{ModeReadingMachine}}}
{\sphinxparam{\DUrole{n}{speaker}\DUrole{p}{:}\DUrole{w}{ }\DUrole{n}{{\hyperref[\detokenize{api/core:core.speaker.Speaker}]{\sphinxcrossref{Speaker}}}}}\sphinxparamcomma \sphinxparam{\DUrole{keyword-only-separator}{\DUrole{o}{\sphinxstyleabbreviation{*}}}}\sphinxparamcomma \sphinxparam{\DUrole{n}{reading\_config}\DUrole{p}{:}\DUrole{w}{ }\DUrole{n}{Dict\DUrole{p}{{[}}str\DUrole{p}{,}\DUrole{w}{ }Any\DUrole{p}{{]}}\DUrole{w}{ }\DUrole{p}{|}\DUrole{w}{ }None}\DUrole{w}{ }\DUrole{o}{=}\DUrole{w}{ }\DUrole{default_value}{None}}\sphinxparamcomma \sphinxparam{\DUrole{n}{ocr\_config}\DUrole{p}{:}\DUrole{w}{ }\DUrole{n}{Dict\DUrole{p}{{[}}str\DUrole{p}{,}\DUrole{w}{ }Any\DUrole{p}{{]}}\DUrole{w}{ }\DUrole{p}{|}\DUrole{w}{ }None}\DUrole{w}{ }\DUrole{o}{=}\DUrole{w}{ }\DUrole{default_value}{None}}\sphinxparamcomma \sphinxparam{\DUrole{n}{camera\_config}\DUrole{p}{:}\DUrole{w}{ }\DUrole{n}{Dict\DUrole{p}{{[}}str\DUrole{p}{,}\DUrole{w}{ }Any\DUrole{p}{{]}}\DUrole{w}{ }\DUrole{p}{|}\DUrole{w}{ }None}\DUrole{w}{ }\DUrole{o}{=}\DUrole{w}{ }\DUrole{default_value}{None}}}
{}
\pysigstopsignatures
\sphinxAtStartPar
Bases : {\hyperref[\detokenize{api/core:core.mode_base.Mode}]{\sphinxcrossref{\sphinxcode{\sphinxupquote{Mode}}}}}

\sphinxAtStartPar
Mode « Machine à lire » : capture caméra \sphinxhyphen{}\textgreater{} OCR \sphinxhyphen{}\textgreater{} synthèse \sphinxhyphen{}\textgreater{} lecture.

\sphinxAtStartPar
Ce mode ne pilote pas directement le keypad matériel : il reçoit les
touches non\sphinxhyphen{}globales via ModeManager.on\_key\_pressed().


\subsubsection{Commandes par défaut}
\label{\detokenize{api/modes:commandes-par-defaut}}\begin{itemize}
\item {} 
\sphinxAtStartPar
\sphinxtitleref{1} : capturer et lancer le pipeline complet.

\item {} 
\sphinxAtStartPar
\sphinxtitleref{3} : annuler la lecture ou le pipeline.

\item {} 
\sphinxAtStartPar
\sphinxtitleref{4} : reculer dans l’audio.

\item {} 
\sphinxAtStartPar
\sphinxtitleref{5} : pause / reprise.

\item {} 
\sphinxAtStartPar
\sphinxtitleref{6} : avancer dans l’audio.

\item {} 
\sphinxAtStartPar
\sphinxtitleref{*} : relire le dernier WAV généré.

\end{itemize}


\subsubsection{Contraintes importantes}
\label{\detokenize{api/modes:contraintes-importantes}}\begin{itemize}
\item {} 
\sphinxAtStartPar
Le pipeline est lancé dans un thread dédié pour ne pas bloquer
l’interface globale.

\item {} 
\sphinxAtStartPar
Les messages d’erreur et de statut sont joués sous forme de sons
courts afin d’éviter une concurrence trop forte avec la lecture
longue du document.

\end{itemize}
\index{KEY\_CAPTURE (attribut modes.mode\_reading\_machine.ModeReadingMachine)@\spxentry{KEY\_CAPTURE}\spxextra{attribut modes.mode\_reading\_machine.ModeReadingMachine}}

\begin{fulllineitems}
\phantomsection\label{\detokenize{api/modes:modes.mode_reading_machine.ModeReadingMachine.KEY_CAPTURE}}
\pysigstartsignatures
\pysigline
{\sphinxbfcode{\sphinxupquote{KEY\_CAPTURE}}\sphinxbfcode{\sphinxupquote{\DUrole{w}{ }\DUrole{p}{=}\DUrole{w}{ }\textquotesingle{}1\textquotesingle{}}}}
\pysigstopsignatures
\end{fulllineitems}

\index{KEY\_CANCEL (attribut modes.mode\_reading\_machine.ModeReadingMachine)@\spxentry{KEY\_CANCEL}\spxextra{attribut modes.mode\_reading\_machine.ModeReadingMachine}}

\begin{fulllineitems}
\phantomsection\label{\detokenize{api/modes:modes.mode_reading_machine.ModeReadingMachine.KEY_CANCEL}}
\pysigstartsignatures
\pysigline
{\sphinxbfcode{\sphinxupquote{KEY\_CANCEL}}\sphinxbfcode{\sphinxupquote{\DUrole{w}{ }\DUrole{p}{=}\DUrole{w}{ }\textquotesingle{}3\textquotesingle{}}}}
\pysigstopsignatures
\end{fulllineitems}

\index{KEY\_BACKWARD (attribut modes.mode\_reading\_machine.ModeReadingMachine)@\spxentry{KEY\_BACKWARD}\spxextra{attribut modes.mode\_reading\_machine.ModeReadingMachine}}

\begin{fulllineitems}
\phantomsection\label{\detokenize{api/modes:modes.mode_reading_machine.ModeReadingMachine.KEY_BACKWARD}}
\pysigstartsignatures
\pysigline
{\sphinxbfcode{\sphinxupquote{KEY\_BACKWARD}}\sphinxbfcode{\sphinxupquote{\DUrole{w}{ }\DUrole{p}{=}\DUrole{w}{ }\textquotesingle{}4\textquotesingle{}}}}
\pysigstopsignatures
\end{fulllineitems}

\index{KEY\_PLAY\_PAUSE (attribut modes.mode\_reading\_machine.ModeReadingMachine)@\spxentry{KEY\_PLAY\_PAUSE}\spxextra{attribut modes.mode\_reading\_machine.ModeReadingMachine}}

\begin{fulllineitems}
\phantomsection\label{\detokenize{api/modes:modes.mode_reading_machine.ModeReadingMachine.KEY_PLAY_PAUSE}}
\pysigstartsignatures
\pysigline
{\sphinxbfcode{\sphinxupquote{KEY\_PLAY\_PAUSE}}\sphinxbfcode{\sphinxupquote{\DUrole{w}{ }\DUrole{p}{=}\DUrole{w}{ }\textquotesingle{}5\textquotesingle{}}}}
\pysigstopsignatures
\end{fulllineitems}

\index{KEY\_FORWARD (attribut modes.mode\_reading\_machine.ModeReadingMachine)@\spxentry{KEY\_FORWARD}\spxextra{attribut modes.mode\_reading\_machine.ModeReadingMachine}}

\begin{fulllineitems}
\phantomsection\label{\detokenize{api/modes:modes.mode_reading_machine.ModeReadingMachine.KEY_FORWARD}}
\pysigstartsignatures
\pysigline
{\sphinxbfcode{\sphinxupquote{KEY\_FORWARD}}\sphinxbfcode{\sphinxupquote{\DUrole{w}{ }\DUrole{p}{=}\DUrole{w}{ }\textquotesingle{}6\textquotesingle{}}}}
\pysigstopsignatures
\end{fulllineitems}

\index{KEY\_REPLAY (attribut modes.mode\_reading\_machine.ModeReadingMachine)@\spxentry{KEY\_REPLAY}\spxextra{attribut modes.mode\_reading\_machine.ModeReadingMachine}}

\begin{fulllineitems}
\phantomsection\label{\detokenize{api/modes:modes.mode_reading_machine.ModeReadingMachine.KEY_REPLAY}}
\pysigstartsignatures
\pysigline
{\sphinxbfcode{\sphinxupquote{KEY\_REPLAY}}\sphinxbfcode{\sphinxupquote{\DUrole{w}{ }\DUrole{p}{=}\DUrole{w}{ }\textquotesingle{}*\textquotesingle{}}}}
\pysigstopsignatures
\end{fulllineitems}

\index{on\_enter() (méthode modes.mode\_reading\_machine.ModeReadingMachine)@\spxentry{on\_enter()}\spxextra{méthode modes.mode\_reading\_machine.ModeReadingMachine}}

\begin{fulllineitems}
\phantomsection\label{\detokenize{api/modes:modes.mode_reading_machine.ModeReadingMachine.on_enter}}
\pysigstartsignatures
\pysiglinewithargsret
{\sphinxbfcode{\sphinxupquote{on\_enter}}}
{}
{{ $\rightarrow$ None}}
\pysigstopsignatures
\sphinxAtStartPar
Prépare le mode à l’utilisation.

\sphinxAtStartPar
Actions réalisées :
\sphinxhyphen{} réinitialise l’état d’annulation ;
\sphinxhyphen{} sonde la disponibilité de la caméra ;
\sphinxhyphen{} lance le thread de surveillance si la caméra est absente.

\end{fulllineitems}

\index{on\_exit() (méthode modes.mode\_reading\_machine.ModeReadingMachine)@\spxentry{on\_exit()}\spxextra{méthode modes.mode\_reading\_machine.ModeReadingMachine}}

\begin{fulllineitems}
\phantomsection\label{\detokenize{api/modes:modes.mode_reading_machine.ModeReadingMachine.on_exit}}
\pysigstartsignatures
\pysiglinewithargsret
{\sphinxbfcode{\sphinxupquote{on\_exit}}}
{}
{{ $\rightarrow$ None}}
\pysigstopsignatures
\sphinxAtStartPar
Nettoie tout ce qui a été démarré par le mode.

\sphinxAtStartPar
Cette méthode doit laisser le mode dans un état neutre pour un
éventuel retour ultérieur, sans fuite de lecture ou de thread.

\end{fulllineitems}

\index{on\_short\_press() (méthode modes.mode\_reading\_machine.ModeReadingMachine)@\spxentry{on\_short\_press()}\spxextra{méthode modes.mode\_reading\_machine.ModeReadingMachine}}

\begin{fulllineitems}
\phantomsection\label{\detokenize{api/modes:modes.mode_reading_machine.ModeReadingMachine.on_short_press}}
\pysigstartsignatures
\pysiglinewithargsret
{\sphinxbfcode{\sphinxupquote{on\_short\_press}}}
{}
{{ $\rightarrow$ None}}
\pysigstopsignatures
\sphinxAtStartPar
Aucun comportement sur le bouton rotatif dans ce mode.

\sphinxAtStartPar
La machine à lire utilise le keypad pour éviter de surcharger le
bouton unique du sélecteur avec trop de commandes.

\end{fulllineitems}

\index{on\_long\_press() (méthode modes.mode\_reading\_machine.ModeReadingMachine)@\spxentry{on\_long\_press()}\spxextra{méthode modes.mode\_reading\_machine.ModeReadingMachine}}

\begin{fulllineitems}
\phantomsection\label{\detokenize{api/modes:modes.mode_reading_machine.ModeReadingMachine.on_long_press}}
\pysigstartsignatures
\pysiglinewithargsret
{\sphinxbfcode{\sphinxupquote{on\_long\_press}}}
{}
{{ $\rightarrow$ None}}
\pysigstopsignatures
\sphinxAtStartPar
Aucun comportement sur l’appui long du bouton rotatif.

\sphinxAtStartPar
Les commandes sont volontairement regroupées sur le clavier 4x4.

\end{fulllineitems}

\index{on\_keypad\_key() (méthode modes.mode\_reading\_machine.ModeReadingMachine)@\spxentry{on\_keypad\_key()}\spxextra{méthode modes.mode\_reading\_machine.ModeReadingMachine}}

\begin{fulllineitems}
\phantomsection\label{\detokenize{api/modes:modes.mode_reading_machine.ModeReadingMachine.on_keypad_key}}
\pysigstartsignatures
\pysiglinewithargsret
{\sphinxbfcode{\sphinxupquote{on\_keypad\_key}}}
{\sphinxparam{\DUrole{n}{key}\DUrole{p}{:}\DUrole{w}{ }\DUrole{n}{str}}}
{{ $\rightarrow$ None}}
\pysigstopsignatures
\sphinxAtStartPar
Route une touche du keypad vers la commande métier correspondante.

\sphinxAtStartPar
Seules les touches non\sphinxhyphen{}globales arrivent ici : volume et vitesse
sont interceptés plus haut par \sphinxtitleref{ModeManager}.

\end{fulllineitems}


\end{fulllineitems}



\subsection{modes.mode\_thermometer module}
\label{\detokenize{api/modes:module-modes.mode_thermometer}}\label{\detokenize{api/modes:modes-mode-thermometer-module}}\index{module@\spxentry{module}!modes.mode\_thermometer@\spxentry{modes.mode\_thermometer}}\index{modes.mode\_thermometer@\spxentry{modes.mode\_thermometer}!module@\spxentry{module}}

\subsubsection{modes/mode\_thermometer.py}
\label{\detokenize{api/modes:modes-mode-thermometer-py}}
\sphinxAtStartPar
Mode « Thermomètre WiFi » pour l’assistant.


\paragraph{Rôle :}
\label{\detokenize{api/modes:id9}}\begin{itemize}
\item {} 
\sphinxAtStartPar
Interface entre l’utilisateur (boutons, sélecteur) et le Driver Thermomètre.

\item {} 
\sphinxAtStartPar
Ce mode suppose que la Pi et l’ESP32 sont sur le même réseau WiFi existant.

\end{itemize}


\paragraph{Comportement UX :}
\label{\detokenize{api/modes:id10}}\begin{itemize}
\item {} 
\sphinxAtStartPar
on\_enter : Démarrage du driver (recherche du capteur sur le réseau).

\item {} 
\sphinxAtStartPar
on\_exit : Arrêt du driver.

\item {} 
\sphinxAtStartPar
Appui court : Lit la température (« Il fait 22 virgule 5 degrés »).

\item {} 
\sphinxAtStartPar
Appui long : Annonce l’état (« Le thermomètre est connecté au réseau local »).

\end{itemize}
\index{ModeThermometre (classe dans modes.mode\_thermometer)@\spxentry{ModeThermometre}\spxextra{classe dans modes.mode\_thermometer}}

\begin{fulllineitems}
\phantomsection\label{\detokenize{api/modes:modes.mode_thermometer.ModeThermometre}}
\pysigstartsignatures
\pysiglinewithargsret
{\sphinxbfcode{\sphinxupquote{\DUrole{k}{class}\DUrole{w}{ }}}\sphinxcode{\sphinxupquote{modes.mode\_thermometer.}}\sphinxbfcode{\sphinxupquote{ModeThermometre}}}
{\sphinxparam{\DUrole{n}{speaker}\DUrole{p}{:}\DUrole{w}{ }\DUrole{n}{{\hyperref[\detokenize{api/core:core.speaker.Speaker}]{\sphinxcrossref{Speaker}}}}}}
{}
\pysigstopsignatures
\sphinxAtStartPar
Bases : {\hyperref[\detokenize{api/core:core.mode_base.Mode}]{\sphinxcrossref{\sphinxcode{\sphinxupquote{Mode}}}}}

\sphinxAtStartPar
Mode Thermomètre connecté via le réseau local.

\sphinxAtStartPar
Le mode s’appuie sur \sphinxtitleref{ThermometerDriver} pour toute la partie
communication. Son rôle est volontairement limité à la logique UX :
formater la valeur pour la voix et répondre aux appuis utilisateur.
\index{\_\_init\_\_() (méthode modes.mode\_thermometer.ModeThermometre)@\spxentry{\_\_init\_\_()}\spxextra{méthode modes.mode\_thermometer.ModeThermometre}}

\begin{fulllineitems}
\phantomsection\label{\detokenize{api/modes:modes.mode_thermometer.ModeThermometre.__init__}}
\pysigstartsignatures
\pysiglinewithargsret
{\sphinxbfcode{\sphinxupquote{\_\_init\_\_}}}
{\sphinxparam{\DUrole{n}{speaker}\DUrole{p}{:}\DUrole{w}{ }\DUrole{n}{{\hyperref[\detokenize{api/core:core.speaker.Speaker}]{\sphinxcrossref{Speaker}}}}}}
{}
\pysigstopsignatures
\sphinxAtStartPar
Initialise le mode avec sa synthèse vocale et son driver.

\end{fulllineitems}

\index{on\_enter() (méthode modes.mode\_thermometer.ModeThermometre)@\spxentry{on\_enter()}\spxextra{méthode modes.mode\_thermometer.ModeThermometre}}

\begin{fulllineitems}
\phantomsection\label{\detokenize{api/modes:modes.mode_thermometer.ModeThermometre.on_enter}}
\pysigstartsignatures
\pysiglinewithargsret
{\sphinxbfcode{\sphinxupquote{on\_enter}}}
{}
{{ $\rightarrow$ None}}
\pysigstopsignatures
\sphinxAtStartPar
Arrivée sur le mode Thermomètre :
On lance le driver qui va chercher l’ESP32 sur le réseau local.

\end{fulllineitems}

\index{on\_exit() (méthode modes.mode\_thermometer.ModeThermometre)@\spxentry{on\_exit()}\spxextra{méthode modes.mode\_thermometer.ModeThermometre}}

\begin{fulllineitems}
\phantomsection\label{\detokenize{api/modes:modes.mode_thermometer.ModeThermometre.on_exit}}
\pysigstartsignatures
\pysiglinewithargsret
{\sphinxbfcode{\sphinxupquote{on\_exit}}}
{}
{{ $\rightarrow$ None}}
\pysigstopsignatures
\sphinxAtStartPar
Sortie du mode :
On arrête le driver pour économiser des ressources réseau.

\end{fulllineitems}

\index{on\_short\_press() (méthode modes.mode\_thermometer.ModeThermometre)@\spxentry{on\_short\_press()}\spxextra{méthode modes.mode\_thermometer.ModeThermometre}}

\begin{fulllineitems}
\phantomsection\label{\detokenize{api/modes:modes.mode_thermometer.ModeThermometre.on_short_press}}
\pysigstartsignatures
\pysiglinewithargsret
{\sphinxbfcode{\sphinxupquote{on\_short\_press}}}
{}
{{ $\rightarrow$ None}}
\pysigstopsignatures
\sphinxAtStartPar
Appui court : Annonce la température.

\end{fulllineitems}

\index{on\_long\_press() (méthode modes.mode\_thermometer.ModeThermometre)@\spxentry{on\_long\_press()}\spxextra{méthode modes.mode\_thermometer.ModeThermometre}}

\begin{fulllineitems}
\phantomsection\label{\detokenize{api/modes:modes.mode_thermometer.ModeThermometre.on_long_press}}
\pysigstartsignatures
\pysiglinewithargsret
{\sphinxbfcode{\sphinxupquote{on\_long\_press}}}
{}
{{ $\rightarrow$ None}}
\pysigstopsignatures
\sphinxAtStartPar
Appui long : Annonce l’état de connexion.

\end{fulllineitems}


\end{fulllineitems}



\subsection{Module contents}
\label{\detokenize{api/modes:module-modes}}\label{\detokenize{api/modes:module-contents}}\index{module@\spxentry{module}!modes@\spxentry{modes}}\index{modes@\spxentry{modes}!module@\spxentry{module}}

\renewcommand{\indexname}{Index des modules Python}
\begin{sphinxtheindex}
\let\bigletter\sphinxstyleindexlettergroup
\bigletter{c}
\item\relax\sphinxstyleindexentry{config}\sphinxstyleindexpageref{api/config:\detokenize{module-config}}
\item\relax\sphinxstyleindexentry{config.config\_loader}\sphinxstyleindexpageref{api/config:\detokenize{module-config.config_loader}}
\item\relax\sphinxstyleindexentry{core}\sphinxstyleindexpageref{api/core:\detokenize{module-core}}
\item\relax\sphinxstyleindexentry{core.config\_loader}\sphinxstyleindexpageref{api/core:\detokenize{module-core.config_loader}}
\item\relax\sphinxstyleindexentry{core.mode\_base}\sphinxstyleindexpageref{api/core:\detokenize{module-core.mode_base}}
\item\relax\sphinxstyleindexentry{core.mode\_manager}\sphinxstyleindexpageref{api/core:\detokenize{module-core.mode_manager}}
\item\relax\sphinxstyleindexentry{core.mode\_registry}\sphinxstyleindexpageref{api/core:\detokenize{module-core.mode_registry}}
\item\relax\sphinxstyleindexentry{core.speaker}\sphinxstyleindexpageref{api/core:\detokenize{module-core.speaker}}
\indexspace
\bigletter{h}
\item\relax\sphinxstyleindexentry{hardware}\sphinxstyleindexpageref{api/hardware:\detokenize{module-hardware}}
\item\relax\sphinxstyleindexentry{hardware.devices}\sphinxstyleindexpageref{api/hardware.devices:\detokenize{module-hardware.devices}}
\item\relax\sphinxstyleindexentry{hardware.devices.caliper}\sphinxstyleindexpageref{api/hardware.devices:\detokenize{module-hardware.devices.caliper}}
\item\relax\sphinxstyleindexentry{hardware.devices.esp32\_thermometer}\sphinxstyleindexpageref{api/hardware.devices:\detokenize{module-hardware.devices.esp32_thermometer}}
\item\relax\sphinxstyleindexentry{hardware.devices.owon\_multimeter}\sphinxstyleindexpageref{api/hardware.devices:\detokenize{module-hardware.devices.owon_multimeter}}
\item\relax\sphinxstyleindexentry{hardware.devices.text\_reader\_device}\sphinxstyleindexpageref{api/hardware.devices:\detokenize{module-hardware.devices.text_reader_device}}
\item\relax\sphinxstyleindexentry{hardware.platform}\sphinxstyleindexpageref{api/hardware.platform:\detokenize{module-hardware.platform}}
\item\relax\sphinxstyleindexentry{hardware.platform.keypad\_4x4}\sphinxstyleindexpageref{api/hardware.platform:\detokenize{module-hardware.platform.keypad_4x4}}
\item\relax\sphinxstyleindexentry{hardware.platform.pi\_camera}\sphinxstyleindexpageref{api/hardware.platform:\detokenize{module-hardware.platform.pi_camera}}
\item\relax\sphinxstyleindexentry{hardware.platform.rotary\_selector}\sphinxstyleindexpageref{api/hardware.platform:\detokenize{module-hardware.platform.rotary_selector}}
\item\relax\sphinxstyleindexentry{hardware.transports}\sphinxstyleindexpageref{api/hardware.transports:\detokenize{module-hardware.transports}}
\item\relax\sphinxstyleindexentry{hardware.transports.ble\_client}\sphinxstyleindexpageref{api/hardware.transports:\detokenize{module-hardware.transports.ble_client}}
\item\relax\sphinxstyleindexentry{hardware.transports.nrf24\_transport}\sphinxstyleindexpageref{api/hardware.transports:\detokenize{module-hardware.transports.nrf24_transport}}
\item\relax\sphinxstyleindexentry{hardware.transports.wifi\_client}\sphinxstyleindexpageref{api/hardware.transports:\detokenize{module-hardware.transports.wifi_client}}
\indexspace
\bigletter{m}
\item\relax\sphinxstyleindexentry{modes}\sphinxstyleindexpageref{api/modes:\detokenize{module-modes}}
\item\relax\sphinxstyleindexentry{modes.mode\_caliper}\sphinxstyleindexpageref{api/modes:\detokenize{module-modes.mode_caliper}}
\item\relax\sphinxstyleindexentry{modes.mode\_datetime}\sphinxstyleindexpageref{api/modes:\detokenize{module-modes.mode_datetime}}
\item\relax\sphinxstyleindexentry{modes.mode\_dummy}\sphinxstyleindexpageref{api/modes:\detokenize{module-modes.mode_dummy}}
\item\relax\sphinxstyleindexentry{modes.mode\_multimeter}\sphinxstyleindexpageref{api/modes:\detokenize{module-modes.mode_multimeter}}
\item\relax\sphinxstyleindexentry{modes.mode\_reading\_machine}\sphinxstyleindexpageref{api/modes:\detokenize{module-modes.mode_reading_machine}}
\item\relax\sphinxstyleindexentry{modes.mode\_thermometer}\sphinxstyleindexpageref{api/modes:\detokenize{module-modes.mode_thermometer}}
\end{sphinxtheindex}

\renewcommand{\indexname}{Index}
\printindex
\end{document}